\chapter{人体化身与动画}

\paragraph*{本章阅读锚点}
人体化身里的 3DGS 与普通静态重建不同：它不是只要求某一时刻渲染得像，而是要求身份、表情、身体姿态、衣物褶皱、手部交互和语音驱动在新姿态下仍能保持一致。原始 3DGS 的锚点是无拓扑的显式高斯；人体方法通常会把它接到 FLAME、SMPL、SMPL-X、MANO、mesh、tetrahedron 或 blendshape 这些可控先验上。读本章时可以抓住三个问题：第一，高斯放在哪里，是自由点、模板表面、UV 图、网格三角面、四面体还是神经点云；第二，高斯怎样动，是 LBS、blendshape、MLP deformation、显式轨迹、扩散生成还是前馈 Transformer；第三，哪些部分需要单独处理，头发、口腔、衣物、手、遮挡区域和动态物体往往不能用同一套刚性人体先验解释。

\section{人脸重建与动画}

人脸头像的基础锚点是 FLAME 与 blendshape：它们提供可控表情和头部姿态，但缺少头发、口腔、毛孔、衣领和复杂光照。3DGS 进入这一方向后，主要变化不是“把脸换成高斯”，而是把高斯绑定到可动画模型上，再用混合表示、UV 空间生成、先验学习或扩散蒸馏补足模板无法表达的部分。



% METHOD: 2403.18784
\paragraph{SplatFace（高斯与可优化 3DMM 表面联合）（2024.03）}
SplatFace~\cite{2403_18784} 的锚点是 3DMM 表面与 Gaussian Splatting 联合优化。它不假设预定几何完全准确，而是让高斯和可变形表面同时精炼。核心改动是 splat-to-surface 距离度量，不只看高斯中心到表面的距离，还考虑协方差形状和方向，使高斯更贴合表面。它又加入 world-space densification，利用表面法线在需要高频皮肤细节的区域增加高斯，避免悬浮杂散点。相对 MixedGaussianAvatar，SplatFace 更强调最后能输出精确网格和可渲染高斯两种结果。
\textbf{实现细节：} SplatFace的核心是将可优化的3DMM表面与高斯溅射联合优化。方法包含：（1）联合优化框架：高斯溅射和3DMM表面使用相同的模板表面初始化。高斯和表面通过非刚性对齐过程紧密耦合，同时优化。（2）Splat-to-Surface距离度量：不同于仅考虑高斯中心到表面距离的简单方法，该度量同时考虑高斯位置和协方差，即考虑高斯的形状和方向，提供更准确的对齐。（3）世界空间密集化：在3DGS原有的视角空间自适应密集化基础上，增加了基于表面信息的世界空间密集化。利用表面的法线信息，在需要更多高频细节的区域（如皮肤纹理）生成额外高斯，同时避免生成悬浮的杂散高斯。（4）最终输出包括用于新视角合成的高斯溅射模型和精确的3D网格表示。 提出SplatFace，将通用3DMM（3D Morphable Model）融入Gaussian Splatting框架，通过联合优化策略同时精炼高斯和可变形表面。提出splat-to-surface距离度量，考虑高斯位置和协方差来改善对齐。引入世界空间密集化过程，利用表面信息增强高频皮肤细节捕获。 依赖3DMM模板的泛化能力，对极端表情和遮挡的建模受限；联合优化增加了训练复杂度；未考虑头发等非刚性组件的精细建模。
\textbf{实现补充：} 实现上还要核对输入、输出、参数、特征、初始化、训练损失、正则、渲染或推理接口，以及优化中的剪枝、分裂、采样、编码或解码步骤；这些抓手决定该方法到底改了 3DGS 管线的哪一环。
\textbf{方法信息：} ByteDance 参与了这篇论文。
\textbf{设计动机与边界：} 相对原始 3DGS 以独立高斯和图像重建误差为主的设定，SplatFace 把问题落在以下可实现的改动上：SplatFace的核心是将可优化的3DMM表面与高斯溅射联合优化。方法包含：（1）联合优化框架：高斯溅射和3DMM表面使用相同的模板表面初始化。高斯和表面通过非刚性对齐过程紧密耦合，同时优化。（2）Splat-to-Surface距离度量：不同于仅考虑高斯中心到表面距离的简单方法，该度量同时考虑高斯位置和协方差，即考虑高斯的形状和方向，提供更准确的对齐。（3）世界空间密集化：在3DGS原有的视角空间自适应密集化基础上，增加了基于表面信息的世界空间密集化。利用表面的法线信息，在需要更多高频细节的区域（如皮肤纹理）生成额外高斯，同时避免生成悬浮的杂散高斯。（4）最终输出包括用于新视角合成的高斯溅射模型和精确的3D网格表示。其设计目的不是单独增加一个模块，而是用上述表示、监督或渲染约束消除原流程在目标任务中的具体失配。提出SplatFace，将通用3DMM（3D Morphable Model）融入Gaussian Splatting框架，通过联合优化策略同时精炼高斯和可变形表面。提出splat-to-surface距离度量，考虑高斯位置和协方差来改善对齐。引入世界空间密集化过程，利用表面信息增强高频皮肤细节捕获。适用时还应保留论文给出的边界：依赖3DMM模板的泛化能力，对极端表情和遮挡的建模受限；联合优化增加了训练复杂度；未考虑头发等非刚性组件的精细建模。
% METHOD: 2404.19026
\paragraph{MeGA（面部网格与头发高斯的混合头像）（2024.04）}
MeGA~\cite{2404_19026} 的关键判断是：面部、头发和编辑需求不适合用单一表示解决。它用增强 FLAME 网格和 delayed neural rendering 表示面部，通过 UV 位移图提供个性化几何细节；头发则用 3DGS 表示静态规范模型，再通过刚性变换和 MLP deformation 处理动态表情影响。遮挡感知混合负责面部网格和头发高斯的边界。相对纯 FLAME 高斯绑定，MeGA 的增量是组件化表示选择：脸用可编辑拓扑，头发用无拓扑高斯。
\textbf{实现细节：} MeGA（2024.04） （Hybrid Mesh-Gaussian Head Avatar）则从根本上质疑了上述所有方法共享的一个假设：能否用同一种表示同时高质量地建模面部的皮肤和头部的头发？锚点回顾：Gaussian Blendshapes 及其后续工作（NPGA、GraphAvatar 等）均用统一的 3DGS 表示建模整个头部，但人头的不同组件具有截然不同的特性：面部皮肤近似光滑曲面且可用低维参数空间驱动，而头发则是包含大量细丝状结构的体积型对象。改了什么：MeGA 放弃了"单一表示"策略，转而采用混合网格-高斯解耦架构：用增强的 FLAME 神经网格建模面部（皮肤、眼睛、嘴内），用静态规范空间中的 3DGS 独立建模头发，两者通过遮挡感知混合渲染（occlusion-aware blending）合成最终图像。设计动机：这一解耦基于一个核心观察：纯 3DGS 方法（如 GaussianAvatars）虽然在头发等体积结构上表现出色，但在面部区域存在抗锯齿缺陷和穿透伪影（如张大嘴时高斯交错），而基于网格的延迟神经渲染恰好擅长产生高质量的面部纹理细节。与其用一种表示的弱点去妥协另一种表示的优势，不如让每种表示各司其职。
\textbf{实现补充：} 实现上还要核对输入、输出、参数、特征、初始化、训练损失、正则、渲染或推理接口，以及优化中的剪枝、分裂、采样、编码或解码步骤；这些抓手决定该方法到底改了 3DGS 管线的哪一环。
\textbf{方法信息：} Tencent 参与了这篇论文。
\textbf{设计动机与边界：} 相对原始 3DGS 以独立高斯和图像重建误差为主的设定，MeGA 把问题落在以下可实现的改动上：MeGA的核心思想是为不同头部组件使用最适合的表示。面部部分：采用增强FLAME网格，学习UV位移图提供逐顶点偏移以获得个性化几何细节；通过延迟神经渲染获得照片级真实感，将神经纹理解耦为三个有意义的部件（如漫反射色、镜面反射、细节法线）。头发部分：使用3DGS构建静态规范头发模型，通过刚性变换和MLP变形场处理动态表情。关键设计是遮挡感知混合，解决面部网格和头发高斯之间的遮挡问题。训练在NeRSemble数据集上进行。其设计目的不是单独增加一个模块，而是用上述表示、监督或渲染约束消除原流程在目标任务中的具体失配。提出混合表示策略，用增强FLAME网格+延迟神经渲染建模面部，3DGS建模头发，解决单一表示无法同时高质量渲染所有头部组件的问题；设计UV位移图提供个性化几何细节；将神经纹理解耦为三个有意义的部件；引入遮挡感知混合实现面部和头发的无缝融合。适用时还应保留论文给出的边界：两个表示的融合需要精细的遮挡处理；训练需要多视角视频数据；面部和头发的解耦依赖分割质量。
% METHOD: 2404.19398
\paragraph{3D Gaussian Blendshapes（高斯混合形状锚点）（2024.04）}
3D Gaussian Blendshapes~\cite{2404_19398} 是头部高斯动画的一条基础路线。它把经典 blendshape 从 mesh 顶点扩展到 3DGS 属性：中性模型和每个表情 blendshape 都由高斯的位置、不透明度、旋转、尺度和 SH 系数组成，任意表情通过线性混合属性差值得到。FLAME 姿态和关节参数再通过 LBS 变换高斯。方法还单独定义口腔内部高斯组，只受下颌控制，以处理牙齿和口腔。相对后续学习式表情网络，它的好处是可解释、可直接由 FLAME 系数驱动。
\textbf{实现细节：} Gaussian Blendshapes（2024.04） 是该方向的先驱工作，提出了基于混合形状（blendshape）的 3D 高斯表示。Baseline 与问题：传统参数化人脸模型（如 FLAME）虽然能通过线性混合形状控制表情，但其基于网格的表示在渲染速度和细节保真度之间存在根本矛盾：网格分辨率限制了高频表面细节的表达，而 NeRF 虽然能表达细节却无法实时渲染。改了什么：Gaussian Blendshapes 将 FLAME 的混合形状机制从网格域迁移到 3D 高斯域，用一组 3D 高斯基元代替网格顶点来表示头部几何和外观。实现原理：具体而言，该方法定义一个中性基础模型 $B_0$ 和 $K$ 个表情混合形状 $\Delta B_1, \Delta B_2, \ldots, \Delta B_K$，全部表示为 3D 高斯集合。每个高斯具有位置 $x$、不透明度 $o$、旋转 $q$、缩放 $s$ 和球谐系数（Spherical Harmonics, SH，用于表示视角依赖颜色的球面基函数，下文简称 SH 系数）$SH$ 五种属性。混合形状定义为各高斯属性相对于中性模型的偏移量 $\Delta B_k = B_k - B_0$。任意表情的化身通过线性混合生成：$\hat{B} = B_0 + \sum_{k=1}^{K} \alpha_k \Delta B_k$，其中 $\{\alpha_k\}$ 为 PCA 表情系数。该模型采用 FLAME 的 50 个基表情，同时利用 FLAME 的关节和姿态参数 $\theta$ 通过线性混合蒙皮（LBS）变换高斯位置：$B^{*} = \mathrm{LBS}(\hat{B}, \theta)$。此外，该方法为嘴部内部定义独立的固定高斯 $B_m$，仅随下颌关节运动而不受表情影响，解决了 FLAME 网格无法覆盖嘴内区域的问题。设计动机：将混合形状应用于高斯域而非网格域，其核心优势在于：高斯基元的半透明核函数天然适合建模头发、透明边界等离散表面难以处理的区域；同时基于 tile 的可微分光栅化器实现了实时渲染。
\textbf{方法信息：} 引用很多，Semantic Scholar 约 82 次。
\textbf{设计动机与边界：} 相对原始 3DGS 以独立高斯和图像重建误差为主的设定，3D Gaussian Blendshapes for Head Avatar Animation 把问题落在以下可实现的改动上：表示由中性基础模型B\_0和K个表情blendshape B\_k组成，全部用3D高斯表示（位置、不透明度、旋转、缩放、SH系数）。每个高斯blendshape与FLAME的mesh blendshape语义一致。任意表情通过线性混合实现：B\_exp = B\_0 + sum(alpha\_k * Delta\_B\_k)，其中Delta\_B\_k = B\_k - B\_0为属性差值。FLAME的关节和姿态参数通过线性混合蒙皮（LBS）变换高斯。定义独立的口腔内部高斯组B\_m，仅受下颌关节控制不受表情影响。初始化通过泊松盘采样在FLAME网格表面分布点，然后联合优化所有高斯属性。其设计目的不是单独增加一个模块，而是用上述表示、监督或渲染约束消除原流程在目标任务中的具体失配。提出3D高斯混合变形表示，将经典blendshape模型扩展到3DGS；通过线性混合中性模型和表情blendshape的高斯属性实现任意表情生成；定义单独的口腔内部高斯组处理牙齿和口腔；支持FLAME表情系数直接驱动，实现实时高保真动画。适用时还应保留论文给出的边界：依赖FLAME跟踪质量；blendshape数量影响表达能力；头发建模受FLAME网格限制。
% METHOD: 2405.12069
\paragraph{Gaussian Head and Shoulders（锚点高斯引导神经纹理）（2024.05）}
Gaussian Head and Shoulders~\cite{2405_12069} 指出一个常被忽略的问题：每个高斯只有有限视角辐射，难以表示衣物这类空间变化纹理。方法因此把头部和上半身分开：头部用 FLAME 驱动的高斯和神经 LBS，身体衣物用神经纹理表示粗颜色与姿态相关细节。稀疏锚点高斯约束图像平面到纹理空间的 neural warping，无需精确 UV 也能把衣物图案对齐。相对纯高斯上半身，它的改变是承认衣物纹理更适合纹理域，而高斯更适合引导几何和渲染。
\textbf{实现细节：} 方法针对3DGS表示穿衣上半身的根本限制——每个高斯只能表示单一视角的颜色，无法捕捉衣物的空间变化纹理——提出分层解决方案：(1) FLAME驱动的头部高斯：使用标准3D高斯通过神经LBS变形跟随FLAME 3DMM，学习个性化表情、姿态混合形状和LBS权重；(2) 身体部分的神经纹理：设计可学习的神经纹理，包含粗粒度基础颜色和姿态相关的细粒度颜色分量，能表示衣物的复杂空间变化图案；(3) 锚点高斯与扭曲场：优化一组稀疏的锚点高斯分布在上半身表面，这些锚点约束神经扭曲场的映射关系——从图像平面坐标映射到纹理空间，无需精确几何或UV映射即可正确渲染身体纹理；(4) 加速推理：提出消除MLP查询的加速推理方法，预计算扭曲映射实现约130FPS的稳定渲染速度。方法支持自重演和交叉重演任务。 发现并解决了3DGS在复杂空间变化纹理建模上的根本局限（每个高斯只能有单一方向辐射），提出锚点高斯引导的神经纹理扭曲方案，实现穿衣上半身高保真重建和130FPS加速推理 头部和身体使用不同的表示方法增加了系统复杂度；对大幅运动或极端姿态的泛化能力有限；神经纹理分辨率受限于GPU内存；当前仅支持上半身，未扩展到全身
\textbf{实现补充：} 实现上还要核对输入、输出、参数、特征、初始化、训练损失、正则、渲染或推理接口，以及优化中的剪枝、分裂、采样、编码或解码步骤；这些抓手决定该方法到底改了 3DGS 管线的哪一环。
\textbf{方法信息：} Google / DeepMind 参与了这篇论文。



% METHOD: 2405.19331
\paragraph{NPGA（NPHM 运动先验蒸馏到高斯）（2024.05）}
NPGA~\cite{2405_19331} 把 Neural Parametric Head Model 的丰富表情空间和 3DGS 的实时渲染结合起来。每个高斯除了标准属性，还带静态 latent feature，用于描述动态行为。动力学模块分为 coarse prior 和 detail network：前者借助 NPHM 运动先验预测中心偏移，后者预测高斯所有属性的细粒度动态，例如皱纹和血色变化。由于 NPHM 的 backward deformation 不适合 rasterization，NPGA 将其蒸馏为 forward deformation。相对 blendshape 线性混合，NPGA 更适合表达非线性细节。
\textbf{实现细节：} NPGA（2024.05） （Neural Parametric Gaussian Avatars）在架构上提出了根本性改进。Baseline 与问题：Gaussian Blendshapes 使用基于 PCA（Principal Component Analysis，主成分分析）的线性表情空间，其表达能力受限于基表情的数量和线性假设，难以捕捉个体特异性的细微表情变化（如特定皱纹模式、血流量变化引起的肤色变化）。改了什么：NPGA 将高斯基元的动态行为条件化于神经参数化头部模型（NPHM）的丰富表情空间，而非传统的基于网格的 3DMM；同时引入逐高斯潜在特征来增强化身的表达能力。实现原理：NPGA 由两个核心组件构成。(1) 规范高斯点云 $A_c$，除了标准的高斯属性（位置、缩放、旋转、不透明度、球谐系数）外，每个高斯还附加一个可学习的潜在特征 $f \in \mathbb{R}^{}8$。该特征与空间坐标无关，提供语义信息描述各基元的动态行为。(2) 动力学模块 $D$，分解为粗粒度先验网络 $D^p$ 和细节网络 $D^a$。
\textbf{方法信息：} 引用很多，Semantic Scholar 约 55 次。
\textbf{设计动机与边界：} 相对原始 3DGS 以独立高斯和图像重建误差为主的设定，NPGA 把问题落在以下可实现的改动上：NPGA的核心是将3DGS与神经参数化头部模型结合。（1）规范高斯表示：在标准3DGS基础上，为每个高斯添加静态潜在特征f（8维），提供描述动态行为的语义信息。类似位置编码，但不与空间坐标相关且具有无限空间分辨率。（2）动力学模块：由两个MLP组成。粗粒度先验网络Π\_coarse基于NPHM的运动先验预测高斯中心偏移；细节网络Π\_detail预测所有高斯属性（位置、旋转、缩放、颜色、不透明度）的偏移，捕捉NPHM无法表达的细粒度动态（如皱纹、血流变化）。（3）NPHM变形蒸馏：NPHM的后向变形场（从观测空间到规范空间）不兼容光栅化渲染。NPGA将其蒸馏为前向变形场（从规范空间到观测空间），通过优化使前向变形近似后向变形的逆。（4）正则化：对潜在特征和预测动态引入拉普拉斯正则化项，防止过拟合和过度动态表达。其设计目的不是单独增加一个模块，而是用上述表示、监督或渲染约束消除原流程在目标任务中的具体失配。提出NPGA，将3DGS的渲染效率与NPHM的丰富表情空间结合。每个高斯引入潜在特征增强动态表达能力。将NPHM的后向变形场蒸馏为前向变形兼容光栅化渲染。在NeRSemble数据集上自重演任务PSNR提升2.6dB，支持单目视频驱动的动画。适用时还应保留论文给出的边界：依赖多视角视频训练数据；NPHM的表情空间仍有局限；计算资源需求较高。
% METHOD: 2406.00440
\paragraph{Topo4D（固定拓扑高斯网格）（2024.06）}
Topo4D~\cite{2406_00440} 关注 4D 头部捕获和可用网格输出。原始 3DGS 渲染质量高，但高斯之间没有拓扑，难以生成规则布线的动态面部网格。Topo4D 的改变是 Gaussian Mesh：把每个高斯中心绑定到预定义拓扑网格顶点上，优化时高斯参数和顶点位置交替更新；几何阶段更新顶点和高斯，纹理阶段在 UV 空间学习高分辨率贴图，最后从学习到的高斯中提取规则布线的动态面部网格。相对 SplatFace，它更强调时序固定拓扑和毛孔级纹理，而不是少量输入下的表面优化。
\textbf{实现细节：} Topo4D（2024.06） 将高斯中心直接绑定到 FLAME 网格顶点而非三角形面片中心，利用预定义拓扑信息（8280 个顶点、16494 个面片）实现跨帧稳定的顶点级追踪。改了什么：不同于 MixedGaussianAvatar 的两阶段 2D+3D 混合策略，Topo4D 采用交替几何-纹理优化方案：在每帧先执行 1000 次几何优化迭代，再做 300 次纹理优化迭代。几何优化阶段引入物理损失和拓扑损失的双重约束。物理损失 $L_{phy} = \lambda_{rot} L_{rot} + \lambda_{iso} L_{iso}$（$\lambda_{rot}=\lambda_{iso}=20$）强制相邻帧间旋转一致性（$L_{rot}$，基于四元数的逐边旋转差异）和长程等距约束（$L_{iso}$，惩罚边长变化 $|d_{0,j} - d_{0,i}|_2 - |d_{t,j} - d_{t,i}|_2$）；拓扑损失 $L_{topo} = \lambda_{pos} L_{pos} + \lambda_{flat} L_{flat}$（$\lambda_{pos}=10^{-3}$, $\lambda_{flat}=2 \times 10^{-4}$）通过顶点-邻居平均位置偏差（$L_{pos}$）和相邻面片夹角余弦变化（$L_{flat}$）维持网格规则性和表面光滑度。
\textbf{方法信息：} Huawei 参与了这篇论文。
\textbf{设计动机与边界：} 相对原始 3DGS 以独立高斯和图像重建误差为主的设定，Topo4D 把问题落在以下可实现的改动上：Topo4D的核心是Gaussian Mesh表示，它将3D高斯的中心绑定到预定义拓扑的网格顶点上。每个高斯点包含可学习参数（位置、旋转、缩放、颜色、不透明度），但其中心受限于网格顶点位置。具体流程为：（1）基于第一帧的面部先验初始化拓扑集成的高斯集合；（2）对每个后续帧交替执行几何优化和纹理优化——几何优化通过梯度下降调整顶点位置和高斯参数，纹理优化利用高分辨率的UV空间学习8K纹理；（3）最后从学习到的高斯中提取具有规则布线的动态面部网格和高保真纹理。关键创新在于利用高斯的渲染能力同时约束其拓扑结构，解决了原始3DGS中高斯无拓扑关系导致无法提取规则网格的问题。其设计目的不是单独增加一个模块，而是用上述表示、监督或渲染约束消除原流程在目标任务中的具体失配。首次实现从校准的多视角时序图像直接优化密集对齐的4D头部和8K纹理贴图。通过将3D高斯与固定拓扑网格绑定，在保持渲染质量的同时解决了高斯缺乏拓扑关系的问题，生成具有规则布线排列的动态面部网格和毛孔级细节纹理。适用时还应保留论文给出的边界：依赖多视角校准的输入和面部先验，处理极端表情变化和头发等细碎结构的挑战仍然存在。
% METHOD: 2406.09377
\paragraph{GGHead（UV 空间 3D Gaussian GAN）（2024.06）}
GGHead~\cite{2406_09377} 是 3DGS 头部生成先验的重要路线。它用 StyleGAN2 类生成器在模板头部 UV 空间预测高斯属性，每个高斯有固定 UV 坐标，再双线性查询位置偏移、尺度、旋转、颜色和不透明度。这样把无结构高斯预测变成规则 2D 图像生成问题。渲染 UV 坐标的 total variation 损失鼓励相邻像素来自相邻高斯，提高几何稳定性。相对 EG23D，它更强调从大规模 2D 图像集合中学习 3D 一致头部生成先验。
\textbf{实现细节：} GGHead（2024.06） 在 3D GAN（Generative Adversarial Network，生成对抗网络）框架中采用 3DGS 表示。Baseline 与问题：3D GAN 训练的核心挑战在于如何用 2D CNN（Convolutional Neural Network，卷积神经网络）高效生成无结构的 3D 点云表示。3D 高斯作为点基表示本身是无序的，直接生成面临顺序模糊性和大量空白区域的困难。改了什么：GGHead 将 3D 高斯绑定到模板头部网格的 UV 空间中，利用 2D CNN 生成器在结构化的 UV 空间中预测高斯属性，巧妙地将 3D 生成问题转化为 2D 纹理生成问题。实现原理：使用 StyleGAN2 架构作为骨干网络 $B$，将潜在编码 $z \in \mathbb{R}^{512}$ 映射到多通道 UV 图 $M \in \mathbb{R}^{256 \times 256 \times 14}$，分别对应位置、缩放、旋转、颜色和不透明度五种高斯属性。对于位置属性，额外添加一个权重初始化为零的可学习 CNN 层 $Z$，确保训练初始时高斯精确落在模板网格上，这是稳定对抗训练的关键设计。3D 高斯通过在 UV 图上均匀采样获得——每个高斯有固定的 UV 坐标 $x_{uv} \in [0,1]^{}2$，通过双线性插值从属性图中查询其属性。
\textbf{方法信息：} 引用很多，Semantic Scholar 约 66 次。
\textbf{设计动机与边界：} 相对原始 3DGS 以独立高斯和图像重建误差为主的设定，GGHead 把问题落在以下可实现的改动上：GGHead采用3D GAN架构（生成器+可微分光栅化器+判别器），核心创新：(1) 模板UV空间生成：利用模板头部网格的UV布局规律性，通过StyleGAN2生成器在UV空间预测高斯属性（位置、尺度、旋转、颜色、透明度）。每个高斯有固定UV坐标，通过双线性插值查询属性。位置偏移初始化为零（高斯从模板表面开始），通过额外CNN层学习。(2) 总变差正则化：在渲染UV坐标上施加总变差损失，鼓励相邻渲染像素来自UV空间中相邻的高斯，提升几何保真度和训练稳定性。(3) 高效渲染：3DGS的tile-based光栅化器实现高分辨率实时渲染。采样方案可在训练中改变以调整高斯数量。其设计目的不是单独增加一个模块，而是用上述表示、监督或渲染约束消除原流程在目标任务中的具体失配。将3D高斯泼溅引入3D GAN框架，通过2D CNN在模板网格UV空间预测高斯属性，利用UV布局的规律性大幅简化非结构化高斯预测。引入渲染UV坐标的总变差损失提升几何保真度。首次实现1024\^{}2分辨率的实时3D一致头部生成和渲染。适用时还应保留论文给出的边界：依赖模板网格的质量和UV展开；对极端表情和姿态的泛化能力有限。
% METHOD: 2408.06019
\paragraph{HeadGAP（少样本头部高斯先验）（2024.08）}
HeadGAP~\cite{2408_06019} 的目标是从少量野外图像快速创建可动画头部。它先在大规模多视角动态数据上学习 3D Gaussian head prior，采用 auto-decoder 架构：身份共享编码、点特定特征和部分 MLP 回归高斯属性，CNN 精炼渲染结果。少样本个性化时，通过反演和微调适配新身份。相对 GGHead 的 GAN 采样，HeadGAP 更偏“可拟合先验”：先学一般头部动态，再用少量图像把先验拉向具体人物。
\textbf{实现细节：} HeadGAP（2024.08） 与 SEGA 的静态-动态双分支解耦思路相似，但采用了不同的先验学习范式。SEGA 用 DINOv2 编码器从 2D 图像直接提取身份特征再经 Transformer 映射到 UV 空间；HeadGAP 则构建基于自动解码器（auto-decoder）的部件级特征点云表示（part-based feature point cloud），从 FLAME 网格 UV 空间初始化 $n$ 个高斯基元，每个语义部件（如面部、头皮、耳朵、口腔）拥有独立的身份码 $z_j \in \mathbb{R}^{c_2}$ 和部件专属 MLP $f^M_p$，而非 SEGA 的单一全局编码器-解码器架构。这一架构选择的设计动机在于：auto-decoder 无需显式编码器即可在多身份联合优化中自然形成身份潜空间，使先验学习不依赖大规模配对数据，且通过共享点特征编码 $f$ 跨身份传递通用几何先验，同时以部件级身份码 $z_j$ 保留各语义区域的独立表达能力——面部和头皮的身份特征解耦度更高，更容易捕获发型等大变异部件的身份信息。
\textbf{方法信息：} ByteDance 参与了这篇论文。
\textbf{设计动机与边界：} 相对原始 3DGS 以独立高斯和图像重建误差为主的设定，HeadGAP 把问题落在以下可实现的改动上：两阶段框架：先验学习阶段使用GAPNet从大规模多视角动态数据学习3D高斯头部先验。GAPNet采用自动解码器架构：初始化部分特征点云（部分身份码+点特定特征编码），通过部分MLP回归高斯属性，CNN模块精炼渲染结果。使用FLAME网格跟踪先验，高斯属性相对于跟踪网格预测。少样本个性化阶段通过反演和微调策略快速适配新身份。部分动态建模使每个点只负责一个语义部分。其设计目的不是单独增加一个模块，而是用上述表示、监督或渲染约束消除原流程在目标任务中的具体失配。提出基于3DGS自动解码器网络的头部先验学习框架HeadGAP，通过部分动态建模学习可泛化的3D高斯头部先验，利用身份共享编码和个性化潜码，实现从少量野外图像快速创建高保真、可动画的3D头部虚拟形象。适用时还应保留论文给出的边界：先验学习需要大规模多视角动态数据；个性化质量受FLAME网格跟踪精度影响；极端表情可能超出先验范围
% METHOD: 2409.16147
\paragraph{Gaussian Deja-vu（先泛化再个性化的快速头像）（2024.09）}
Gaussian Deja-vu~\cite{2409_16147} 也走先验加个性化路线，但它用 Rectification Blendmaps 替代测试时神经网络。泛化模型基于 FLAME 构建 UV Gaussian map，并在大规模合成和真实 2D 图像上训练。个性化时，给定单目视频，学习表情感知修正混合图来调整高斯的位置、尺度、旋转、透明度和颜色。Blendmaps 是 UV 空间 2D 图，可按表情系数动态修正。相对 HeadGAP，它的增量是个性化更轻，减少长时间 per-subject 训练。
\textbf{实现细节：} Gaussian Deja-vu（2024.09） 与 SEGA 在 UV 空间参数化高斯的思路一致，但采用了"先通用化后个性化"的两阶段流程而非 SEGA 的单阶段前馈预测。这一流程的设计动机在于：通用化阶段通过大规模数据学习到的单图重建能力提供了稳健的初始化，使后续个性化阶段仅需在此基础上做小幅修正，而非从零开始优化整个 UV 高斯图——这与 HeadGAP 利用身份码本反演作为初始化的思路异曲同工，都是"先获取合理起点再精修"的策略。在 UV 参数化上，SEGA 使用标准 FLAME UV 对应关系；Gaussian Deja-vu 修改了 FLAME UV 映射以引入口腔内部形状并移除颈部区域。这一修改的设计考量在于：标准 FLAME UV 缺少口腔内部的几何表示，导致牙齿和舌头无法在 UV 空间中被编码；而颈部区域在头部动画中运动模式与面部差异过大，将其排除可减少无关区域对 UV 图空间的浪费，使有限的 UV 分辨率集中用于面部和头皮细节。
\textbf{方法信息：} Huawei 参与了这篇论文。
\textbf{设计动机与边界：} 相对原始 3DGS 以独立高斯和图像重建误差为主的设定，Gaussian Deja-vu 把问题落在以下可实现的改动上：方法分为两阶段：(1)泛化模型：基于FLAME模型构建UV高斯图表示(U\_gaussian R\^{}{H*W*13})，在大规模合成和真实2D图像数据集上训练，得到通用的3D高斯头部模型；(2)个性化：给定单目视频，使用可学习的表情感知修正混合图(Rectification Blendmaps)修正初始3D高斯的位置、缩放、旋转、alpha和颜色参数。Blendmaps是UV空间的可学习2D图，根据表情系数动态调整高斯参数，无需神经网络即可快速收敛。使用0阶球谐函数表示颜色，简化表示。其设计目的不是单独增加一个模块，而是用上述表示、监督或渲染约束消除原流程在目标任务中的具体失配。提出Gaussian Deja-vu框架，通过先泛化再个性化的策略大幅加速3DGS头部虚拟形象创建。引入可学习的表情感知修正混合图(Rectification Blendmaps)，无需神经网络即可修正初始3D高斯，训练时间降至现有方法的四分之一以下，同时渲染质量超越SOTA。适用时还应保留论文给出的边界：依赖FLAME模型的拓扑约束，对极端表情和复杂发型的建模可能受限；个性化阶段仍需单目视频输入。
% METHOD: 2412.04955
\paragraph{MixedGaussianAvatar（2D-3D 混合高斯头像）（2024.12）}
MixedGaussianAvatar~\cite{2412_04955} 解决 3DGS 渲染好但几何弱、2DGS 几何好但外观弱的矛盾。它先把 2D Gaussians 附着到 FLAME 三角网格上，用 ray-splat intersection 保证多视图一致几何；再在 2DGS 渲染不足处绑定额外 3D Gaussians，增强颜色和非表面细节。训练先优化 2D 高斯建立几何基础，再微调整个 2D-3D 混合系统。相对 SplatFace 的表面联合优化，它更直接地把 2DGS 作为几何底座，把 3DGS 作为渲染补丁。
\textbf{实现细节：} MixedGaussianAvatar（2024.12） 针对这一盲点提出了混合 2D-3D 高斯表示，与 NPGA 纯 3D 高斯的策略形成互补。改了什么：该方法放弃 NPGA 的逐高斯潜在特征和 MLP 动力学模块，转而采用两阶段混合表示。先用 2D 高斯（surfels）贴附于 FLAME 网格三角形面片重建几何表面，再通过基于瓦片损失的自适应筛选机制识别渲染质量不足的区域并附加 3D 高斯补偿外观。具体而言，渲染图像与真实图像被分割为 $w \times h$ 大小的瓦片，计算逐瓦片 MSE 损失 $L_{ij} = \|o_{ij} - x_{ij}\|_2^{}2$，选取损失最大的前 $k$ 个瓦片并利用高斯树结构回溯其关联的 2D 高斯父节点（同一三角形面片上的兄弟节点一并纳入），经 $n$ 次不同 FLAME 参数和相机参数下的迭代取并集 $U = \bigcup_{i=1}^{n} S'_i$ 后，为 $U$ 中的每个 2D 高斯生成一个继承其属性的 3D 高斯子节点。混合渲染：两种高斯通过局部-全局坐标变换统一由 FLAME 参数驱动，其中位置变换公式为 $\mu_g^{2D} = R \mu_l^{2D} + T + p^{2D}_*(t)$（$T$ 为三角形三顶点平均位置，$p^{2D}_*(t)$ 为可学习扰动项以自适应修正位置误差）。
\textbf{设计动机与边界：} 相对原始 3DGS 以独立高斯和图像重建误差为主的设定，MixedGaussianAvatar 把问题落在以下可实现的改动上：核心设计：(1) 2D高斯基础：将2D高斯(去除3D高斯的第三维尺度)附着于FLAME头部模型的三角网格，利用2DGS的射线-面片交叉方法确保多视图一致的几何精度。(2) 3D高斯增强：识别2DGS渲染质量不足的2D高斯，在其上绑定额外的3D高斯构建混合2D-3D高斯。3D高斯提供更丰富的颜色表征弥补2DGS的渲染短板。(3) FLAME动画驱动：通过局部到全局变换方法，利用FLAME参数驱动混合2D-3D高斯创建动态表情。(4) 渐进训练：先训练2D高斯建立几何基础，再微调整个混合2D-3D高斯系统。为了消除原始 3DGS 在该目标任务中的具体失配，方法把上述表示、监督或渲染约束作为可学习和可执行的改动，而不是单独堆叠模块。提出混合2D-3D高斯表征解决3DGS渲染质量高但几何差、2DGS几何好但渲染差的问题。将2D高斯附着于FLAME三角网格确保几何精度，在2DGS渲染不足处附加3D高斯增强渲染质量。渐进训练策略先训练2D后微调混合表征。用统一混合高斯表征整合2D图像和3D网格两种模态。适用时还应保留论文给出的边界：FLAME模型对极端表情的表达能力有限；2D-3D高斯的绑定策略可能不适合所有面部区域；训练流程较复杂。
\textbf{方法信息：} Intel 参与了这篇论文。



% METHOD: 2412.19149
\paragraph{3D Gaussian GAN Head Avatars（可编辑 GAN 头像）（2024.12）}
Generating Editable Head Avatars with 3D Gaussian GANs~\cite{2412_19149} 把 3D-aware GAN 和 3DGS 结合。传统 NeRF-GAN 可控性差，3DGS 又有尺度与距离混淆、点云无结构等歧义；该方法用 FLAME 约束面部区域，形成 Editable Gaussian Head，通过 albedo map 和 bump map 控制颜色和细粒度几何。头发等非面部区域用有限空间变化的高斯点云和 tri-plane 特征。相对 GGHead 的 UV 空间生成，它更突出编辑性：纹理和表情可由 3DMM 语义直接控制。
\textbf{实现细节：} EG2-3D框架将共享点云分为两组：(1) 面部区域使用Editable Gaussian Head(EG-Head)，基于FLAME 3DMM模型生成面部几何，albedo贴图控制颜色属性，bump贴图表示细粒度几何偏移，实现直观的纹理编辑和表情动画；(2) 头发等非面部区域使用有限空间变化的3DGS点云结合tri-plane特征表示。两组点云合并后通过高斯泼溅高分辨率渲染。为解决3DGS在GAN框架中的歧义性（尺度与距离的混淆、点云的无结构性质），利用3DMM的opacity属性表示点的可见性。光照控制使用screen-space缓冲区结合SH光照和轻量CNN。 提出 EG2-3D，首个在3D感知GAN框架中利用3DGS生成可编辑、可动画化3D头部的方法。通过EG-Head解耦高斯点的位置与属性，结合albedo贴图和bump贴图实现灵活纹理编辑，在保持身份特征的同时实现最先进的可控性。 头发区域的建模仍受限于tri-plane表示的表达能力；训练需要大量计算资源；对于极端姿态的渲染质量可能下降。
\textbf{方法信息：} Alibaba 参与了这篇论文。
\textbf{设计动机与边界：} 相对原始 3DGS 以独立高斯和图像重建误差为主的设定，Generating Editable Head Avatars with 3D Gaussian GANs 把问题落在以下可实现的改动上：EG\^{}2\^{}3D框架将共享点云分为两组：(1) 面部区域使用Editable Gaussian Head(EG-Head)，基于FLAME 3DMM模型生成面部几何，albedo贴图控制颜色属性，bump贴图表示细粒度几何偏移，实现直观的纹理编辑和表情动画；(2) 头发等非面部区域使用有限空间变化的3DGS点云结合tri-plane特征表示。两组点云合并后通过高斯泼溅高分辨率渲染。为解决3DGS在GAN框架中的歧义性（尺度与距离的混淆、点云的无结构性质），利用3DMM的opacity属性表示点的可见性。光照控制使用screen-space缓冲区结合SH光照和轻量CNN。其设计目的不是单独增加一个模块，而是用上述表示、监督或渲染约束消除原流程在目标任务中的具体失配。提出EG\^{}2\^{}3D，首个在3D感知GAN框架中利用3DGS生成可编辑、可动画化3D头部的方法。通过EG-Head解耦高斯点的位置与属性，结合albedo贴图和bump贴图实现灵活纹理编辑，在保持身份特征的同时实现最先进的可控性。适用时还应保留论文给出的边界：头发区域的建模仍受限于tri-plane表示的表达能力；训练需要大量计算资源；对于极端姿态的渲染质量可能下降。
% METHOD: 2502.19782
\paragraph{Open-Vocabulary Human Part Segmentation（人体部件开放词汇分割）（2025.02）}
Open-Vocabulary Semantic Part Segmentation of 3D Human~\cite{2502_19782} 不直接生成头像，而是给 3D 人体表示提供开放词汇部件分割。方法先从多个渲染视角用 SAM 产生 mask proposals，再用 HumanCLIP 获取更适合人体文本对齐的嵌入，最后用 MaskFusion 将多视图掩码和文本提示聚合成 3D 语义掩码。它支持多种 3D 表示，包括高斯人体。相对 COSY 或 LAGA 的编辑生成，它解决的是编辑前的可定位问题：读者要能用文本指定“头发”“袖口”“眼镜”这类人体部件。
\textbf{实现细节：} 方法包含三个核心模块：1）多视图掩码提议：利用SAM在预定义相机视角渲染的图像上生成类别无关的重叠掩码提议；2）HumanCLIP编码：提出专门针对人体内容的CLIP变体，在统一CLIP特征空间中为每个掩码生成嵌入，相比原始CLIP产生更准确的文本对齐嵌入；3）MaskFusion模块：将HumanCLIP编码的文本提示与掩码嵌入同时进行分类和多视图聚合，无需复杂的投票和分组机制。掩码提议和嵌入只需生成一次，后续分割可在毫秒级完成。 提出了首个面向3D人体的开放词汇分割方法，设计了HumanCLIP模型生成更准确的人体相关嵌入，以及简单高效的MaskFusion模块实现多视图特征到3D语义掩码的分类与融合，支持多种3D表示形式。 依赖预定义相机视角的渲染质量；对极端姿态的人体可能效果下降；HumanCLIP需要额外训练数据。
\textbf{方法信息：} Qualcomm 参与了这篇论文。
\textbf{设计动机与边界：} 相对原始 3DGS 以独立高斯和图像重建误差为主的设定，Open-Vocabulary Semantic Part Segmentation of 3D Human 把问题落在以下可实现的改动上：方法包含三个核心模块：1）多视图掩码提议：利用SAM在预定义相机视角渲染的图像上生成类别无关的重叠掩码提议；2）HumanCLIP编码：提出专门针对人体内容的CLIP变体，在统一CLIP特征空间中为每个掩码生成嵌入，相比原始CLIP产生更准确的文本对齐嵌入；3）MaskFusion模块：将HumanCLIP编码的文本提示与掩码嵌入同时进行分类和多视图聚合，无需复杂的投票和分组机制。掩码提议和嵌入只需生成一次，后续分割可在毫秒级完成。其设计目的不是单独增加一个模块，而是用上述表示、监督或渲染约束消除原流程在目标任务中的具体失配。提出了首个面向3D人体的开放词汇分割方法，设计了HumanCLIP模型生成更准确的人体相关嵌入，以及简单高效的MaskFusion模块实现多视图特征到3D语义掩码的分类与融合，支持多种3D表示形式。适用时还应保留论文给出的边界：依赖预定义相机视角的渲染质量；对极端姿态的人体可能效果下降；HumanCLIP需要额外训练数据。
% METHOD: 2504.10486
\paragraph{DNF-Avatar（神经场蒸馏到可重光照高斯）（2025.04）}
DNF-Avatar~\cite{2504_10486} 的锚点是 IntrinsicAvatar 一类高质量神经 SDF 与 PBR 场。它的改变是教师-学生蒸馏：教师用神经 SDF 和物理外观场通过 Monte Carlo 光线追踪得到高质量结果，学生用关节化 2DGS 与 split-sum 近似替代昂贵光追。每个身体部位设置环境光遮蔽探针，像素只需一次查询即可估计阴影；学生再从教师蒸馏几何和外观属性。骨骼变形采用 SMPL 与逆 LBS。相对普通头像高斯只做颜色重建，DNF-Avatar 把重光照和 PBR 外观压缩到实时高斯学生模型中。
\textbf{实现细节：} DNF-Avatar采用教师-学生知识蒸馏架构：(1) 教师模型：基于IntrinsicAvatar，使用神经SDF场和PBR外观场，通过蒙特卡洛光线追踪渲染，质量高但速度慢。(2) 学生模型：使用关节化2DGS表征，扩展外观模型引入近似PBR模型（split-sum近似），避免光线追踪。提出新颖的部件级环境光遮蔽探针（Part-wise Ambient Occlusion Probes）：每个身体部位设置探针，每个像素仅需一次查询即可预测阴影。学生模型通过从教师模型蒸馏学习几何和外观属性。骨骼变形基于SMPL模型和逆线性混合蒙皮（Fast-SNARF）。 提出从隐式神经场（教师）蒸馏到2D高斯泼溅（学生）的知识蒸馏框架，结合split-sum近似和部件级环境光遮蔽探针，实现67FPS实时重光照。 依赖教师模型的质量；重光照质量受split-sum近似精度限制；对极端光照条件（如强聚光灯）的建模有限。
\textbf{实现补充：} 实现上还要核对输入、输出、参数、特征、初始化、训练损失、正则、渲染或推理接口，以及优化中的剪枝、分裂、采样、编码或解码步骤；这些抓手决定该方法到底改了 3DGS 管线的哪一环。
\textbf{方法信息：} Intel 参与了这篇论文。
\textbf{设计动机与边界：} 相对原始 3DGS 以独立高斯和图像重建误差为主的设定，DNF-Avatar 把问题落在以下可实现的改动上：DNF-Avatar采用教师-学生知识蒸馏架构：(1) 教师模型：基于IntrinsicAvatar，使用神经SDF场和PBR外观场，通过蒙特卡洛光线追踪渲染，质量高但速度慢。(2) 学生模型：使用关节化2DGS表征，扩展外观模型引入近似PBR模型（split-sum近似），避免光线追踪。提出新颖的部件级环境光遮蔽探针（Part-wise Ambient Occlusion Probes）：每个身体部位设置探针，每个像素仅需一次查询即可预测阴影。学生模型通过从教师模型蒸馏学习几何和外观属性。骨骼变形基于SMPL模型和逆线性混合蒙皮（Fast-SNARF）。其设计目的不是单独增加一个模块，而是用上述表示、监督或渲染约束消除原流程在目标任务中的具体失配。提出从隐式神经场（教师）蒸馏到2D高斯泼溅（学生）的知识蒸馏框架，结合split-sum近似和部件级环境光遮蔽探针，实现67FPS实时重光照。适用时还应保留论文给出的边界：依赖教师模型的质量；重光照质量受split-sum近似精度限制；对极端光照条件（如强聚光灯）的建模有限。
% METHOD: 2504.12999
\paragraph{GSAC（从单目视频到 Unity 的头像管线）（2025.04）}
GSAC~\cite{2504_12999} 是工程管线型方法。它继承 ExAvatar 和 SMPL-X 绑定高斯的路线，但把重点放在从单目视频到 Unity 集成。预处理阶段用 DECA 提取 FLAME 参数、Hand4Whole 估计 SMPL-X，并用 mmpose 检测关键点；手部不可见时，根据运动学链和角速度估计缺失手部位置，改善 SMPL-X 初始化。头像训练侧基于 3DGS 生成绑定骨架的写实化身，Unity 编辑器负责交互和 VR/AR 接入。项目实现更像端到端制作工具，而不是单一表示创新。
\textbf{实现细节：} GSAC包含三大组件：(1) 定制数据预处理：从单目视频提取帧，使用DECA提取FLAME面部参数，Hand4Whole估计SMPL-X参数，mmpose检测关键点。创新性地解决了手部不可见时的关键点预测问题：利用运动学链和角速度粗略估计缺失手部位置，改善SMPL-X参数初始化。(2) 头像训练：基于ExAvatar的方法，使用3DGS训练绑定了SMPL-X骨骼的写实化身，支持详细面部表情重建。定制预处理支持野外视频输入。(3) Unity集成：提供Unity编辑器中的高斯泼溅头像编辑工具，支持VR/AR应用开发。用户友好的环境降低了技术门槛。 提出从单目视频到Unity游戏引擎的端到端3DGS头像创建管线，包含定制预处理、表情重建和Unity集成编辑器。 单目视频质量限制最终头像质量；复杂服装和头发建模有限；Unity集成中高斯渲染的性能开销。
\textbf{方法信息：} 代码：\url{https://github.com/vu-rasl/gsac}。
\textbf{设计动机与边界：} 相对原始 3DGS 以独立高斯和图像重建误差为主的设定，GSAC 把问题落在以下可实现的改动上：GSAC包含三大组件：(1) 定制数据预处理：从单目视频提取帧，使用DECA提取FLAME面部参数，Hand4Whole估计SMPL-X参数，mmpose检测关键点。创新性地解决了手部不可见时的关键点预测问题：利用运动学链和角速度粗略估计缺失手部位置，改善SMPL-X参数初始化。(2) 头像训练：基于ExAvatar的方法，使用3DGS训练绑定了SMPL-X骨骼的写实化身，支持详细面部表情重建。定制预处理支持野外视频输入。(3) Unity集成：提供Unity编辑器中的高斯泼溅头像编辑工具，支持VR/AR应用开发。用户友好的环境降低了技术门槛。其设计目的不是单独增加一个模块，而是用上述表示、监督或渲染约束消除原流程在目标任务中的具体失配。提出从单目视频到Unity游戏引擎的端到端3DGS头像创建管线，包含定制预处理、表情重建和Unity集成编辑器。适用时还应保留论文给出的边界：单目视频质量限制最终头像质量；复杂服装和头发建模有限；Unity集成中高斯渲染的性能开销。
% METHOD: 2504.14373
\paragraph{SEGA（单图可驱动高斯头部）（2025.04）}
SEGA~\cite{2504_14373} 从单张图像生成可驱动 3D Gaussian head avatar。它继承 FLAME 驱动的人脸路线，但不只依赖参数模型，而是把 2D 身份先验和 3D 多视角先验结合起来。静态分支用 DINOv2 与 LRM 生成全头高斯资产，保留身份、头发和耳朵等相对稳定部分；动态分支基于 FLAME，把嘴部和眼睛这类表情区域放到多分辨率哈希编码中更新。两支最后融合并用输入图像微调，静态组件可以预计算，人物特定微调再提升身份保真度。相对 3D Gaussian Blendshapes，它更偏单图前馈生成；相对 FastAvatar，它的关键是静态组件和动态表情组件在 UV 空间分层解耦。
\textbf{实现细节：} SEGA（2025.04） 将重建问题拆解为"静态身份"和"动态表情"两个独立子问题。FaceLift 用扩散模型生成多视角再重建，SEGA 则完全绕过多视角生成：静态分支用 DINOv2（自监督视觉Transformer，通过自蒸馏学习强大视觉特征）编码器提取 2D 身份特征，经 UV-Alignment Transformer（10 层）映射到 FLAME UV 空间，再由静态解码器直接预测全头部高斯属性（$1024 \times 1024$ UV 分辨率），其中位置偏移量 $M_{\mathrm{offset}}$ 捕捉身份特有的几何细节（发型、面部轮廓等）；动态分支用 VQ-VAE（Vector Quantized Variational AutoEncoder，向量量化变分自编码器）编码器提取表情特征，通过位移 VAE 预测表情驱动的形变图（$400 \times 400$ UV 分辨率），仅覆盖嘴唇、眉毛等表情敏感区域。推理时静态分支仅计算一次并缓存，动态分支逐帧更新，通过预定义二值掩码 $M_d$ 将两路输出在 UV 空间中融合，并用线性插值权重掩码确保过渡区域无缝拼接。
\textbf{方法信息：} ByteDance 参与了这篇论文。
\textbf{设计动机与边界：} 相对原始 3DGS 以独立高斯和图像重建误差为主的设定，SEGA 把问题落在以下可实现的改动上：SEGA包含三个核心部分：(1)静态分支：结合DINOv2骨干（2D身份一致性先验）和大型3D重建模型LRM（3D几何一致性先验），生成全头3D高斯资产以忠实保留身份信息；(2)动态分支：基于FLAME结构先验，使用双分支架构将静态（表情不变区域如头发、耳朵）和动态（表情驱动区域如嘴部、眼睛）面部组件解耦，动态分支使用多分辨率哈希编码实现表情驱动的精细细节更新；(3)混合阶段：融合两个分支输出并使用输入图像进一步精调。静态组件可预计算，确保实时性能。还包含人物特定微调以增强保真度。其设计目的不是单独增加一个模块，而是用上述表示、监督或渲染约束消除原流程在目标任务中的具体失配。提出SEGA方法从单张图像生成可驱动3D高斯头部化身，结合2D大规模数据集先验和3D多视角先验，设计分层UV空间高斯溅射框架将静态和动态面部组件解耦，实现实时动画渲染。适用时还应保留论文给出的边界：依赖FLAME模型的精度；对极端姿态和夸张表情的泛化能力有待提升；单图像输入下背面信息可能不够准确。
% METHOD: 2507.18155
\paragraph{GeoAvatar（刚柔高斯预分配与嘴部结构）（2025.07）}
GeoAvatar~\cite{2507_18155} 继承 FLAME 加 3DGS 的头部 avatar 路线，但它发现不同面部区域对变形的需求不同：额头、耳朵等更接近刚性，嘴部和表情区域需要更强自由度。方法通过 adaptive pre-allocation stage 无监督地把高斯分成刚性和柔性集合，再分别施加不同 offset 正则。它还按嘴部解剖和运动设计专门结构，并用高斯到 3DMM 的绑定正则稳定动画。项目实现包含 DynamicFace 数据集和部件级变形。相对普通 FLAME 绑定，它的增量是先判断高斯该“硬”还是“软”，再驱动。
\textbf{实现细节：} GeoAvatar（2025.07） 从正则化粒度的角度补充了 ScaffoldAvatar 的设计。锚点回顾：ScaffoldAvatar 的所有锚点采用统一的 MLP 架构和优化策略，对不同几何复杂度的面部区域施加同等强度的约束——但面部并非各区域同等刚性：鼻梁应紧贴网格，而头发和嘴角需要更大的偏移自由度。改了什么：GeoAvatar 引入自适应预分配阶段（Adaptive Pre-allocation Stage, APS），在初始训练 $N$ 步后根据各面部区域高斯到网格的平均距离 $\mu^{(k)} = \frac{1}{|G_k|} \sum_{i,j \in G_k} \| \mu_{i,j} \|$ 将 FLAME 三角面片无监督地分割为刚性集 $F_r$（$\mu^{(k)} < \bar\mu$）和柔性集 $F_f$（$\mu^{(k)} > \bar\mu$），对两组施加差异化的位置正则 $\delta_r=0.1$（刚性，严格约束高斯贴近网格）和 $\delta_f=2.0$（柔性，允许较大偏移以建模头发等非刚性区域）。此外，GeoAvatar 还引入角度正则 $L_{\mathrm{angle}}(\theta) = \mathbf{1}_{r > \sigma_r}(\theta - \theta_0)$（$\theta_0 = 45°$），防止柔性集中距离网格较远的高斯偏移到错误的相邻面片区域。
\textbf{实现补充：} 实现上还要核对输入、输出、参数、特征、初始化、训练损失、正则、渲染或推理接口，以及优化中的剪枝、分裂、采样、编码或解码步骤；这些抓手决定该方法到底改了 3DGS 管线的哪一环。
\textbf{方法信息：} 代码：\url{https://github.com/seungjun-moon/geoavatar}。
\textbf{设计动机与边界：} 相对原始 3DGS 以独立高斯和图像重建误差为主的设定，GeoAvatar 把问题落在以下可实现的改动上：基于FLAME参数模型和3DGS构建头部Avatar系统。核心创新：(1)自适应预分配阶段(APS)，无监督地将高斯分割为刚性和柔性集合，实现自适应偏移正则化；(2)基于嘴部解剖学和动力学设计新型嘴部结构和部件级变形策略，增强嘴部动画保真度；(3)提出精确的高斯与3DMM面部绑定正则化损失。发布DynamicFace视频数据集，包含高表情面部运动。其设计目的不是单独增加一个模块，而是用上述表示、监督或渲染约束消除原流程在目标任务中的具体失配。提出自适应几何高斯泼溅框架GeoAvatar，通过自适应预分配阶段(APS)将高斯分为刚性和柔性集合实现自适应偏移正则化，设计新型嘴部结构和部件级变形策略，并发布DynamicFace高表情数据集。适用时还应保留论文给出的边界：依赖FLAME追踪器的精度；极端表情泛化能力待验证；实时渲染性能未明确报告。
% METHOD: 2508.18389
\paragraph{FastAvatar-Template（任意姿态单图三秒头像）（2025.08）}
另一篇 FastAvatar~\cite{2508_18389} 更强调从单张任意姿态照片快速生成 3DGS 人脸头像。它先在 FLAME 三角面中心放置高斯，对多个对象优化后取平均，得到数据驱动 3DGS template face model。前馈阶段，编码器把输入图像映射到姿态无关身份嵌入，解码器预测模板高斯的残差位置、不透明度、尺度、旋转和外观；测试时只轻量优化 SH 外观参数。相对 LGRT 版 FastAvatar 的多输入聚合，这里更像“模板加残差”的单图推断路线，用几何先验换取极快重建。
\textbf{实现细节：} FastAvatar（2025.08） 实现了约 3 秒内的单图像 3D 人脸重建，与 FaceLift 的两阶段扩散+重建策略形成对比。FaceLift 依赖扩散模型先"想象"不可见视角再重建，虽能处理任意视角但生成过程较慢且可能引入扩散模型的伪影。FastAvatar 的核心思路转变是：与其让模型从零"脑补"缺失视角，不如直接利用数据驱动的平均模板作为强先验，只需预测个体与模板之间的偏差即可。与锚点的核心差异：FaceLift 通过扩散模型在 2D 图像空间生成缺失视角、再由 Transformer 重建 3DGS；FastAvatar 完全绕过视角生成步骤，转而在 3D 参数空间建立显式模板先验，通过残差预测直接从单张图像回归完整 3DGS 模型。模板构建：对 $N=410$ 个训练身份各自优化一个 3DGS 模型（每个高斯放置在 FLAME 网格三角面片中心，加上 168 个牙齿锚点，共 $K=10{,}144$ 个语义一致的高斯），然后对对应高斯的参数取平均得到模板 $T$。这种简单平均策略相比联合优化模板或随机初始化能产生更稳定的先验。残差预测：前馈编码器-解码器不直接回归绝对参数值，而是预测相对于模板的残差。
\textbf{方法信息：} Samsung 参与了这篇论文。
\textbf{设计动机与边界：} 相对原始 3DGS 以独立高斯和图像重建误差为主的设定，FastAvatar 把问题落在以下可实现的改动上：方法采用两阶段设计：(1)模板构建：受经典形变模型启发，在FLAME网格每个三角面中心放置高斯球，对多个对象优化3DGS模型后取平均构建模板T，提供几何先验；(2)前馈预测：编码器将输入图像映射到姿态无关的身份嵌入，解码器预测每个高斯的残差参数（位置偏移、不透明度、缩放、旋转、SH外观残差）将模板变形为特定对象模型；(3)轻量测试时精化：优化外观参数（SH系数）实现照片级渲染。支持FLAME引导的表情动画实现可控重演。其设计目的不是单独增加一个模块，而是用上述表示、监督或渲染约束消除原流程在目标任务中的具体失配。提出FastAvatar，从单张任意姿态输入图像在约3秒内重建高质量3DGS人脸头像。两阶段设计：前馈编码器-解码器从姿态不变身份嵌入预测粗糙面部几何（高斯结构），轻量测试时精化优化外观参数。构建数据驱动3DGS模板面模型作为几何先验，比现有逐对象优化方法(FlashAvatar、GaussianAvatars)快600倍以上，PSNR 24.01dB、SSIM 0.91。适用时还应保留论文给出的边界：模板的泛化能力受限于训练数据多样性；对极端表情和配饰（如眼镜）的处理有限；3秒重建时间仍非严格实时。
% METHOD: 2508.19754
\paragraph{FastAvatar-LGRT（多输入统一头像重建）（2025.08）}
FastAvatar~\cite{2508_19754} 面向单图、多视角或单目视频的统一快速头像重建。它的局部锚点是大型高斯重建 Transformer，而不是逐人优化。系统先提取面部特征和相机、表情、头姿等多粒度条件，再由 3DGS Transformer 聚合多帧信息，预测规范空间中的 registered 3DGS。增量高斯聚合利用关键点追踪和切片融合损失，使更多输入带来更好质量。项目实现围绕 FastAvatar 和 LGRT 展开。相对 SEGA 的单图静动分支，FastAvatar 更强调输入数量可变和质量-速度可调。
\textbf{实现细节：} FastAvatar的核心是大型高斯重建Transformer（LGRT），包含6个阶段：面部特征提取、面部编码、面部聚合与注册、3DGS属性生成、规范3DGS模型融合、3DGS光栅化。三个关键设计：(1) 3DGS Transformer聚合多帧信息并注入初始3D提示来预测注册的规范3DGS表示；(2) 多粒度引导编码（相机姿态、表情系数、头部姿态），通过轻量MLP处理缓解动画引起的错位，支持变长输入；(3) 增量高斯聚合：通过关键点追踪和切片融合损失实现增量重建，即更多观测带来更高质量而不浪费输入数据。模型输出包括颜色、不透明度、缩放因子、旋转、重要性分数和点偏移等3D高斯属性。 提出FastAvatar框架，基于大型高斯重建Transformer（LGRT），可在数秒内从任意数量的输入（单图/多视角/视频）重建高质量3D高斯头像。核心创新包括：3DGS Transformer注入初始3D提示、多粒度引导编码、增量高斯聚合策略，实现了质量-速度可调节的头像建模范式。 最大输入帧数受GPU容量限制（通常16帧）；依赖FLAME参数化模型进行粗略姿态估计；增量重建可能引入累积误差
\textbf{方法信息：} 代码：\url{https://github.com/tyrionwuyue/fastavatar}。
\textbf{设计动机与边界：} 相对原始 3DGS 以独立高斯和图像重建误差为主的设定，FastAvatar 把问题落在以下可实现的改动上：FastAvatar的核心是大型高斯重建Transformer（LGRT），包含6个阶段：面部特征提取、面部编码、面部聚合与注册、3DGS属性生成、规范3DGS模型融合、3DGS光栅化。三个关键设计：(1) 3DGS Transformer聚合多帧信息并注入初始3D提示来预测注册的规范3DGS表示；(2) 多粒度引导编码（相机姿态、表情系数、头部姿态），通过轻量MLP处理缓解动画引起的错位，支持变长输入；(3) 增量高斯聚合：通过关键点追踪和切片融合损失实现增量重建，即更多观测带来更高质量而不浪费输入数据。模型输出包括颜色、不透明度、缩放因子、旋转、重要性分数和点偏移等3D高斯属性。其设计目的不是单独增加一个模块，而是用上述表示、监督或渲染约束消除原流程在目标任务中的具体失配。提出FastAvatar框架，基于大型高斯重建Transformer（LGRT），可在数秒内从任意数量的输入（单图/多视角/视频）重建高质量3D高斯头像。核心创新包括：3DGS Transformer注入初始3D提示、多粒度引导编码、增量高斯聚合策略，实现了质量-速度可调节的头像建模范式。适用时还应保留论文给出的边界：最大输入帧数受GPU容量限制（通常16帧）；依赖FLAME参数化模型进行粗略姿态估计；增量重建可能引入累积误差
% METHOD: 2511.17747
\paragraph{AEGIS（高斯人脸头像隐私保护）（2025.11）}
AEGIS~\cite{2511_17747} 关注 3D Gaussian facial avatars 的隐私。它不修改几何，也不重训头像，而是作为后训练步骤扰动 SH 表示的 DC 系数，也就是高斯基色。对抗目标来自 ArcFace 和 AdaFace 人脸验证嵌入，优化通过可微渲染器反传到头像参数，使去身份化在多视角和动画中保持一致。项目实现围绕 adversarial perturbation。相对 COSY 的语义编辑，AEGIS 是安全层：它改变可识别身份，同时尽量保留年龄、表情、种族和性别等非身份属性。
\textbf{实现细节：} AEGIS（2025.11） 提出了面向 3D 高斯化身的隐私保护身份遮蔽框架。Baseline 与问题：基于 3DGS 的光真实 3D 面部化身日益普及，但这也带来了身份盗用风险——恶意方可以从在线共享的化身模型中提取面部特征，绕过生物识别认证系统。2D 图像领域已有成熟的对抗性遮蔽方法，但 3D 化身的动态性和多视角一致性要求使得直接迁移不可行：2D 扰动仅对特定视角有效，而 3D 高斯从任意角度渲染时扰动模式可能失效或产生可见伪影。改了什么：AEGIS 通过对抗性扰动 3D 高斯的颜色系数（球谐函数参数），在保持化身感知质量的同时实现跨视角一致的去识别。核心设计利用预训练的人脸验证网络作为对抗性损失的身份判别器，引导颜色扰动方向使验证网络无法匹配原始身份，但不修改化身的几何（位置、缩放、旋转、不透明度），从而保持面部结构和表情完整性。优化仅作用于高斯颜色系数，几何属性完全冻结。扰动后的高斯化身从任意视角渲染时，身份信息被遮蔽，但年龄、种族、性别、表情等非身份属性尽量保留；因此它更像后训练安全层，可接到已有 3D 高斯头像模型之后。
\textbf{方法信息：} 代码：\url{https://github.com/woleek/aegis}。
\textbf{设计动机与边界：} 相对原始 3DGS 以独立高斯和图像重建误差为主的设定，AEGIS 把问题落在以下可实现的改动上：AEGIS是一种后训练方法，仅扰动球谐表示的第一系数（DC系数，编码Gaussian原语的基色）。通过约束投影梯度下降(PGD)优化对人脸验证目标进行对抗优化，使用ArcFace和AdaFace嵌入作为验证目标。优化直接在头像参数空间进行并通过可微分渲染管线传播，确保去身份化在视角上一致、动画中稳定且感知连贯。设计三个核心组件：(1)Gaussian Avatars的可微分渲染器；(2)预训练人脸验证网络（内置人脸对齐预处理）；(3)对抗优化框架。参考身份使用原始头像的正面渲染。方法不修改头像几何或高阶视角依赖性。其设计目的不是单独增加一个模块，而是用上述表示、监督或渲染约束消除原流程在目标任务中的具体失配。提出AEGIS，首个3D Gaussian人脸头像的隐私保护身份掩码框架。通过对Gaussian颜色系数施加对抗扰动实现视角一致的去身份化，面部检索和验证准确率降至0%，同时保持高感知质量（SSIM=0.9555, PSNR=35.52dB）和关键面部属性（年龄、种族、性别、情感）。适用时还应保留论文给出的边界：对未知的验证网络可能泛化有限；强攻击者（如集成攻击）可能部分恢复身份；仅保护颜色信息，几何信息未保护。
% METHOD: 2512.02648
\paragraph{PoreTrack3D（毛孔级动态高斯基准）（2025.12）}
PoreTrack3D~\cite{2512_02648} 是基准型工作。它面向动态 3DGS 的毛孔级面部轨迹追踪，构建多视角面部运动序列和密集 3D 轨迹，包含传统关键点和毛孔级关键点。具体流程是使用多视角相机阵列采集 RGB 与精确几何，在面部表面提取密集关键点，并通过时序关联形成可评估的 3D 轨迹；评估阶段再把 Dynamic 3DGS、Deformable GS 等方法的输出与这些轨迹对齐，测量非刚性面部点能否跨时间稳定跟踪。相对 Topo4D 直接产出规则网格，PoreTrack3D 的价值在于评价：高保真头像不只看单帧 PSNR，还要看细粒度表面点能否跨时间稳定追踪。
\textbf{实现细节：} PoreTrack3D 的实现重点是数据集和评测协议，而不是新的头像渲染器。输入包括多视角同步面部视频、相机参数、重建几何、传统 landmark 与毛孔级关键点；输出是跨时间关联的 3D 轨迹、逐点跟踪误差和不同动态 3DGS 基线的指标表。它不新增高斯位置、尺度、旋转、不透明度或协方差参数化，而是把这些参数由 Dynamic 3DGS、Deformable GS 等方法训练后，再投影或对齐到密集面部轨迹上评价。
\textbf{实现补充：} 可复现流程应按采集、标注、训练和评测四层读：先用多相机阵列和精确几何建立面部表面，再提取毛孔级特征并做时序关联；随后运行不同动态头像或动态 3DGS 基线，保持相同输入视角、训练 split、优化步数和渲染设置；最后计算轨迹稳定性、点位误差、动态区域质量和失配案例。它的适用边界是高保真面部运动，不是通用人体或场景 benchmark；因此结论应落在表情、皮肤细节、非刚性运动和细粒度表面追踪上。
\textbf{方法信息：} 代码：\url{https://github.com/jhxion9/poretrack3d}。
\textbf{设计动机与边界：} 相对原始 3DGS 以独立高斯和图像重建误差为主的设定，PoreTrack3D 把问题落在以下可实现的改动上：基准数据集构建管道：(1)数据采集：使用多视角相机阵列捕获面部运动序列，同时记录RGB图像和精确的3D几何信息。(2)轨迹标注：在面部表面提取密集关键点(包括传统面部关键点和毛孔级精细关键点)，通过时序关联形成3D轨迹。(3)数据规模：超过44万条面部轨迹，其中5.2万+条长度超过10帧，68条经人工审核跨越全150帧(时间一致性的金标准)。(4)评估协议：系统评估SOTA动态3DGS方法(如Dynamic 3DGS、Deformable GS等)在轨迹追踪精度上的表现。建立首个毛孔级面部动态3D重建的性能基线。管道本身也为高保真面部运动捕获建立了新框架。其设计目的不是单独增加一个模块，而是用上述表示、监督或渲染约束消除原流程在目标任务中的具体失配。提出首个面向毛孔级非刚性3D面部轨迹追踪的动态3DGS基准PoreTrack3D。包含超过44万条面部轨迹，其中5.2万条长度超过10帧，68条经人工审核跨越全150帧。系统性评估SOTA动态3DGS方法，建立该领域首个性能基线。同时捕捉传统面部关键点和毛孔级关键点轨迹。适用时还应保留论文给出的边界：仅覆盖面部运动；轨迹标注可能存在噪声；未包含极端表情和遮挡场景；动态3DGS方法在该基准上的表现普遍有待提升。
% METHOD: 2512.15711
\paragraph{Gaussian Pixel Codec Avatars（网格与高斯混合编码）（2025.12）}
Gaussian Pixel Codec Avatars~\cite{2512_15711} 的局部锚点是 Pixel Codec Avatars。它指出皮肤这类表面区域适合网格和纹理，头发、胡须等非表面区域更适合 3D Gaussians。方法因此使用混合表示：网格分支生成面部网格和纹理，高斯分支生成非表面区域的各向异性高斯，两者共享 UV 卷积特征。渲染器把网格当作半透明层并入 3DGS 体渲染范式。相对 MeGA，它更偏压缩和移动端渲染效率；相对纯高斯头像，它用网格承担稳定表面，用高斯补非表面细节。
\textbf{实现细节：} GPiCA的核心是混合表示策略和统一渲染管线。三角网格高效表示面部皮肤等表面区域，3D高斯处理头发、胡须等非表面区域。开发统一可微渲染管线，将网格视为半透明层融入3DGS体渲染范式。训练变分自编码器，编码器从粗网格和平均纹理生成潜在码z，解码器分两个分支：Dm生成网格和纹理，Dg生成3D高斯参数，两者共享UV空间的卷积特征。网格采用半透明渲染与3D高斯联合渲染，所有组件在同一渲染引擎中同时渲染。方法基于PiCA构建以确保移动设备上的高效运行。 提出GPiCA混合表示方法，将三角网格与各向异性3D高斯结合，网格高效表示皮肤等表面区域，高斯处理毛发等非表面区域，开发统一可微渲染管线将网格视为半透明层融入3DGS体渲染范式，实现接近纯高斯头像的真实感和接近网格头像的渲染性能。 依赖多视角图像训练数据；需要粗网格跟踪预处理；动态光照适应性有限。
\textbf{方法信息：} Meta 参与了这篇论文。
\textbf{设计动机与边界：} 相对原始 3DGS 以独立高斯和图像重建误差为主的设定，Gaussian Pixel Codec Avatars 把问题落在以下可实现的改动上：GPiCA的核心是混合表示策略和统一渲染管线。三角网格高效表示面部皮肤等表面区域，3D高斯处理头发、胡须等非表面区域。开发统一可微渲染管线，将网格视为半透明层融入3DGS体渲染范式。训练变分自编码器，编码器从粗网格和平均纹理生成潜在码z，解码器分两个分支：Dm生成网格和纹理，Dg生成3D高斯参数，两者共享UV空间的卷积特征。网格采用半透明渲染与3D高斯联合渲染，所有组件在同一渲染引擎中同时渲染。方法基于PiCA构建以确保移动设备上的高效运行。其设计目的不是单独增加一个模块，而是用上述表示、监督或渲染约束消除原流程在目标任务中的具体失配。提出GPiCA混合表示方法，将三角网格与各向异性3D高斯结合，网格高效表示皮肤等表面区域，高斯处理毛发等非表面区域，开发统一可微渲染管线将网格视为半透明层融入3DGS体渲染范式，实现接近纯高斯头像的真实感和接近网格头像的渲染性能。适用时还应保留论文给出的边界：依赖多视角图像训练数据；需要粗网格跟踪预处理；动态光照适应性有限。
% METHOD: 2512.16893
\paragraph{Instant Expressive Gaussian Head Avatar（3D 感知表情蒸馏）（2025.12）}
Instant Expressive Gaussian Head Avatar~\cite{2512_16893} 继承单图前馈头像生成，但不依赖预定义参数化模型约束。它把 2D diffusion model 中学到的动画知识蒸馏到 3DGS 编码器：输入图像被编码为 triplanes，再从中采样高斯；运动基向量从 triplanes 检索，经紧凑 MLP 预测每个高斯的运动更新。局部融合替代全局注意力，降低推理开销。相对 SEGA 的 FLAME 动态分支，它的改变是动画表示与人脸 3D 表示解耦，从合成扩散数据中隐式学习运动。
\textbf{实现细节：} 方法核心是将2D扩散模型学到的动画知识蒸馏到3DGS前馈框架。架构包含：(1)即时编码器E基于LP3B改编，将输入图像编码为triplanes中间表示，再从中采样3D高斯基元；(2)解耦的动画表示——运动基向量mi从triplanes中检索，通过紧凑MLP预测逐高斯基元的运动更新，避免昂贵的全局注意力融合；(3)高效局部融合策略替代计算密集的全局融合机制；(4)从预训练扩散模型(zerox)渲染合成数据集用于训练。动画表示与人脸3D表示解耦，从数据中隐式学习运动，消除对预定义参数化模型的依赖。推理时仅需一次编码即可从任意驱动图像进行动画。 提出将2D扩散模型知识蒸馏到前馈编码器的方法，实现从单张图像即时转换为3D一致、快速且富有表现力的可动画表示。设计解耦动画表示和高效轻量级局部融合策略，动画速度达107.31FPS，无需预定义参数化模型约束。 动画表现力受限于蒸馏源模型；极端姿态下可能出现伪影；对输入图像质量敏感。
\textbf{方法信息：} NVIDIA 参与了这篇论文。
\textbf{设计动机与边界：} 相对原始 3DGS 以独立高斯和图像重建误差为主的设定，Instant Expressive Gaussian Head Avatar via 3D-Aware Expression Distillation 把问题落在以下可实现的改动上：方法核心是将2D扩散模型学到的动画知识蒸馏到3DGS前馈框架。架构包含：(1)即时编码器E基于LP3B改编，将输入图像编码为triplanes中间表示，再从中采样3D高斯基元；(2)解耦的动画表示——运动基向量mi从triplanes中检索，通过紧凑MLP预测逐高斯基元的运动更新，避免昂贵的全局注意力融合；(3)高效局部融合策略替代计算密集的全局融合机制；(4)从预训练扩散模型(zerox)渲染合成数据集用于训练。动画表示与人脸3D表示解耦，从数据中隐式学习运动，消除对预定义参数化模型的依赖。推理时仅需一次编码即可从任意驱动图像进行动画。其设计目的不是单独增加一个模块，而是用上述表示、监督或渲染约束消除原流程在目标任务中的具体失配。提出将2D扩散模型知识蒸馏到前馈编码器的方法，实现从单张图像即时转换为3D一致、快速且富有表现力的可动画表示。设计解耦动画表示和高效轻量级局部融合策略，动画速度达107.31FPS，无需预定义参数化模型约束。适用时还应保留论文给出的边界：动画表现力受限于蒸馏源模型；极端姿态下可能出现伪影；对输入图像质量敏感。
% METHOD: 2604.04787
\paragraph{AvatarPointillist（自回归 4D 高斯头像生成）（2026.04）}
AvatarPointillist~\cite{2604_04787} 把单图头像生成改成自回归点云生成。DINOv2 图像特征和 FLAME 参数进入 decoder-only Transformer，模型按 token 顺序生成高斯点坐标和绑定信息；随后 Gaussian Decoder 结合 detokenize 坐标、位置编码和 Transformer hidden states 预测完整高斯属性。绑定信息让生成头像可随 FLAME 网格运动变形。相对前馈模板残差方法，它的改变是自适应密度：模型不是在固定模板点上回归属性，而是按序生成需要的点云。
\textbf{实现细节：} AvatarPointillist（2026.04） 保留 GGHead 基于 FLAME 网格驱动的动画框架，但将点云获取方式从 UV 空间的 2D CNN 单次前向生成替换为仅解码器 Transformer 的自回归逐点生成。核心差异：GGHead 的高斯数量由 UV 图分辨率固定（如 $256^{}2$ 即 65K），AvatarPointillist 则将规范空间点云表示为长度 $4N$ 的 token 序列——$N$ 个点各含 4 个量化 token（$x, y, z$ 坐标经 1024 级量化 + FLAME 三角面片绑定索引），由 Transformer 依条件概率 $p(T_n | T_{<n})$ 逐 token 预测，点数量 $N$ 不预设而是由自回归过程的结束标记决定，可根据主体复杂度自适应调整。训练数据通过 GaussianAvatars 为每个身份拟合高斯-绑定表示，坐标按 $y$-$z$-$x$ 字典序排列后量化为离散序列。
\textbf{方法信息：} Alibaba 参与了这篇论文。
\textbf{设计动机与边界：} 相对原始 3DGS 以独立高斯和图像重建误差为主的设定，AvatarPointillist 把问题落在以下可实现的改动上：两阶段流程：(1)Autoregressive Model：decoder-only Transformer包含交叉注意力层、自注意力层和前馈网络。输入图像经DINOv2提取特征，3D人脸重建模型获取FLAME参数，采样顶点经点云编码器获得特征，两者拼接通过交叉注意力注入。输出通过条件概率p(T\_n|T\_<n)顺序生成每个token，完整点云表示为4N个token序列(N个点，每个4个量化token: x,y,z,绑定)。训练使用标准交叉熵损失。(2)Gaussian Decoder：detokenize恢复坐标后经位置编码和MLP产生逐点几何特征，结合AR Transformer的隐藏状态(对生成质量至关重要)通过Transformer解码器预测完整Gaussian参数(位置偏移、尺度、旋转、透明度、颜色)。绑定信息支持通过FLAME网格运动变形实现动画其设计目的不是单独增加一个模块，而是用上述表示、监督或渲染约束消除原流程在目标任务中的具体失配。提出自回归3DGS点云生成新范式，通过decoder-only Transformer自回归生成自适应密度的Gaussian头像点云，联合预测绑定信息实现可动画化，显著提升单图像头像生成的保真度适用时还应保留论文给出的边界：自回归生成速度受序列长度限制，绑定信息依赖FLAME模板网格的表达能力
% METHOD: 2605.24114
\paragraph{COSY（组合式 3DGS GAN 头部编辑）（2026.05）}
COSY~\cite{2605_24114} 解决 3DGS GAN 潜空间纠缠导致编辑不稳定的问题。它把头发、面部皮肤、眼镜和躯干分配给独立子生成器，每个子生成器产生对应语义区域的高斯。眼镜这类加性组件用布尔条件控制存在与否，皮肤和头发用颜色直方图调制。少量 context tokens 在子生成器之间传递形状和光照信息，最终所有高斯统一渲染。相对 GGHead 的单生成器，COSY 的核心是组合性：编辑哪个语义部件，只动对应子生成器。
\textbf{实现细节：} COSY（2026.05） 从组件可控性角度出发，解决了 GGHead 和 CGS-GAN 等单一统一生成器的根本局限——无条件潜在空间中所有语义属性（发型、肤色、配饰）高度纠缠，编辑某一属性时不可避免地改变其他区域。锚点回顾与核心差异：GGHead 和 CGS-GAN 使用单一生成器将所有语义区域混合在同一个潜在空间中；COSY 则将头部合成显式分解为四个独立的 CGS-GAN 风格子生成器——头发、皮肤、眼镜和躯干，各自拥有独立的潜在编码，编辑时仅修改目标组件的潜在向量而保持其余不变。无监督组件分离：组件间的分离完全通过隐式信号实现，无需分割掩码或几何先验。具体地，永久组件（皮肤、头发、躯干）通过稀疏的 10-bin RGB 颜色直方图（每个通道独立统计 10 个 bin，三通道拼接为 30 维向量）调制子生成器中高斯基元的颜色属性来分离——该信号仅输入对应子生成器及其判别器，使生成器隐式学会仅负责指定区域。可添加组件（如眼镜）通过布尔标志（0 或 1）控制存在性：存在时正常生成并加入合成场景，不存在时将高斯不透明度减去一个大值以隐藏；判别器同步接收该标志，惩罚错误缺失或错误出现的生成结果。全局一致性机制：多个独立子生成器的核心挑战是确保组合后几何与光照的匹配。
\textbf{实现补充：} 实现上还要核对输入、输出、参数、特征、初始化、训练损失、正则、渲染或推理接口，以及优化中的剪枝、分裂、采样、编码或解码步骤；这些抓手决定该方法到底改了 3DGS 管线的哪一环。
\textbf{方法信息：} NVIDIA 参与了这篇论文。



\section{全身人体化身}

全身人体方法的锚点通常是 SMPL 或 SMPL-X。高斯从模板表面初始化，随后通过 LBS、非刚性 MLP、triplane、四面体约束或前馈 Transformer 表达衣服、头发和姿态依赖细节。代表性方法 HUGS 和 3DGS-Avatar 奠定了“单目视频到可动画高斯人体”的基本范式；后续方法再解决单图前馈、遮挡、模糊、全身面部细节和长时间外观漂移。


% METHOD: 2311.10812
\paragraph{SplatArmor（正向蒙皮与 SE(3) 非刚性场）（2023.11）}
SplatArmor~\cite{2311_10812} 从单目 RGB 视频恢复可动画人体。它的核心改变是把 SMPL 的 LBS 扩展到 3D 空间任意位置，给规范空间中的每个高斯定义正向蒙皮，避免逆向蒙皮的多解歧义。高斯经正向蒙皮到观察空间后，再由 SE(3) deformation field 预测姿态相关非刚性变形，处理衣服和肌肉动态。颜色不是直接逐高斯自由优化，而由神经颜色场从 3D 位置预测，提供颜色正则和位置监督。相对 HUGS，SplatArmor 更强调空间中任意点的可控蒙皮。
\textbf{实现细节：} HUGS（2023.11） （Human Gaussian Splats）从仅 50--100 帧的单目视频中自动分离静态场景和可动画人体化身。Baseline 与问题：3DGS-Avatar 需要预拟合的 SMPL 参数作为输入，且人体和场景无法自动分离——限制了其在真实视频中的适用性。改了什么：HUGS 实现了人体-场景的自动分离，将整个场景分解为人体高斯（由 SMPL 初始化并参数化）和场景高斯（由 COLMAP 点云初始化），两个子系统在统一框架中联合优化。实现原理：利用 SMPL 体模型参数化人体高斯：每个高斯在 SMPL 规范空间中定义，通过学习到的 LBS 权重和 SMPL 骨架变换到观测空间。关键设计是联合优化 LBS 权重——预定义的 SMPL 蒙皮权重仅基于关节位置，无法区分同一关节影响范围内的衣物和皮肤，学习到的权重可以捕捉这种差异。场景高斯则遵循标准 3DGS 流程，通过 COLMAP 稀疏点云初始化。两套高斯共享同一光栅化器进行渲染，通过前景/背景分离损失和 SMPL 前景掩码实现自动分割。设计动机：人体-场景分离是实现化身重用的关键——只有将人体从场景中分离出来，才能将化身放到新场景中。利用 SMPL 的参数化能力将高斯绑定到人体骨架上，使得动画驱动仅需改变姿态参数即可。
\textbf{方法信息：} Amazon 参与了这篇论文。
\textbf{设计动机与边界：} 相对原始 3DGS 以独立高斯和图像重建误差为主的设定，SplatArmor 把问题落在以下可实现的改动上：SplatArmor的核心方法包含四个部分。第一，扩展混合蒙皮：将SMPL网格的LBS混合蒙皮扩展到3D空间中任意点。通过在规范空间中为每个3D位置插值最近的SMPL顶点的蒙皮权重，定义任意点的正向蒙皮变换。正向蒙皮避免了逆向蒙皮的多解歧义。第二，SE(3)变形场：在规范空间的高斯通过正向蒙皮变换到观察空间后，SE(3)场提供额外的姿态相关非刚性变形（位置偏移、旋转和缩放调整），捕获衣服飘动和肌肉变形等SMPL无法建模的效果。第三，神经颜色场：不直接优化每个高斯的颜色，而是训练一个神经场从3D位置预测颜色。该场为邻近高斯提供颜色正则化，并为高斯均值提供3D位置监督。第四，初始化：高斯从SMPL表面采样，颜色场从粗糙纹理初始化。其设计目的不是单独增加一个模块，而是用上述表示、监督或渲染约束消除原流程在目标任务中的具体失配。提出SplatArmor方法，用3D高斯「武装」SMPL人体模型实现可动画人体重建。将SMPL的混合蒙皮扩展到3D空间中任意位置定义正向蒙皮，避免逆向蒙皮的歧义。引入SE(3)场建模姿态相关的非刚性变形，捕获高斯的位置和各向异性。设计神经颜色场提供颜色正则化和3D监督。适用时还应保留论文给出的边界：依赖SMPL拟合质量；SE(3)场的表达能力有限；对极端姿态和快速运动可能不够稳定。
% METHOD: 2311.17910
\paragraph{HUGS（人体与场景分离的代表性锚点）（2023.11）}
HUGS~\cite{2311_17910} 是全身人体 3DGS 的早期锚点之一，它解决的是原始 3DGS 只能处理静态摄影测量场景的问题：视频里人会动，背景相对静止，若把二者塞进同一套静态高斯，人体姿态变化会变成拖影和错位。HUGS 的输入是一段带相机运动的短单目视频，先用人体估计器得到每帧 SMPL 姿态和共享体型，再把表示拆成两部分：静态场景高斯从 COLMAP/SfM 点云初始化，人体高斯从 SMPL 顶点初始化并放到规范空间。这样，背景仍沿用普通 3DGS 的优化逻辑，而人体由骨架驱动。

关键增量在人体高斯的变形方式。SMPL 只给出光滑人体模板，不能直接解释衣服、头发和局部外形，所以 HUGS 不把高斯固定在 SMPL 表面，而是让每个高斯学习均值偏移、旋转和尺度，从模板附近长出非模板细节。动画时，它也不直接套用 SMPL 自带 LBS 权重，而是为高斯联合优化 LBS 权重，使相邻高斯在新姿态下随骨架协调移动，减少孔洞和漂浮。训练时利用前景 mask 分开人体与背景监督；渲染时先把规范空间人体高斯通过学习到的蒙皮变到目标姿态，再与静态背景高斯合成。这个流程把 SMPL 的可控骨架、3DGS 的实时显式外观和可学习蒙皮放到同一条链路里，因此后续许多人体方法都会把它作为“骨架驱动高斯人体”的参照。

从参数化看，HUGS 不是只给 SMPL 顶点挂一层固定贴图，而是在规范空间为人体高斯学习位置、旋转、尺度、球谐颜色和蒙皮权重。论文中用 feature triplane 与多个小 MLP 从高斯中心位置预测这些属性，训练过程中还允许高斯增密和裁剪：观测到衣服褶皱、头发或鞋子边界时可以增加细节，落在无效区域的高斯则被删掉。这样做的结果是，SMPL 负责提供人体拓扑和姿态坐标系，高斯负责补足模板没有的外观与几何自由度。

实验上，HUGS 在 NeuMan 与 ZJU-MoCap 等设置中报告约 30 分钟训练、60 FPS 渲染，并相对 NeRF 型人体方法大幅缩短训练时间。Apple 参与了这篇论文；代码发布在 \url{https://github.com/apple/ml-hugs}，GitHub 约 391 stars；Semantic Scholar 约 201 次引用。
\textbf{方法信息：} Apple 参与；代码：\url{https://github.com/apple/ml-hugs}；GitHub 约 391 stars；Semantic Scholar 约 201 次引用。



% METHOD: 2312.09228
\paragraph{3DGS-Avatar（等距正则的可动画人体）（2023.12）}
3DGS-Avatar~\cite{2312_09228} 与 HUGS 同属单目视频到可动画人体的锚点，但它更集中地处理“规范空间高斯怎样在新姿态下保持局部形状”的问题。输入包括单目视频、标定相机、SMPL 参数和前景 mask。方法先在规范 SMPL 表面随机采样约 50k 个点作为高斯初始化，每个高斯维护位置、尺度、旋转、不透明度和颜色特征；随后把人体变形分成两级：先做姿态依赖的非刚性变形，再做骨架驱动的刚性 LBS 变换。

非刚性层用多级哈希网格编码规范空间位置，并与由层级姿态编码器得到的姿态潜码拼接，送入浅层 MLP，输出位置偏移、尺度偏移、旋转偏移和颜色相关特征。刚性层再用 skinning-weight MLP 预测每个高斯的蒙皮权重，将非刚性变形后的高斯通过 forward LBS 变到观察空间。相对 HUGS 强调人体与背景分离，3DGS-Avatar 的关键增量是 as-isometric-as-possible 正则：它不仅约束高斯均值的相对距离，也约束由尺度和旋转组成的协方差变形，避免高度关节化或未见姿态下局部高斯被拉伸、压扁或错位。颜色由可学习特征解码，而不是简单依赖固定模板纹理，使局部非刚性变形和外观变化能一起优化。

实验报告显示，它可在约 30 分钟内训练单人 avatar，并以 50+ FPS 渲染，在 ZJU-MoCap、PeopleSnapshot/MonoCap 类设置中相比 HumanNeRF 等方法显著加速。代码发布在官方仓库 \url{https://github.com/mikeqzy/3dgs-avatar-release}，GitHub 约 436 stars；Semantic Scholar 约 257 次引用。
\textbf{锚点细化：} DGS-Avatar：针对 HumanNeRF 等隐式方法体渲染导致每帧数秒推理、无法满足实时交互需求的局限，用\textbf{规范空间 3D 高斯集合 + 可学习非刚性变形网络}替代 NeRF 体渲染，核心思想是\textbf{将复杂人体变形分解为非刚性（衣物褶皱等姿态依赖变形）和刚性（骨骼驱动整体变换）两个层级}——多级哈希网格编码规范位置与姿态潜编码拼接后送入浅层 MLP 输出位置、缩放、旋转偏移，蒙皮 MLP 预测权重、通过 LBS 变换到观测空间，30 分钟训练并以超过 50 FPS 渲染。
\textbf{方法信息：} 代码：\url{https://github.com/mikeqzy/3dgs-avatar-release}；GitHub 约 436 stars；Semantic Scholar 约 257 次引用。



% METHOD: 2405.12806
\paragraph{MOSS（运动感知穿衣人体合成）（2024.05）}
MOSS~\cite{2405_12806} 针对单目视频中衣物随运动变形的问题。它用 SMPL kinematic tree 的关节旋转，通过 matrix-Fisher 分布在 $SO(3)$ 上建模运动概率，得到密度因子和旋转因子来控制高斯分布。运动变形大的关节和衣物区域增加高斯密度，椭圆方向也随全局运动调整。Surface Deformation Detector 再识别单目遮挡造成的信息缺失，在变形显著区域补充几何。项目实现围绕 motion-based clothed human synthesis。相对 3DGS-Avatar 的等距正则，MOSS 更直接用运动学概率调度高斯。
\textbf{实现细节：} MOSS框架包含两个核心模块：(1) Kinematic Gaussian Locating Splatting(KGAS)：将SMPL运动学树中的关节旋转通过matrix-Fisher分布建模为SO(3)流形上的概率分布，提取密度因子和旋转因子来控制3D高斯点的空间位置和方向。密度因子决定在运动变形大的区域（如关节附近）增加高斯点密度，旋转因子根据全局运动调整高斯椭圆的朝向。这确保了全局运动信息传播到整个身体表面，改善运动中人体关节处和紧身衣物区域的重建质量；(2) Surface Deformation Detector(UID)：针对单目视角中局部遮挡导致的表面变形丢失问题，基于KGAS的输出识别变形显著的表面区域，在这些区域进行额外的几何重建以补偿变形信息。UID关注最大几何变形方向上的高斯点，补充被遮挡区域的表面细节。整体框架从SMPL参数出发，以初始点云为输入，通过运动感知的高斯控制实现单目视频的3D穿衣人体合成。
\textbf{方法信息：} 代码：\url{https://github.com/fly-pluche/moss}。
\textbf{设计动机与边界：} 相对原始 3DGS 以独立高斯和图像重建误差为主的设定，MOSS 把问题落在以下可实现的改动上：MOSS框架包含两个核心模块：(1) Kinematic Gaussian Locating Splatting(KGAS)：将SMPL运动学树中的关节旋转通过matrix-Fisher分布建模为SO(3)流形上的概率分布，提取密度因子和旋转因子来控制3D高斯点的空间位置和方向。密度因子决定在运动变形大的区域（如关节附近）增加高斯点密度，旋转因子根据全局运动调整高斯椭圆的朝向。这确保了全局运动信息传播到整个身体表面，改善运动中人体关节处和紧身衣物区域的重建质量；(2) Surface Deformation Detector(UID)：针对单目视角中局部遮挡导致的表面变形丢失问题，基于KGAS的输出识别变形显著的表面区域，在这些区域进行额外的几何重建以补偿变形信息。UID关注最大几何变形方向上的高斯点，补充被遮挡区域的表面细节。整体框架从SMPL参数出发，以初始点云为输入，通过运动感知的高斯控制实现单目视频的3D穿衣人体合成。其设计目的不是单独增加一个模块，而是用上述表示、监督或渲染约束消除原流程在目标任务中的具体失配。首次将运动学树的全局运动信息与matrix-Fisher概率分布结合控制高斯点的密度和旋转，提出UID模块处理单目视角的局部遮挡，在Human NeRF和Gaussian Splatting的LPIPS上分别提升33.94%和16.75%适用时还应保留论文给出的边界：依赖SMPL参数和标定的相机参数作为输入；matrix-Fisher分布的参数回归增加了计算开销；对极端运动（如翻转）的处理能力有限
% METHOD: 2407.21686
\paragraph{ExAvatar（面部、身体和手的完整 SMPL-X 驱动）（2024.07）}
Expressive Whole-Body 3D Gaussian Avatar~\cite{2407_21686} 也称 ExAvatar。它把 3DGS 与 SMPL-X 参数化 mesh 结合，每个高斯被视作 mesh 表面顶点，利用 SMPL-X 拓扑连接性做正则。训练从短单目视频估计 SMPL-X 参数，在 mesh 表面初始化高斯，并联合优化高斯属性和 mesh 参数。连接性正则约束相邻高斯距离和一致性，减少未见表情、身体姿态和手部动作下的伪影。相对 HUGS 的 SMPL 人体，ExAvatar 的增量是完整 SMPL-X 驱动，包括面部表情和手部。
\textbf{实现细节：} DGS-Avatar（2023.12） 首次展示了从单目视频使用 3DGS 创建可动画化身穿身的可行性。Baseline 与问题：此前的 HumanNeRF 等隐式方法虽然能从单目视频重建人体，但其体渲染管线导致训练和推理速度极慢（每帧数秒），完全无法满足实时交互需求。改了什么：3DGS-Avatar 将人体建模为规范空间中的 3D 高斯集合，通过可学习的非刚性变形网络将高斯从规范空间映射到观测空间，利用 3DGS 的可微分光栅化器实现实时渲染。实现原理：核心架构包含三个模块。(1) 非刚性变形模块：分解为姿态依赖的非刚性变形 $F_{nr}$ 和刚性变换 $F_r$。$F_{nr}$ 使用多级哈希网格编码规范位置 $x_c$，与姿态潜编码 $Z_p$（由轻量级层次姿态编码器从 SMPL 姿态参数提取）拼接后送入浅层 MLP（2 个隐藏层，宽度 128），输出高斯位置、缩放、旋转的偏移量以及特征向量 $z$。变形规则为：$x_d = x_c + \Delta x$、$s_d = s_c \cdot \Delta s$、$q_d = q_c \cdot \Delta q$（四元数乘法等价于旋转矩阵相乘）。(2) 刚性变换模块：学习一个蒙皮 MLP 从变形位置 $x_d$ 预测蒙皮权重，通过前向 LBS 变换高斯位置和旋转到观测空间。(3) 颜色 MLP：将变形网络输出的特征 $z$ 通过球谐基函数解码为视角依赖的颜色。初始化时在 SMPL 规范网格表面随机采样 $N=50k$ 个点作为高斯初始位置。设计动机：将复杂的人体变形分解为非刚性（衣物褶皱等姿态依赖变形）和刚性（骨骼驱动的整体变换）两个层级，借鉴了 HumanNeRF 的设计思想但用 3DGS 替代 NeRF 实现了数量级的加速。公式与参数：引入尽可能等距正则化（as-isometric-as-possible），约束规范空间中相邻高斯在变形后保持距离，提升对未见姿态的泛化。
\textbf{方法信息：} 引用很多，Semantic Scholar 约 108 次。
\textbf{设计动机与边界：} 相对原始 3DGS 以独立高斯和图像重建误差为主的设定，Expressive Whole-Body 3D Gaussian Avatar 把问题落在以下可实现的改动上：ExAvatar的核心设计：(1) 混合表示：将3DGS和SMPL-X参数化mesh结合。每个3D高斯被视作mesh表面上的顶点，具有SMPL-X拓扑定义的预连接信息（三角面）。这种连接性使高斯能够随SMPL-X的形状、姿态和表情参数自然变形。(2) 训练过程：从短视频估计SMPL-X参数，在mesh表面初始化3D高斯，然后联合优化高斯属性和mesh参数。利用视频的多视角信息训练，无需3D扫描或RGBD数据。(3) 连接性正则化：利用mesh面片的连接关系约束相邻高斯之间的距离和一致性，显著减少训练视频中未见过的新表情和新姿态下的伪影。(4) 驱动方式：训练后可用SMPL-X的表情编码和3D姿态参数驱动avatar，支持面部表情、身体和手部动作的完整驱动。其设计目的不是单独增加一个模块，而是用上述表示、监督或渲染约束消除原流程在目标任务中的具体失配。提出ExAvatar，首个从单目短视频学习支持面部表情+身体+手部完整驱动的3D人体avatar。SMPL-X+3DGS混合表示利用mesh拓扑连接性作为正则化，显著减少新表情和姿态下的伪影。训练数据仅需单目视频，无需3D扫描。适用时还应保留论文给出的边界：训练视频中表情和姿态的多样性限制了avatar的表现范围；对宽松衣物的建模仍有挑战。
% METHOD: 2412.02684
\paragraph{AniGS（不一致多视图的 4DGS 单图人体）（2024.12）}
AniGS~\cite{2412_02684} 从单张图像生成可动画人体。第一阶段改造视频扩散 Transformer，给定参考图像和 SMPL-X 姿态条件，生成多视角规范姿态 RGB 和法线图。由于这些生成视图可能不一致，第二阶段不把它们当作静态多视图，而是把视角维度建成 4D 问题：用变形 SMPL-X 粗网格初始化 4DGS，先优化 3DGS 参数，再优化时序维度来吸收视角不一致。相对 AdaHuman 的联合扩散与高斯细化，AniGS 的关键是把不一致多视图转成可优化的 4DGS。
\textbf{实现细节：} AniGS包含两个阶段。(1)多视角规范图像生成：改造CogVideo扩散Transformer模型，输入参考图像和SMPL-X姿态条件，生成多视角规范姿态RGB图像和对应法线图。在大型野外视频数据上预训练以增强泛化性，再用合成数据微调并采用数据混合策略（10
\textbf{设计动机与边界：} 相对原始 3DGS 以独立高斯和图像重建误差为主的设定，AniGS 把问题落在以下可实现的改动上：AniGS包含两个阶段。(1)多视角规范图像生成：改造CogVideo扩散Transformer模型，输入参考图像和SMPL-X姿态条件，生成多视角规范姿态RGB图像和对应法线图。在大型野外视频数据上预训练以增强泛化性，再用合成数据微调并采用数据混合策略（10%野外+90%合成）保持泛化。同时生成法线图作为额外几何监督。(2)不一致图像的4DGS重建：将视角不一致的多视图重建建模为4D任务（时间维度对应视角）。用变形SMPL-X粗网格初始化4DGS，先优化3000次3DGS参数，再优化4000次时序维度处理不一致性。最终规范空间的3DGS模型即为高保真3D头像，通过线性混合蒙皮实现动画。为了消除原始 3DGS 在该目标任务中的具体失配，方法把上述表示、监督或渲染约束作为可学习和可执行的改动，而不是单独堆叠模块。提出AniGS，从单张图像生成可动画化3D高斯头像的框架。改造Transformer视频生成模型生成多视角规范姿态图像和法线图，利用大规模视频数据预训练增强泛化能力。将不一致多视图重建问题转化为4D任务，引入高效4D高斯泼溅处理视角不一致性，实现实时推理渲染。适用时还应保留论文给出的边界：生成的多视角图像仍存在不一致性；依赖SMPL-X参数化模型；对极端姿态的动画效果有限。
\textbf{方法信息：} Alibaba 参与了这篇论文。



% METHOD: 2412.07739
\paragraph{GASP（合成先验辅助的 Gaussian Avatar）（2024.12）}
GASP~\cite{2412_07739} 继承 Mesh Attached Gaussians，但把合成数据的完美标注用作 avatar prior。先验模型采用 auto-decoder，学习身份编码和 MLP 解码器；为避免每个高斯参数量爆炸，每个高斯带少量可学习特征，再由小网络映射为属性。推理时先验只用于拟合单图或短视频，最终 avatar 直接渲染高斯，支持实时 360 度渲染。相对 HUGS 从真实单目视频直接优化，GASP 的改变是用合成先验弥补少视角和单图输入。
\textbf{实现细节：} 方法基于Mesh Attached Gaussians框架，将每个高斯基元绑定到FLAME网格的三角形上。先验模型采用auto-decoder架构，联合学习每身份编码z（512维）和MLP解码器D(z)。为解决参数量爆炸问题，为每个高斯引入8维可学习特征向量 $f_i$，通过小网络将特征映射为高斯属性。先验在合成数据集（完美标注、多视角）上训练，学习身份分布。推理时将先验拟合到单张照片或视频，通过优化身份编码z和微调解码器，得到目标人物的高质量头像。先验模型仅在拟合阶段使用，推理时直接渲染高斯基元，实现70fps实时渲染。 提出GASP方法，利用合成数据的像素级完美特性训练Gaussian Avatar先验模型。先验仅需用于拟合阶段而非推理阶段，支持实时渲染。从单张图像或短视频即可生成高质量可动画头像，支持360度渲染，在商用硬件上达到70fps。 依赖合成数据训练先验，合成数据与真实数据之间存在域差距；FLAME模型的表达能力限制了对极端表情和细节的捕捉；对非正面输入数据的质量可能下降。
\textbf{方法信息：} Microsoft 参与了这篇论文；项目页：\url{https://microsoft.github.io/GASP/}；暂未确认完整公开代码。
\textbf{设计动机与边界：} 相对原始 3DGS 以独立高斯和图像重建误差为主的设定，GASP 把问题落在以下可实现的改动上：方法基于Mesh Attached Gaussians框架，将每个高斯基元绑定到FLAME网格的三角形上。先验模型采用auto-decoder架构，联合学习每身份编码z（512维）和MLP解码器D(z)。为解决参数量爆炸问题，为每个高斯引入8维可学习特征向量f\_i，通过小网络将特征映射为高斯属性。先验在合成数据集（完美标注、多视角）上训练，学习身份分布。推理时将先验拟合到单张照片或视频，通过优化身份编码z和微调解码器，得到目标人物的高质量头像。先验模型仅在拟合阶段使用，推理时直接渲染高斯基元，实现70fps实时渲染。其设计目的不是单独增加一个模块，而是用上述表示、监督或渲染约束消除原流程在目标任务中的具体失配。提出GASP方法，利用合成数据的像素级完美特性训练Gaussian Avatar先验模型。先验仅需用于拟合阶段而非推理阶段，支持实时渲染。从单张图像或短视频即可生成高质量可动画头像，支持360度渲染，在商用硬件上达到70fps。适用时还应保留论文给出的边界：依赖合成数据训练先验，合成数据与真实数据之间存在域差距；FLAME模型的表达能力限制了对极端表情和细节的捕捉；对非正面输入数据的质量可能下降。
% METHOD: 2501.13335
\paragraph{Deblur-Avatar（运动模糊人体曝光建模）（2025.01）}
Deblur-Avatar~\cite{2501_13335} 面向运动模糊单目视频。普通 avatar 训练把每帧当作瞬时姿态，但模糊图像来自曝光时间内人体运动积分。方法把曝光时间离散为多个时间戳，插值得到每个时刻的 SMPL 参数，把规范高斯通过刚性和非刚性变形渲染为多张清晰图像，再平均合成模糊图像与输入监督。姿态依赖融合机制预测模糊区域，把模糊区域交给曝光模型渲染，清晰区域直接渲染后再混合。相对 SLAM 章节的 Deblur-SLAM，这里运动主体是人体而不是相机。
\textbf{实现细节：} Deblur-Avatar（2025.01） 处理运动模糊这一常见但被忽视的退化因素。Baseline 与问题：上述方法（包括 AHOY、JOintGS）均假设输入视频帧是清晰的。然而在真实拍摄中，快速运动的人体（如舞蹈、体育场景）在曝光时间内会产生运动模糊——模糊方向取决于运动轨迹，模糊程度取决于运动速度和曝光时间。直接用模糊图像作为光度监督训练 3DGS 会导致化身吸收模糊模式，渲染结果永久模糊。改了什么：Deblur-Avatar 将基于人体运动的运动模糊模型集成到 3DGS 优化循环中。核心思想是"正向生成模糊"：在曝光窗口内采样多个时间步的骨架姿态，对每个时间步进行 LBS 变换和光栅化，然后将多帧渲染结果平均合成模糊图像。生成的模糊图像与真实模糊输入比对，光度损失反向传播同时优化运动轨迹和 3D 高斯参数。设计动机：与外部去模糊预处理不同，将模糊模型嵌入优化循环使化身直接学习清晰外观。因为模糊是可微渲染管线的一部分，梯度信号会引导高斯参数收敛到清晰状态。压缩流式与扩展应用。
\textbf{方法信息：} Huawei 参与了这篇论文。
\textbf{设计动机与边界：} 相对原始 3DGS 以独立高斯和图像重建误差为主的设定，Deblur-Avatar 把问题落在以下可实现的改动上：Deblur-Avatar将人体运动模糊模型集成到3DGS中：(1) 人体运动轨迹建模：在曝光时间T内将人体运动轨迹离散化为n个时间戳，利用插值函数计算每个时刻的SMPL参数，通过刚性和非刚性变形从规范空间变换到各个时刻的姿态。(2) 模糊图像形成：将曝光时间内n个时刻渲染的清晰图像取平均，近似模糊图像的形成过程，用于与实际模糊输入计算损失。(3) 姿态依赖融合机制：预测每个视角的模糊区域，分别优化模糊区域（使用模糊模型渲染）和清晰区域（直接渲染），最终混合得到最优结果。其设计目的不是单独增加一个模块，而是用上述表示、监督或渲染约束消除原流程在目标任务中的具体失配。提出首个针对运动模糊单目视频的可动画3D人体头像重建框架。将人体运动轨迹建模融入3DGS训练，通过显式建模曝光时间内的人体运动轨迹，联合优化轨迹和3D高斯以重建清晰的高质量人体头像。引入姿态依赖融合机制区分运动区域和静止区域。适用时还应保留论文给出的边界：依赖SMPL人体模型的准确性。对于极端复杂的非刚性变形（如宽松衣物）效果可能受限。运动轨迹离散化精度和计算效率之间存在权衡。
% METHOD: 2503.10625
\paragraph{LHM（单图秒级大人体重建模型）（2025.03）}
LHM~\cite{2503_10625} 把人体 3DGS 从 per-subject 优化推向大规模前馈模型。它用多模态 Body-Head Transformer：几何 tokens 从 SMPL-X 表面采样，body 图像 tokens 来自 Sapiens-1B，头部特征金字塔负责高频面部细节。几何 token 与图像 token 注意力融合后，解码为高斯的变形、尺度、旋转和 SH 系数。训练时通过 LBS 把规范空间 avatar 变到目标视角。相对 3DGS-Avatar 的单人训练，LHM 的增量是秒级单图生成和更强泛化。
\textbf{锚点细化：} LHM：针对从单图像创建全身化身的多阶段优化范式耗时数小时且难以泛化的局限，用\textbf{大规模前馈 Transformer 直接从单图回归全身 3DGS}替代多阶段优化，核心思想是\textbf{一次前向传播生成全身化身}——基于多视角大规模数据预训练强先验，无需测试时优化即可在秒级时间内输出可驱动全身化身。
	extbf{实现细节：} LHM 是前馈式人体高斯生成器。它从单张人体图像和 SMPL-X 结构先验构造三类 token：由身体表面采样得到的几何 token、预训练 Sapiens 特征形成的身体图像 token，以及多尺度头部特征金字塔 token。MBHT 用注意力融合三者，并逐高斯解码规范空间的位移、尺度、旋转和球谐颜色；LBS 再依据目标姿态把高斯变形到渲染空间。训练以多视角重建监督约束渲染外观和几何/姿态一致性，推理则一次前向预测高斯后再通过 LBS 驱动动画，避免逐人物从头优化。
\textbf{实现细节：} LHM（2025.03） （Large Animatable Human Reconstruction Model）修改了 3DGS 的化身创建管线——用多模态 Transformer 前向推理替代多视角优化循环——3D 高斯的参数化表示和可微分光栅化渲染管线保持不变。Baseline 与问题。上文介绍的 3DGS-Avatar 和 GoMAvatar 等先驱方法都依赖单目视频序列进行逐帧优化（3DGS-Avatar 需 30 分钟训练），且需要预拟合的 SMPL 参数。核心瓶颈在于：标准 3DGS 的化身创建本质是一个优化问题（从 COLMAP 点云或 SMPL 网格初始化高斯，然后迭代优化参数），而非推理问题——给定一张图像，系统无法"一步"生成化身。改了什么。LHM 将化身创建从优化问题重新定义为回归问题：设计多模态 Transformer 架构，以单张 RGB 图像为输入，在单次前向传播中直接预测规范空间中每个 3D 高斯的全部参数（位置偏移、旋转、缩放、球谐系数）。旧机制——多视角输入 + COLMAP 初始化 + 迭代优化——被新机制——单图像输入 + SMPL-X 采样点初始化 + Transformer 回归——替代。实现原理。LHM 的核心是多模态体-头 Transformer（Multimodal Body-Head Transformer, MBHT），它融合三类 token 来回归高斯参数。(1) \textit{几何 token}：从 SMPL-X 规范网格表面采样 3D 点 $\{X_i\}_{i=1}^{N_{\mathrm{points}}}$，经正弦位置编码和 MLP 投影得到几何 query token： \[ T\_{3D} = \phi\_{\text{proj}}(\gamma(X)) \in \mathbb{R}\^{}{N\_{\text{points}} \times C} \] 直觉上，SMPL-X 网格上的采样点充当"问题点"——每个点询问图像"我在这个位置应该长什么样"。$C$ 是 token 维度。(2) \textit{身体图像 token}：使用在千万级人体图像上预训练的 Sapiens-1B 编码器提取语义身体特征。由于 Sapiens 是专门针对人体预训练的视觉 Transformer，它能捕获比通用编码器更丰富的衣物纹理和身体形状信息。(3) \textit{头部特征金字塔 token}：人的头部在输入图像中只占很小的区域，经过编码器的空间下采样后关键面部细节（皱纹、瞳孔）几乎完全丢失。三类 token 在 MBHT 块中通过注意力机制交互融合。
\textbf{方法信息：} Alibaba 参与了这篇论文；引用很多，Semantic Scholar 约 55 次。



% METHOD: 2504.20403
\paragraph{TetGS（四面体约束的可编辑人体）（2025.04）}
Creating Your Editable 3D Photorealistic Avatar with Tetrahedron-constrained Gaussian Splatting~\cite{2504_20403} 关注可编辑性。普通自由高斯很难定位和编辑局部区域，TetGS 把高斯显式嵌入四面体网格中，由四面体 SDF 值和局部重心坐标控制高斯位置。训练分为实例化、局部空间适配和基于几何的外观生成三阶段。局部空间适配通过更新 SDF 值驱动高斯重新定位，从而支持较大几何变化。相对 SMPL 表面绑定，四面体给高斯一个可编辑的体空间控制结构。
\textbf{实现细节：} TetGS将高斯核显式嵌入四面体网格中，每个高斯的位置由四面体顶点的SDF值确定。方法包含三个渐进阶段：(1) 3D化身实例化：从单目视频重建人体先验，为TetGS提供准确的初始化。(2) 局部空间适配：利用显式划分的四面体引导高斯核的重新分配，通过更新SDF值驱动高斯重新定位，实现大范围几何变化。(3) 基于几何的外观生成：采用从粗到细的激活策略逐步生成逼真外观。高斯的位置通过父三角形顶点的重心权重计算，加上沿法线的可学习位移。 提出TetGS混合表示，将可控的四面体网格显式结构与3DGS高精度渲染结合，通过渐进式三阶段优化实现从单目视频创建可编辑3D逼真化身，支持精确区域定位和几何自适应。 依赖单目视频质量；对极端姿态变化的处理能力有限；四面体网格的分辨率影响最终质量。
\textbf{方法信息：} Alibaba 参与了这篇论文。
\textbf{设计动机与边界：} 相对原始 3DGS 以独立高斯和图像重建误差为主的设定，Creating Your Editable 3D Photorealistic Avatar with Tetrahedron-constrained Gaussian Spla 把问题落在以下可实现的改动上：TetGS将高斯核显式嵌入四面体网格中，每个高斯的位置由四面体顶点的SDF值确定。方法包含三个渐进阶段：(1) 3D化身实例化：从单目视频重建人体先验，为TetGS提供准确的初始化。(2) 局部空间适配：利用显式划分的四面体引导高斯核的重新分配，通过更新SDF值驱动高斯重新定位，实现大范围几何变化。(3) 基于几何的外观生成：采用从粗到细的激活策略逐步生成逼真外观。高斯的位置通过父三角形顶点的重心权重计算，加上沿法线的可学习位移。其设计目的不是单独增加一个模块，而是用上述表示、监督或渲染约束消除原流程在目标任务中的具体失配。提出TetGS混合表示，将可控的四面体网格显式结构与3DGS高精度渲染结合，通过渐进式三阶段优化实现从单目视频创建可编辑3D逼真化身，支持精确区域定位和几何自适应。适用时还应保留论文给出的边界：依赖单目视频质量；对极端姿态变化的处理能力有限；四面体网格的分辨率影响最终质量。
% METHOD: 2505.18342
\paragraph{Pose Splatter（动物姿态与外观的高斯前馈模型）（2025.05）}
Pose Splatter~\cite{2505_18342} 扩展到动物姿态与外观量化。它不依赖 SMPL 这类人体模板，而是用少量校准相机的 mask 图像做 shape carving，得到体素粗表示，再用堆叠 3D U-Net 预测 occupancy 和特征通道。有效体素渲染为高斯，其余特征经 MLP 预测位移、协方差和外观。项目实现面向动物建模。相对人体高斯 avatar，它的意义在于去掉特定人体先验，只用多视图 silhouette 和旋转不变视觉嵌入学习姿态外观。
\textbf{实现细节：} Pose Splatter框架利用校准相机的少量遮罩图像，通过形状雕刻技术创建动物的体素化粗表示。粗体积经过堆叠3D U-Net架构处理，生成包含占用通道和附加特征通道的清洁体积。占用通道决定是否为每个体素渲染高斯，其余特征通道通过小型MLP映射为3D位移向量（反体素化）以及协方差和外观特征。最终通过高斯泼溅进行场景渲染，使用基于图像的损失传播梯度。整体框架仅使用约2.5GB GPU显存。核心组件包括：1）使用SAM2生成遮罩视频；2）通过鲁棒三角测量确定动物3D中心位置和方位角；3）形状雕刻过程利用多视图遮罩反投影确定体素有效性；4）旋转不变视觉嵌入技术用于编码姿态和外观。 提出首个基于3D高斯泼溅的前馈式动物3D姿态外观建模框架，无需先验动物几何知识、逐帧优化或手动标注即可学习准确的3D动物几何，并提出旋转不变视觉嵌入技术替代3D关键点数据。 依赖多相机设置（通常3-6个校准相机）；形状雕刻的分辨率受体素网格大小限制；SAM2遮罩质量直接影响重建效果；对于毛发丰富的动物可能难以精确重建。
\textbf{方法信息：} 代码：\url{https://github.com/jackgoffinet/pose-splatter}。
\textbf{设计动机与边界：} 相对原始 3DGS 以独立高斯和图像重建误差为主的设定，Pose Splatter 把问题落在以下可实现的改动上：Pose Splatter框架利用校准相机的少量遮罩图像，通过形状雕刻技术创建动物的体素化粗表示。粗体积经过堆叠3D U-Net架构处理，生成包含占用通道和附加特征通道的清洁体积。占用通道决定是否为每个体素渲染高斯，其余特征通道通过小型MLP映射为3D位移向量（反体素化）以及协方差和外观特征。最终通过高斯泼溅进行场景渲染，使用基于图像的损失传播梯度。整体框架仅使用约2.5GB GPU显存。核心组件包括：1）使用SAM2生成遮罩视频；2）通过鲁棒三角测量确定动物3D中心位置和方位角；3）形状雕刻过程利用多视图遮罩反投影确定体素有效性；4）旋转不变视觉嵌入技术用于编码姿态和外观。其设计目的不是单独增加一个模块，而是用上述表示、监督或渲染约束消除原流程在目标任务中的具体失配。提出首个基于3D高斯泼溅的前馈式动物3D姿态外观建模框架，无需先验动物几何知识、逐帧优化或手动标注即可学习准确的3D动物几何，并提出旋转不变视觉嵌入技术替代3D关键点数据。适用时还应保留论文给出的边界：依赖多相机设置（通常3-6个校准相机）；形状雕刻的分辨率受体素网格大小限制；SAM2遮罩质量直接影响重建效果；对于毛发丰富的动物可能难以精确重建。
% METHOD: 2505.24877
\paragraph{AdaHuman（姿态条件扩散与组合式细化）（2025.05）}
AdaHuman~\cite{2505_24877} 从单张野外图像生成高保真可动画人体。它先用姿态条件 3D 联合扩散模型同时生成多视角一致图像和 3DGS 重建；输入条件包括局部身体视角、SMPL 渲染的 2D pose map 和相机 ray map。随后组合式 3DGS 细化模块把初始 avatar 渲染成局部放大视角，用多视角扩散做 image-to-image refinement，再通过裁剪感知 ray map 合回全局高斯。相对 LHM 的纯前馈 Transformer，AdaHuman 更依赖扩散生成来补足未见视角和局部细节。
\textbf{实现细节：} AdaHuman（2025.05） 进一步提升了扩散辅助的化身创建质量。Baseline 与问题：PERSONA 的扩散模型生成整体人体图像，但不同身体部位（头部、躯干、手部）在渲染质量上的重要性差异很大——面部和手部是视觉注意力集中的区域，质量要求远高于躯干。改了什么：AdaHuman 引入姿态条件化 3D 联合扩散模型，在任意姿态下合成一致的多视角图像。同时通过组合式 3DGS 精修模块分别处理头部、躯干、手部等局部身体部位的细节增强——为不同部位分配独立的精修网络，使计算资源集中在视觉上最敏感的区域。在单目视频设定下，几何信息的匮乏使得化身重建更加困难。以下方法通过引入基础模型先验和修复策略来弥补这一缺陷。
\textbf{实现补充：} 实现上还要核对输入、输出、参数、特征、初始化、训练损失、正则、渲染或推理接口，以及优化中的剪枝、分裂、采样、编码或解码步骤；这些抓手决定该方法到底改了 3DGS 管线的哪一环。
\textbf{方法信息：} NVIDIA 参与了这篇论文。
\textbf{设计动机与边界：} 相对原始 3DGS 以独立高斯和图像重建误差为主的设定，AdaHuman 把问题落在以下可实现的改动上：AdaHuman框架由两个核心模块组成。第一个是姿态条件3D联合扩散模型：给定全身输入图像，首先生成不同身体部位（头部、上半身、下半身）的局部视角图像，结合输入图像形成输入视角集合。这些视角以RGB图像、2D语义姿态图（从SMPL模型渲染得到）和相机参数（编码为正弦嵌入的相机光线图）的元组表示。扩散模块在每一步同时进行图像去噪和3DGS重建，确保多视角一致性。关键的是，通过简单改变姿态条件即可在任意姿态下生成图像，从而支持生成最小化自遮挡的A姿态头像。第二个是组合式3DGS细化模块：从初始3DGS头像渲染局部身体部位的放大视角，利用多视角扩散模型进行图像到图像细化以提升细节和分辨率。然后通过裁剪感知相机光线图建立局部和全局视角间的精确对应关系，将细化后的局部视角无缝融合为高细节的整体3D头像。该模块使用光栅化器中修改的alpha-blending公式处理重叠的高斯区域。其设计目的不是单独增加一个模块，而是用上述表示、监督或渲染约束消除原流程在目标任务中的具体失配。本文提出AdaHuman框架，包含两个核心创新：(1) 姿态条件3D联合扩散模型，在扩散过程的每一步同时合成多视角一致图像和对应3DGS重建；(2) 组合式3DGS细化模块，通过图像到图像细化增强局部身体部位细节，并使用新颖的裁剪感知相机光线图无缝集成，生成具有高细节的统一3D头像。该方法可生成标准化A姿态头像，支持任意输入动作的绑定和动画。适用时还应保留论文给出的边界：对极端姿态和严重遮挡的处理仍有挑战；多视角扩散模型计算开销较大；非常规服装（如飘逸长裙）的细节保持有限。
% METHOD: 2506.13766
\paragraph{LHM++（姿态无标注图像到高效人体）（2025.06）}
LHM++~\cite{2506_13766} 延续 LHM，但把计算效率和姿态无标注输入作为重点。它使用 Encoder-Decoder Point-Image Transformer：从 SMPL-X 表面采样几何 tokens，编码阶段逐层 grid pooling 下采样点 token 并与图像 token 注意力交互，解码阶段逐步上采样恢复细节。图像 token 合并借鉴 ToME，减少冗余注意力计算。最后用轻量神经动画渲染器改善 LBS 动画结果。相对 LHM，LHM++ 的改变不是换任务，而是把大模型注意力压到更可部署的层次结构。
\textbf{实现细节：} LHM++（2025.06） 进一步改进 LHM 的效率和输入灵活性。Baseline 与问题：LHM 虽然实现了单图像前馈化身重建，但其多模态体-头 Transformer 中的全局注意力机制计算复杂度随点 token 数量二次增长（规范网格上采样点越多，注意力矩阵越大），限制了可处理的点密度和推理速度。此外，LHM 要求输入图像附带准确的相机参数和 SMPL-X 姿态标注，在用户随意拍摄的照片中这些信息通常不可用。改了什么：LHM++ 在架构层面引入编码器-解码器点-图像 Transformer（Encoder-Decoder PIT），将 LHM 的单层全局注意力替换为渐进式编解码流程；在输入层面去除了对相机参数和姿态标注的依赖，支持从一张或多张无标注的随意拍摄图像直接生成可动画化身。实现原理：核心架构 PIT 的工作方式如下。编码器阶段采用 Point Transformer v3 的分块自注意力（partitioned self-attention），将规范空间中的 3D 几何点 token 按空间邻近性分块处理，通过网格池化（grid pooling）逐层压缩点 token 数量——浅层处理密集的局部点特征，深层处理稀疏的全局语义特征。
\textbf{方法信息：} Alibaba 参与了这篇论文。
\textbf{设计动机与边界：} 相对原始 3DGS 以独立高斯和图像重建误差为主的设定，LHM++ 把问题落在以下可实现的改动上：方法的核心架构包含：1）SMPL-X锚点初始化：从标准SMPL-X表面均匀采样空间点作为3D几何Token的初始化锚点；2）编码器-解码器PIT框架：编码阶段通过多个PIT块逐步利用Grid Pooling下采样几何Token，每层下采样后的点Token与图像Token通过多模态注意力交互；解码阶段逐步上采样恢复精细几何细节。这种层次化设计将注意力计算从O((N\_p+N)\^{}2)降至更可控的级别；3）图像Token合并：受ToME启发，在帧内注意力后基于Key相似度按固定比例r合并冗余图像Token，在全局注意力后再解除合并，意外发现适度合并甚至能轻微提升性能；4）轻量级神经动画渲染器：在3DGS渲染结果上应用实时神经渲染器，通过线性混合蒙皮（LBS）实现姿态驱动动画。最终几何点特征通过轻量MLP头回归3D高斯参数。其设计目的不是单独增加一个模块，而是用上述表示、监督或渲染约束消除原流程在目标任务中的具体失配。提出高效的Encoder-Decoder Point-Image Transformer（PIT）框架，在保持重建质量的同时显著降低计算开销；引入基于ToME的图像Token合并策略减少冗余计算；设计轻量级3D感知神经动画渲染器提升动画渲染质量；实现从一张或多张无姿态图像在数秒内生成高保真可动画3D人体。适用时还应保留论文给出的边界：仍依赖SMPL-X参数化模型作为几何先验；极度遮挡区域恢复有限；多图像输入时内存仍随图像数量增长。
% METHOD: 2601.02098
\paragraph{InpaintHuman（遮挡人体的身份保持修复）（2026.01）}
InpaintHuman~\cite{2601_02098} 从遮挡单目视频重建完整可动画人体。它用多尺度 UV 参数化表示人体：粗尺度 UV 有更大感受野，适合插值遮挡区域；细尺度保留可见细节；粗到细融合形成规范表示。遮挡补全不使用 SDS 间接优化，而是用 identity-preserving diffusion inpainting 生成像素级监督，文本反转捕获身份概念，语义条件引导未见身体部位修复。每个采样点关联高斯并通过 SMPL 蒙皮动画。相对 LAGA 的服装生成，InpaintHuman 的重点是遮挡情况下仍保持身份和时序一致。
\textbf{实现细节：} InpaintHuman（2026.01） 专门解决遮挡场景下的化身重建，与 AHOY 从野外视频处理遮挡的路线互补。Baseline 与问题：AHOY 用视频扩散模型为未观察区域"幻觉"出监督，但视频扩散的逐帧随机性可能引入时序不一致。InpaintHuman 从图像修复而非视频生成的角度解决遮挡问题。改了什么：InpaintHuman 通过多尺度 UV 参数化表示和身份保持扩散修复模块实现完整化身的生成。该模块对个性化扩散模型进行身份微调后，为未观察到的身体区域生成像素级监督信号——即使人体的一半被遮挡物覆盖，也能推断出完整的 3D 外观。身份微调确保修复结果与已观察区域保持身份一致性，不会生成属于其他人的外观。物理与衣物动态。本节讨论 2 条技术路线：（1）物理模拟与衣物动态；（2）遮挡、模糊和野外场景的处理。
\textbf{实现补充：} 实现上还要核对输入、输出、参数、特征、初始化、训练损失、正则、渲染或推理接口，以及优化中的剪枝、分裂、采样、编码或解码步骤；这些抓手决定该方法到底改了 3DGS 管线的哪一环。
\textbf{方法信息：} Alibaba 参与了这篇论文。
\textbf{设计动机与边界：} 相对原始 3DGS 以独立高斯和图像重建误差为主的设定，InpaintHuman 把问题落在以下可实现的改动上：方法包含两大核心：(1) 多尺度UV参数化表示：在SMPL模板的规范姿态空间中利用UV参数化表示人体，不同尺度的UV图提供不同的遮挡鲁棒性-细节保持权衡。粗尺度UV图通过更大感受野更好地插值遮挡区域，细尺度保持细节。层次化粗到细特征融合策略综合两者优势；(2) 身份保持扩散修复模块：利用文本反转捕获概念级身份特征，通过语义引导的个性化扩散修复完成未见身体部位。修复结果作为像素级监督精炼规范表示。每个采样点关联3D高斯原语，通过SMPL蒙皮变形实现动画化。其设计目的不是单独增加一个模块，而是用上述表示、监督或渲染约束消除原流程在目标任务中的具体失配。提出首个针对遮挡单目视频的完整3D头像重建方法。两大核心创新：(1) 多尺度UV参数化表示，通过层次化粗到细特征插值对遮挡区域提供鲁棒重建；(2) 身份保持的扩散修复模块，结合文本反转和语义条件引导实现主体特定、时间一致的遮挡补全。不同于SDS方法，采用直接像素级监督确保身份保真度。适用时还应保留论文给出的边界：依赖SMPL参数估计质量；扩散修复可能引入与真实几何不一致的细节；对极端遮挡（全身大部分被遮挡）效果有限。
% METHOD: 2602.04317
\paragraph{JOintGS（相机、身体和高斯联合优化）（2026.02）}
JOintGS~\cite{2602_04317} 面向野外单目视频中的粗初始化。它把前景人体、背景场景、相机外参和人体姿态放到同一优化中。前景用 triplane 定义规范高斯人体并通过 LBS 变形，偏移模块和残差颜色场建模衣物褶皱与光照变化；背景用静态 3DGS 锚定相机估计。背景光度一致性精化相机，精化相机又改善人体对齐，姿态优化减少动态伪影。项目实现包含 Joint Optimization。相对 HUGS 的人体场景分离，JOintGS 更强调二者协同反向修正粗相机和粗 SMPL。
\textbf{实现细节：} JOintGS（2026.02） （Joint Optimization of Cameras, Bodies and 3D Gaussians）正是针对这一前提的脆弱性，提出相机外参、人体姿态和 3D 高斯表示的统一联合优化框架，在规范空间中用三平面表示解码人体高斯属性，通过前景-背景分离实现协同精修——多视角光栅化渲染管线保持不变。Baseline 与问题。现有 3DGS 化身方法（3DGS-Avatar、HUGS、GaussianAvatar、ExAvatar 等）严重依赖精确的相机标定和姿态标注作为固定输入。然而在野外单目视频场景中，COLMAP 等运动恢复结构工具因动态人体运动导致静态特征对应不足，产生噪声相机轨迹；HMR2.0 等单目姿态估计器因深度歧义和遮挡产生不准确的 3D SMPL 参数。即使是微小的初始误差也会级联放大为对齐不准、时序不一致和人体-场景穿透等严重伪影。改了什么。
\textbf{方法信息：} 代码：\url{https://github.com/mililab/jointgs}。
\textbf{设计动机与边界：} 相对原始 3DGS 以独立高斯和图像重建误差为主的设定，JOintGS 把问题落在以下可实现的改动上：框架包含四个核心组件：(1)前景人体表示——在规范空间用Triplane表示定义3D高斯人体，通过LBS蒙皮变形到任意姿态，引入偏移模块和残差颜色场建模时序动态（服装褶皱、外观变化）；(2)背景场景表示——使用静态3D高斯建模背景，利用多视图一致性约束相机估计；(3)协同精化机制——静态背景高斯通过光度一致性约束相机轨迹（不受动态人体影响），精化相机通过准确时序对应改善人体对齐，优化人体姿态通过消除动态伪影提升静态约束下的场景重建；(4)残差颜色场——建模光照变化。初始化使用COLMAP粗相机和HMR2.0粗SMPL参数，通过可微渲染联合精化。其设计目的不是单独增加一个模块，而是用上述表示、监督或渲染约束消除原流程在目标任务中的具体失配。提出JOintGS统一框架，通过前景-背景解耦实现相机外参、人体姿态和3D高斯表示的联合优化。静态背景高斯锚定相机估计，精化相机改善人体对齐，优化姿态提升场景重建，三者协同增强。在NeuMan数据集上PSNR提升2.1dB。适用时还应保留论文给出的边界：依赖粗初始化质量；极端遮挡或运动模糊场景仍有限制；计算开销因联合优化而增加。
% METHOD: 2602.05572
\paragraph{ShapeGaussian（视觉先验驱动的无模板 4D 人体重建）（2026.02）}
ShapeGaussian~\cite{2602_05572} 面向随手单目视频中的高保真 4D human reconstruction。通用 4DGS 方法在人体这种大形变对象上容易缺少多视角约束；HUGS 等模板化方法依赖 SMPL，如果人体姿态估计有误，衣物、肢体和遮挡区域会产生伪影。ShapeGaussian 的核心选择是 template-free：不把 SMPL 当作唯一几何骨架，而是把视觉先验、多参考帧和神经形变模型结合起来，重建动态人体高斯。
方法采用两阶段流程。第一阶段利用预训练视觉模型估计数据驱动先验，学习粗糙但稳定的可形变几何；这些先验帮助系统在单目输入中推断不可见关键点和大幅姿态变化。第二阶段用 neural deformation model 细化几何和动态细节，把多参考帧中的外观线索聚合到随时间变化的人体表示里。训练时，图像重建损失、几何一致性约束和视觉先验共同限制高斯位置、尺度、外观和变形场，避免单帧监督把人体拉成局部错误形状。相对 JOintGS 主要联合优化相机、身体和背景，ShapeGaussian 更强调用无模板视觉先验替代强 SMPL 依赖。
\textbf{方法信息：} 作为 2026 年扩展召回条目纳入；当前仅依据本地 arXiv 元数据与分类材料写作，具体实验数字未在正文复述。
\textbf{设计动机与边界：} 相对原始 3DGS 以独立高斯和图像重建误差为主的设定，ShapeGaussian 把问题落在以下可实现的改动上：方法采用两阶段流程：先用预训练模型估计数据驱动先验，学习粗糙可形变几何；再用神经 deformation model 细化几何并捕获动态细节。它针对通用 4DGS 在单目高形变人体运动中缺少多视角约束的问题，将 2D vision priors 和多参考帧用于处理不可见关键点与姿态估计误差。其设计目的不是单独增加一个模块，而是用上述表示、监督或渲染约束消除原流程在目标任务中的具体失配。用模板无关视觉先验替代对 SMPL 姿态估计的强依赖，以缓解模板方法在姿态误差下的伪影。适用时还应保留论文给出的边界：当前仅基于摘要分类；具体 Gaussian 参数化、损失和数据集表格需要全文核验。
% METHOD: 2604.00928
\paragraph{Autoregressive Appearance Prediction（长期外观潜变量预测）（2026.04）}
Autoregressive Appearance Prediction for 3D Gaussian Avatars~\cite{2604_00928} 处理长时间捕获中的姿态-外观虚假相关。它基于 sparse anchors、spatial MLPs、control points 和 dense Gaussians 的层次表示，引入局部化姿态条件：每个锚点只接收由 skinning weight mask 选出的局部姿态特征，并用身体区域 PCA 减少长程假相关。每帧外观潜变量从 UV 纹理编码得到；测试时没有 UV 纹理，于是用 Transformer 从短历史姿态和前一潜变量自回归预测下一潜变量。相对普通姿态条件 MLP，它让外观随时间平滑演化，而不是被单帧姿态硬编码。
\textbf{实现细节：} 基于Spatial MLPs层次表示构建：(1)稀疏锚点+空间MLP+控制点+密集3D高斯层次结构；(2)局部化姿态条件：基于皮肤权重局部掩码每个锚点接收的姿态特征，加上7个身体区域的局部PCA，减少虚假长程姿态相关性；(3)外观潜变量：每帧学习一个外观潜变量，通过卷积外观编码器从UV纹理（多视角投影到模板网格获得）编码为2D特征图，通过锚点UV坐标双线性采样获取局部潜变量；(4)自回归外观预测器：测试时UV纹理不可用，使用transformer从短历史掩码姿态和前一潜变量回归下一潜变量，实现平滑外观演化。训练损失：L1、感知损失、不透明度正则化、尺度正则化、变形平滑正则化、KL散度外观潜变量正则化。 提出带外观潜变量的3DGS虚拟形象模型，通过局部化姿态条件减少虚假相关性，自回归外观预测器实现时间平滑稳定驱动，解决姿态-外观歧义问题。 依赖模板网格和UV映射；外观编码器训练需要多视角数据；长序列自回归可能累积误差
\textbf{方法信息：} Meta 参与了这篇论文。
\textbf{设计动机与边界：} 相对原始 3DGS 以独立高斯和图像重建误差为主的设定，Autoregressive Appearance Prediction for 3D Gaussian Avatars 把问题落在以下可实现的改动上：基于Spatial MLPs层次表示构建：(1)稀疏锚点+空间MLP+控制点+密集3D高斯层次结构；(2)局部化姿态条件：基于皮肤权重局部掩码每个锚点接收的姿态特征，加上7个身体区域的局部PCA，减少虚假长程姿态相关性；(3)外观潜变量：每帧学习一个外观潜变量，通过卷积外观编码器从UV纹理（多视角投影到模板网格获得）编码为2D特征图，通过锚点UV坐标双线性采样获取局部潜变量；(4)自回归外观预测器：测试时UV纹理不可用，使用transformer从短历史掩码姿态和前一潜变量回归下一潜变量，实现平滑外观演化。训练损失：L1、感知损失、不透明度正则化、尺度正则化、变形平滑正则化、KL散度外观潜变量正则化。其设计目的不是单独增加一个模块，而是用上述表示、监督或渲染约束消除原流程在目标任务中的具体失配。提出带外观潜变量的3DGS虚拟形象模型，通过局部化姿态条件减少虚假相关性，自回归外观预测器实现时间平滑稳定驱动，解决姿态-外观歧义问题。适用时还应保留论文给出的边界：依赖模板网格和UV映射；外观编码器训练需要多视角数据；长序列自回归可能累积误差
% METHOD: 2604.09835
\paragraph{F3G-Avatar（面部优先的全身高斯）（2026.04）}
F3G-Avatar~\cite{2604_09835} 解决全身 avatar 中面部特写质量不足的问题。它从 Momentum Human Rig 模板出发，渲染前后 position maps，并用双分支解码高斯：身体分支预测姿态依赖全身属性，面部聚焦分支用轻量 StyleGAN 网络处理头部 position map，生成更高分辨率面部高斯。训练加入面部特定 adversarial loss。项目实现围绕 face-focused full-body avatar。相对 ExAvatar 的统一全身 SMPL-X 表面绑定，F3G 的改变是给脸单独一条高容量生成分支。
\textbf{实现细节：} 从穿着衣物的MHR模板出发，渲染前后位置图并通过双分支架构解码为3D高斯。身体分支使用BodyUNet预测姿态依赖的高斯属性图；面部聚焦分支使用三个轻量级StyleGAN网络处理头部特定的位置图，预测高分辨率姿态依赖面部高斯图。预测的高斯通过线性混合蒙皮(LBS)进行姿态变形，最后通过可微分高斯泼溅渲染。训练结合重建损失、感知损失和面部特定对抗损失，增强特写视图的真实感。 提出F3G-Avatar，首个面向面部的全身高斯化身合成方法。采用双分支架构：全局身体分支捕获姿态依赖非刚性变形，面部聚焦分支细化头部几何和外观。引入Momentum Human Rig (MHR)模板和面部特定对抗损失，在AvatarReX数据集上面部视角达到PSNR 26.243/SSIM 0.964/LPIPS 0.084的优异性能。 依赖多视角RGB输入，无法从单目视频工作，面部建模仍受限于SMPL-X参数化精度，对极端表情建模有局限。
\textbf{方法信息：} 代码：\url{https://github.com/wjmenu/f3g-avatar}。
\textbf{设计动机与边界：} 相对原始 3DGS 以独立高斯和图像重建误差为主的设定，F3G-Avatar 把问题落在以下可实现的改动上：从穿着衣物的MHR模板出发，渲染前后位置图并通过双分支架构解码为3D高斯。身体分支使用BodyUNet预测姿态依赖的高斯属性图；面部聚焦分支使用三个轻量级StyleGAN网络处理头部特定的位置图，预测高分辨率姿态依赖面部高斯图。预测的高斯通过线性混合蒙皮(LBS)进行姿态变形，最后通过可微分高斯泼溅渲染。训练结合重建损失、感知损失和面部特定对抗损失，增强特写视图的真实感。其设计目的不是单独增加一个模块，而是用上述表示、监督或渲染约束消除原流程在目标任务中的具体失配。提出F3G-Avatar，首个面向面部的全身高斯化身合成方法。采用双分支架构：全局身体分支捕获姿态依赖非刚性变形，面部聚焦分支细化头部几何和外观。引入Momentum Human Rig (MHR)模板和面部特定对抗损失，在AvatarReX数据集上面部视角达到PSNR 26.243/SSIM 0.964/LPIPS 0.084的优异性能。适用时还应保留论文给出的边界：依赖多视角RGB输入，无法从单目视频工作，面部建模仍受限于SMPL-X参数化精度，对极端表情建模有局限。
% METHOD: 2604.10259
\paragraph{HumanGS（稀疏多视角前馈重建与实时动画）（2026.04）}
Real-Time Human Reconstruction and Animation using Feed-Forward Gaussian Splatting~\cite{2604_10259} 也叫 HumanGS。它从稀疏多视角 RGB 和 SMPL-X 姿态直接预测可动画高斯人体。每个 SMPL-X 顶点关联受约束高斯和若干无约束高斯，前者贴近模板表面，后者捕获衣服、头发等偏离表面的结构。Transformer 编码器提取多视角 3D 特征，神经纹理和姿态图建立输入与几何的对应。项目实现支持前馈运行。相对 LHM 单图输入，HumanGS 更依赖多视角稀疏观测，但动画时可直接用 LBS，不需逐帧网络推理。
\textbf{实现细节：} 基于LVSM(Large View Synthesis Model)的Transformer编码器架构提取多视角3D特征。引入SMPL-X结构先验，在规范姿态预测与每个SMPL-X顶点关联的3D高斯原语。每个顶点分配一个受约束的高斯(靠近SMPL-X表面)和若干无约束高斯(捕获偏离参数化表面的几何如衣物和头发)。使用神经纹理和姿态图像机制建立稀疏输入与3D几何的对应关系，通过可微分光栅化渲染到输入视角。动画通过线性混合蒙皮(LBS)直接驱动高斯，无需额外网络评估。 提出HumanGS，首个统一的前馈式框架，从多视角RGB和SMPL-X姿态实现人体3D重建和实时动画。每个SMPL-X顶点关联一组3D高斯原语，一次前向传播即生成可通过LBS直接动画化的化身，无需HumanRAM等方法的逐帧网络推理。在THuman 2.1、AvatarReX和THuman 4.0上达到SOTA可比质量，同时独有地支持实时动画。 依赖SMPL-X参数估计的准确性，对严重遮挡和极端姿态的处理有限，面部和手部细节不如专门的精细建模方法。
\textbf{方法信息：} 代码：\url{https://github.com/devdoot57/humangs}。
\textbf{设计动机与边界：} 相对原始 3DGS 以独立高斯和图像重建误差为主的设定，Real-Time Human Reconstruction and Animation using Feed-Forward Gaussian Splatting 把问题落在以下可实现的改动上：基于LVSM(Large View Synthesis Model)的Transformer编码器架构提取多视角3D特征。引入SMPL-X结构先验，在规范姿态预测与每个SMPL-X顶点关联的3D高斯原语。每个顶点分配一个受约束的高斯(靠近SMPL-X表面)和若干无约束高斯(捕获偏离参数化表面的几何如衣物和头发)。使用神经纹理和姿态图像机制建立稀疏输入与3D几何的对应关系，通过可微分光栅化渲染到输入视角。动画通过线性混合蒙皮(LBS)直接驱动高斯，无需额外网络评估。其设计目的不是单独增加一个模块，而是用上述表示、监督或渲染约束消除原流程在目标任务中的具体失配。提出HumanGS，首个统一的前馈式框架，从多视角RGB和SMPL-X姿态实现人体3D重建和实时动画。每个SMPL-X顶点关联一组3D高斯原语，一次前向传播即生成可通过LBS直接动画化的化身，无需HumanRAM等方法的逐帧网络推理。在THuman 2.1、AvatarReX和THuman 4.0上达到SOTA可比质量，同时独有地支持实时动画。适用时还应保留论文给出的边界：依赖SMPL-X参数估计的准确性，对严重遮挡和极端姿态的处理有限，面部和手部细节不如专门的精细建模方法。

\section{关节物体、人群与平台化方法}

这一节包含人体相邻但不完全等同的工作：铰接物体、人群渲染、人-物交互、头发、压缩和 Web 平台。它们的共同点是把高斯从“单个人体表面”推广到可移动部件、多个角色、物体接触或运行时系统。
% METHOD: 2410.22070
\paragraph{FreeGaussian（从光流微分发现可控部件）（2024.10）}
FreeGaussian~\cite{2410_22070} 进一步降低标注需求。它从连续帧光流和相机自运动解析推导 dynamic Gaussian flow，利用高斯轨迹和 HDBSCAN 聚类发现可交互结构，不需要 mask 或控制信号。控制方式也不使用传统一维关节参数，而是用 3D 球面向量表示高斯轨迹状态，与 splatting 光栅化管线直接对齐。项目实现围绕 flow derivatives。相对 SplArt 的两状态重建，FreeGaussian 更依赖视频运动微分来自动发现可动部件。
\textbf{实现细节：} FreeGaussian（2024.10） 是首个免标注的铰接物体可控重建方法，其核心创新在于通过微分分析解析推导动态高斯流（dynamic Gaussian flow，即每个高斯基元在连续帧间的 3D 运动向量场）。Baseline 与问题：此前方法需要 mask 标注才能识别动态区域。改了什么：FreeGaussian 建立了光流、相机运动与高斯流之间的数学关系（Lemma 1），从而将相机自运动（平移速度 $\mathbf{v}$ 和旋转速度 $\boldsymbol{\omega}$）与铰接运动在数学上完全解耦，无需 mask 即可自动发现可交互结构。具体流程为：先从连续帧间的光流和相机自运动推导动态高斯流，再通过 HDBSCAN 聚类获取运动轨迹实现部件分割。在此基础上，FreeGaussian 引入 3D 球面向量控制方案（3D spherical vector control），利用高斯轨迹作为状态表示直接在 3D 空间定义控制参数，与 splatting 光栅化管线天然对齐，替代了传统需要手动定义的 1D 控制信号。
\textbf{实现补充：} 实现上还要核对输入、输出、参数、特征、初始化、训练损失、正则、渲染或推理接口，以及优化中的剪枝、分裂、采样、编码或解码步骤；这些抓手决定该方法到底改了 3DGS 管线的哪一环。
\textbf{方法信息：} 代码：\url{https://github.com/tavish9/freegaussian}。
\textbf{设计动机与边界：} 相对原始 3DGS 以独立高斯和图像重建误差为主的设定，FreeGaussian 把问题落在以下可实现的改动上：方法核心是动态高斯流分析：1)从连续帧间的光流和相机自运动（平移速度v和旋转速度omega）通过微分分析解析推导动态高斯流（Lemma 1建立了光流、相机流和高斯流的关系）。2)通过追踪动态高斯流定位可交互的动态高斯，使用HDBSCAN聚类获取运动轨迹，消除对mask标注的依赖。3)3D球面向量控制方案：利用3D高斯的轨迹作为状态表示（球面向量），直接在3D空间定义控制参数，与splatting光栅化管线对齐，替代传统1D控制信号计算。4)训练时联合优化动态高斯流约束和光度损失。其设计目的不是单独增加一个模块，而是用上述表示、监督或渲染约束消除原流程在目标任务中的具体失配。提出FreeGaussian，首个无需标注的铰接物体可控3DGS重建方法。核心创新是通过微分分析从光流和相机运动解析推导动态高斯流，无需mask即可自动发现可交互结构。引入3D球面向量控制方案利用高斯轨迹作为状态表示，消除对复杂1D控制信号的依赖，简化可控高斯建模。适用时还应保留论文给出的边界：依赖光流估计质量；快速运动时光流可能不准确；多物体交互场景处理有限；复杂铰接（如多关节）仍有挑战。
% METHOD: 2501.17792
\paragraph{CrowdSplat（LOD 人群高斯渲染）（2025.01）}
CrowdSplat~\cite{2501_17792} 把高斯 avatar 用于大规模人群。重建阶段基于 GaussianAvatar 生成少量人体模板；合成阶段，多名角色共享同一模板的大部分高斯属性，只改变均值随姿态运动，以节省内存。LOD 策略按视距选择高斯数量，近处保留完整细节，远处使用少量高斯。CUDA 实现支持数千角色实时渲染。相对单人体 HUGS，CrowdSplat 的核心问题是“许多个高斯人体如何同时跑”。
\textbf{实现细节：} CrowdSplat框架分为两阶段：(1)头像重建：基于GaussianAvatar方法，在SMPL模板上构建头像，训练神经网络预测姿态相关的3DGS以模拟衣物褶皱等细节。从4个受试者的单目视频中重建14个头像模板。(2)人群合成：优化CUDA实现，对相同模板但不同姿态的角色复用除均值外的所有高斯参数（固定不透明度和旋转），减少内存约11
\textbf{方法信息：} 代码：\url{https://github.com/rockyxu66/crowdsplat}。
\textbf{设计动机与边界：} 相对原始 3DGS 以独立高斯和图像重建误差为主的设定，CrowdSplat 把问题落在以下可实现的改动上：CrowdSplat框架分为两阶段：(1)头像重建：基于GaussianAvatar方法，在SMPL模板上构建头像，训练神经网络预测姿态相关的3DGS以模拟衣物褶皱等细节。从4个受试者的单目视频中重建14个头像模板。(2)人群合成：优化CUDA实现，对相同模板但不同姿态的角色复用除均值外的所有高斯参数（固定不透明度和旋转），减少内存约11%。LOD策略根据视距调整高斯数量：5米内202k、5-10米12k、10米外3k。使用AMASS数据集中的15个运动序列驱动人群动画。其设计目的不是单独增加一个模块，而是用上述表示、监督或渲染约束消除原流程在目标任务中的具体失配。提出CrowdSplat框架，利用3DGS从单目视频重建可控制的人体头像，引入LOD策略（近处202k、中距12k、远距3k高斯）并优化CUDA实现以共享相同模板的高斯参数，实现3500个角色31FPS的实时渲染，内存优化约11%。适用时还应保留论文给出的边界：依赖SMPL模板，对非标准体型适配有限；LOD切换可能产生视觉跳变；仅测试了静态相机场景；运动多样性受限于AMASS数据集
% METHOD: 2502.14004
\paragraph{Inter3D（交互物体多状态基准与张量表示）（2025.02）}
Inter3D~\cite{2502_14004} 面向人机交互 3D 物体的新状态合成。具有 $n$ 个可动部件的物体理论上有 $2^n$ 种状态，逐状态训练高斯成本过高。方法用 Space Discrepancy Tensors 和多分辨率哈希编码表示状态差异，只需观测标准状态和单部件状态，就能合成未见组合状态。Mutual State Regularization 约束跨状态可动部件的空间密度一致性。相对 SPLATART，它更偏状态组合合成和评估。
\textbf{实现细节：} Inter3D针对具有n个可动部件的人机交互物体建模难题，提出Space Discrepancy Tensors结合多分辨率哈希编码(MHE)的方法。核心思想是：不使用 $2^n$ 个独立模型，而是用一套参数表示所有状态。具体而言，MHE初始化用于捕捉场景细节，Space Discrepancy Tensors则编码不同部件状态之间的差异。训练时仅观测标准状态和各部件的个体状态，组合状态完全未参与训练。为解决跨状态相机轨迹不一致的问题，提出Mutual State Regularization机制，通过约束不同状态下可动部件的空间密度一致性来提升合成质量。还探索了两种占用网格采样策略以提高训练效率。自建了包含Car、Display Case、Drawer、Furniture等日常交互物体的数据集，并设计了新的评估流水线。提出Inter3D基准和方法，用于人机交互3D物体的新状态合成。引入Space Discrepancy Tensors高效建模n个可动部件物体的$2^n$种状态，提出Mutual State Regularization机制增强可动部件的空间密度一致性，避免为每个状态训练独立模型的高昂开销。
\textbf{方法信息：} 代码：\url{https://github.com/inter3d-ui/inter3d}。
\textbf{设计动机与边界：} 相对原始 3DGS 以独立高斯和图像重建误差为主的设定，Inter3D 把问题落在以下可实现的改动上：Inter3D针对具有n个可动部件的人机交互物体建模难题，提出Space Discrepancy Tensors结合多分辨率哈希编码(MHE)的方法。核心思想是：不使用2\^{}n个独立模型，而是用一套参数表示所有状态。具体而言，MHE初始化用于捕捉场景细节，Space Discrepancy Tensors则编码不同部件状态之间的差异。训练时仅观测标准状态和各部件的个体状态，组合状态完全未参与训练。为解决跨状态相机轨迹不一致的问题，提出Mutual State Regularization机制，通过约束不同状态下可动部件的空间密度一致性来提升合成质量。还探索了两种占用网格采样策略以提高训练效率。自建了包含Car、Display Case、Drawer、Furniture等日常交互物体的数据集，并设计了新的评估流水线。其设计目的不是单独增加一个模块，而是用上述表示、监督或渲染约束消除原流程在目标任务中的具体失配。提出Inter3D基准和方法，用于人机交互3D物体的新状态合成。引入Space Discrepancy Tensors高效建模n个可动部件物体的2\^{}n种状态，提出Mutual State Regularization机制增强可动部件的空间密度一致性，避免为每个状态训练独立模型的高昂开销。适用时还应保留论文给出的边界：目前仅适用于离散状态的人机交互物体，对连续变形物体的扩展性未验证；自建数据集规模有限，泛化能力需进一步验证。
% METHOD: 2503.06677
\paragraph{REArtGS（SDF 正则的铰接物体重建与生成）（2025.03）}
REArtGS~\cite{2503_06677} 用 3DGS 重建铰接物体并生成未见状态。它先用 MLP 学习 SDF 场，通过 Sigmoid 转为高斯不透明度，并结合 GOF 渲染贡献设计 unbiased SDF regularization，使高斯中心更贴近真实表面。随后建立受运动学结构约束的变形场，把高斯从观测状态映射到任意状态。相对 SplArt 的 mobility 优化，REArtGS 更强调表面几何质量和连续未见状态生成。
\textbf{实现细节：} REArtGS框架分为重建和生成两个阶段：(1)无偏SDF引导重建——引入MLP学习SDF场，通过Sigmoid转换获得高斯不透明度，结合GOF的渲染贡献设计无偏SDF正则化。该正则化鼓励高斯中心处的SDF值趋近于零（即高斯位于表面上），当t接近t*时|S|收敛到零，增强几何约束。相比原始不透明度场，SDF提供了更严格的表面约束，产生更高质量的密集点云；(2)运动学约束变形场——基于重建得到的密集点云和运动学结构先验，建立时间连续的变形场。每个高斯基元通过变形场映射到任意状态s∈[0,1]的位置，变形受铰接运动学结构约束而非简单的时间输入。这使得无需对中间状态的监督即可生成未见状态的表面网格。 提出REArtGS框架，引入无偏SDF引导正则化高斯不透明度场增强几何约束，建立受铰接运动学结构约束的可变形场实现未见状态的无监督表面网格生成，仅用两个任意状态的多视角RGB图像即可实现高保真重建和生成。 需要两个不同状态的输入；对复杂多部件铰接结构的运动学建模仍有挑战；生成状态的范围受运动学先验限制。
\textbf{方法信息：} ByteDance 参与了这篇论文。



% METHOD: 2503.09342
\paragraph{GASPACHO（人-物交互可控高斯）（2025.03）}
GASPACHO~\cite{2503_09342} 同时恢复可动画人和物体模板。人用 SMPL-X 驱动高斯，物体高斯则学习在底层 2D surface manifold 上，而不是普通 3D 体空间中，从而保留薄结构和细节。方法还在高斯空间定义接触约束，正则化人-物关系，避免动画后穿透或接触不合理。相对 BIGS 这类手-物交互重建，GASPACHO 更关注重建后可控合成：人和物体作为独立高斯集合，可分别动画并重新组合。
\textbf{实现细节：} GASPACHO（2025.03） 修改了 3DGS 的物体表示与训练策略——用 2D 表面流形上的高斯图（Gaussian Maps）替代 3D 体积中的自由高斯学习，用遮挡感知的组合损失替代标准光度损失——基于 SMPL 的人体 LBS 蒙皮变形和基于 6D 刚体变换的物体运动建模保持不变。Baseline 与问题。读者已经从前文 SplArt 了解到铰接物体重建需要参数化运动模型。GASPACHO 面临的核心挑战更进一步：在人体-物体交互中，物体通常被人体严重遮挡（即使多相机设置下，大量视角中的物体仍不可见），直接使用标准 3DGS 优化物体时无法收敛；此外，先前方法要么将物体当作背景处理（无法动画化），要么假设物体静态不动，均无法实现可控制的人-物交互合成。改了什么。
\textbf{方法信息：} Huawei 参与了这篇论文。
\textbf{设计动机与边界：} 相对原始 3DGS 以独立高斯和图像重建误差为主的设定，GASPACHO 把问题落在以下可实现的改动上：方法包含三个核心设计：(1) 2D流形物体高斯：不同于传统在3D体空间中学习高斯，物体高斯在底层2D表面流形上学习，产生更锐利、更精细的物体细节。这种表示特别适合动态物体重建中保留薄结构和细小特征。(2) 接触约束：在高斯空间中定义接触约束来正则化人-物关系，确保动画时人物交互的物理合理性。约束基于高斯间的空间关系而非表面网格。(3) 可控渲染：人和物体作为独立的高斯集合，可分别动画化后重新组合，支持从新相机视角渲染新型人-物交互。使用SMPL-X参数化人体模型驱动人物高斯，物体高斯通过刚体/非刚体变形实现动画。其设计目的不是单独增加一个模块，而是用上述表示、监督或渲染约束消除原流程在目标任务中的具体失配。首个同时恢复可动画化的人和物体模板的3DGS方法，创新性地在2D表面流形上学习物体高斯实现更清晰的细节，引入接触约束实现物理合理的动画，支持可控的新人物交互合成。适用时还应保留论文给出的边界：对快速运动和大变形的处理仍有挑战；2D流形表示需要合理的初始化；多物体交互的支持有限。
% METHOD: 2506.03594
\paragraph{SplArt（自监督部件分割与铰接估计）（2025.06）}
SplArt~\cite{2506_03594} 从两组不同铰接状态 RGB 图像重建部件级物体。它把优化拆成三阶段：先对两个状态分别训练 3DGS，避免端到端联合优化陷入局部最优；再通过可微 mobility 参数估计每个高斯是否属于移动部件；最后联合优化旋转或平移关节参数与 mobility。几何自监督在两个状态之间建立一致性，不需要 3D 标注或类别先验。项目实现面向 articulation estimation。相对 REArtGS 的 SDF 表面约束，SplArt 更强调可微移动性带来的部件分割。
\textbf{实现细节：} SplArt（2025.06） 提出了自监督、类别无关的铰接物体重建框架。Baseline 与问题：此前的铰接物体重建方法需要部件级标注或已知铰接类型，且联合优化重建和铰接估计容易陷入局部最优。改了什么：SplArt 为每个高斯增加一个可微移动性参数（mobility parameter），实现精细的部件分割，并通过解耦的多阶段优化策略将重建、部件分割和铰接估计逐步处理。实现原理：方法假设物体由静态父部件和可移动子部件组成，子部件通过旋转或平移铰接相对于父部件运动。核心流程分三个阶段：阶段 1 对铰接的两个端状态分别独立优化标准 3DGS 模型；阶段 2 引入跨静态公式（cross-static formulation），通过比较两个独立重建的高斯属性差异，估计每个高斯的移动性参数（0 表示静态、1 表示可动），实现部件级的软分割；阶段 3 使用跨移动公式（cross-mobile formulation），利用移动性参数将可动部件的高斯从一个状态变换到另一个状态，同时优化铰接参数（旋转轴 $a$、枢轴点 $p$、旋转角度 $\theta$），并通过渲染损失精修移动性。
\textbf{方法信息：} 代码：\url{https://github.com/ripl/splart}。
\textbf{设计动机与边界：} 相对原始 3DGS 以独立高斯和图像重建误差为主的设定，SplArt 把问题落在以下可实现的改动上：方法流程（三阶段解耦优化）：(1) Stage 1 - 分离重建：对两个铰接状态分别独立优化3DGS模型，避免端到端联合优化的局部极小值问题。(2) Stage 2 - 移动性估计：通过跨静态公式估计每个高斯的移动性参数（可微分，范围0-1），实现静态和移动部件的精细化分割。移动性参数通过梯度优化，无需手动标注。(3) Stage 3 - 铰接估计与细化：基于跨移动公式联合优化铰接运动参数和移动性参数。支持旋转（轴+轴心+角度）和平移（轴+距离）两种铰接类型。利用几何自监督（两状态重建间的几何一致性约束）消除对3D标注或类别先验的需求。最终模型支持新视角和新铰接状态的实时光真实渲染，比NeRF-based方法快100倍以上。其设计目的不是单独增加一个模块，而是用上述表示、监督或渲染约束消除原流程在目标任务中的具体失配。提出自监督、类别无关的铰接物体3DGS重建框架，为每个高斯添加可微移动性参数实现精细化部件分割，多阶段优化策略避免局部最优，仅需两组不同铰接状态的RGB图像即可实现实时渲染。适用时还应保留论文给出的边界：假设已知铰接类型（旋转或平移）；目前仅支持两部分物体；需要两个极端铰接状态的观测。
% METHOD: 2506.12184
\paragraph{SPLATART（部件分离与关节估计解耦）（2025.06）}
SPLATART~\cite{2506_12184} 处理更深运动学树的铰接物体。它先训练带语义参数的 Splatfacto 高斯模型，用 Softmax 和交叉熵学习部件归属；再在不同配置的高斯之间估计部件 6DoF 姿态，最后从部件姿态推断关节轴、角度和树结构。这样部件分离与关节估计被解耦，允许支持七自由度机器人臂这类更深树结构。相对 ScrewSplat 的端到端螺旋轴优化，SPLATART 更强调可解释的分步 pipeline。
\textbf{实现细节：} SPLATART管线包含五个主要步骤：(1)逐场景语义高斯泼溅训练：使用Splatfacto模型训练带有语义参数的高斯泼溅，语义参数通过Softmax+交叉熵损失学习部件归属；(2)跨场景部件姿态估计：在不同配置的高斯泼溅间建立部件对应关系，估计每个部件的6DoF姿态；(3)关节参数估计：从估计的部件姿态中推断关节参数（关节轴、旋转角度等）；(4)树结构生成：构建部件间的运动学树，支持比以往更深的树结构（如7自由度机器人臂）；(5)新姿态/配置渲染：在用户指定配置下渲染物体，包括训练中未见过的配置。关键创新是将部件分离与铰接估计解耦，允许先完成部件分割再进行关节估计。 提出将部件分离与铰接估计解耦的管线，支持事后确定关节估计和更深运动学树的铰接物体表示，通过语义高斯泼溅训练和跨场景部件姿态估计实现复杂铰接物体重建。 目前仅假设两个观测状态；需要部分图像空间分割标注作为监督；对复杂关节类型（如螺旋关节）的支持有限。
\textbf{方法信息：} Amazon 参与了这篇论文。
\textbf{设计动机与边界：} 相对原始 3DGS 以独立高斯和图像重建误差为主的设定，SPLATART 把问题落在以下可实现的改动上：SPLATART管线包含五个主要步骤：(1)逐场景语义高斯泼溅训练：使用Splatfacto模型训练带有语义参数的高斯泼溅，语义参数通过Softmax+交叉熵损失学习部件归属；(2)跨场景部件姿态估计：在不同配置的高斯泼溅间建立部件对应关系，估计每个部件的6DoF姿态；(3)关节参数估计：从估计的部件姿态中推断关节参数（关节轴、旋转角度等）；(4)树结构生成：构建部件间的运动学树，支持比以往更深的树结构（如7自由度机器人臂）；(5)新姿态/配置渲染：在用户指定配置下渲染物体，包括训练中未见过的配置。关键创新是将部件分离与铰接估计解耦，允许先完成部件分割再进行关节估计。其设计目的不是单独增加一个模块，而是用上述表示、监督或渲染约束消除原流程在目标任务中的具体失配。提出将部件分离与铰接估计解耦的管线，支持事后确定关节估计和更深运动学树的铰接物体表示，通过语义高斯泼溅训练和跨场景部件姿态估计实现复杂铰接物体重建。适用时还应保留论文给出的边界：目前仅假设两个观测状态；需要部分图像空间分割标注作为监督；对复杂关节类型（如螺旋关节）的支持有限。
% METHOD: 2507.17029
\paragraph{StreamME（直播流中的在线头像重建）（2025.07）}
StreamME~\cite{2507_17029} 面向 live stream 中同步录制和重建 3D 头像。它完全去掉可变形 3DGS 常用的 MLP 运动建模，将高斯点云附着在追踪头部网格表面，随几何驱动直接运动。主点简化策略保留捕捉面部运动的关键锚点，减少点数；Gaussian Properties warm-up 用纹理和光照参数辅助几何重定向。系统可边直播边训练，并传输 3DGS 参数而不是完整图像。相对 GSAC 的离线 Unity 管线，StreamME 是实时流媒体化身系统。
\textbf{实现细节：} StreamME的核心设计：(1)无MLP运动建模：将3D高斯点云附着在追踪的头部网格表面，随网格变形移动，完全消除了传统可变形3DGS中的MLP层，显著提高表情适应速度；(2)主点简化策略：基于面部表面更稀疏地分布点云，在优化点数的同时保持渲染质量。通过运动感知锚点策略，保留捕捉面部运动的关键锚点；(3)Gaussian Properties热身：引入可学习纹理和光照参数辅助顶点变形，20秒内完成几何重定向；(4)on-the-fly训练：与直播流同步进行逐帧训练，无需预缓存数据。3DGS参数直接在流中传输而非完整图像（约70
\textbf{实现补充：} 实现上还要核对输入、输出、参数、特征、初始化、训练损失、正则、渲染或推理接口，以及优化中的剪枝、分裂、采样、编码或解码步骤；这些抓手决定该方法到底改了 3DGS 管线的哪一环。
\textbf{方法信息：} Adobe 参与了这篇论文。
\textbf{设计动机与边界：} 相对原始 3DGS 以独立高斯和图像重建误差为主的设定，StreamME 把问题落在以下可实现的改动上：StreamME的核心设计：(1)无MLP运动建模：将3D高斯点云附着在追踪的头部网格表面，随网格变形移动，完全消除了传统可变形3DGS中的MLP层，显著提高表情适应速度；(2)主点简化策略：基于面部表面更稀疏地分布点云，在优化点数的同时保持渲染质量。通过运动感知锚点策略，保留捕捉面部运动的关键锚点；(3)Gaussian Properties热身：引入可学习纹理和光照参数辅助顶点变形，20秒内完成几何重定向；(4)on-the-fly训练：与直播流同步进行逐帧训练，无需预缓存数据。3DGS参数直接在流中传输而非完整图像（约70%压缩），保护面部隐私并降低通信带宽。其设计目的不是单独增加一个模块，而是用上述表示、监督或渲染约束消除原流程在目标任务中的具体失配。提出StreamME，首个实现直播流中同步录制和重建3D头像的方法（on-the-fly训练）。基于3DGS但完全消除可变形3DGS中对MLP的依赖，仅依赖几何驱动运动。引入基于主点的简化策略，在保持渲染质量的同时优化点云数量。支持隐私保护（约70%压缩）和下游应用（动画、卡通化、重光照）。适用时还应保留论文给出的边界：依赖头部网格追踪质量；简化策略可能在极端表情下丢失细节；on-the-fly训练对GPU资源有实时要求。
% METHOD: 2508.02146
\paragraph{ScrewSplat（端到端螺旋轴铰接识别）（2025.08）}
ScrewSplat~\cite{2508_02146} 将铰接物体识别建成联合优化问题。每个高斯除几何外附加 part segmentation label；每个关节用 screw axis 参数化，连续表示旋转、平移或螺旋运动，并用置信度分数表示关节类型和数量不确定性。渲染损失重建外观，运动学约束让同一刚体部件运动一致，角度参数编码不同配置。相对 REArtGS++ 的平面高斯和 Taylor 正则，ScrewSplat 更强调从纯 RGB 到关节数量、类型和轴的端到端推断。
\textbf{实现细节：} ScrewSplat将铰接物体识别建模为联合优化问题，同时优化部分感知几何、关节轴/类型和关节角度。核心设计：(1)3DGS几何表示：使用3DGS重建物体外观和几何，为每个高斯基元附加分割标签实现部分分割；(2)螺旋轴模型：为每个运动关节参数化为螺旋轴（旋量），提供关节类型的连续参数化（旋转、平移、螺旋）。引入螺旋轴的置信度分数表示关节类型和数量的不确定性，避免离散优化；(3)联合优化：渲染损失驱动几何重建，运动学约束确保分割部分的运动一致性，角度参数编码不同配置下的运动状态；(4)优化过程中，螺旋轴从随机初始化收敛到正确的运动学结构。无需任何中间步骤（如点对应匹配或部分聚类），整体框架简洁高效。 提出ScrewSplat，首个纯RGB输入的端到端铰接物体识别方法。随机初始化螺旋轴并迭代优化恢复运动学结构，集成3DGS同时重建3D几何并分割刚体可移动部分。无需已知关节类型、关节数量或深度输入。实现SOTA识别精度并支持零样本文本引导操控。 复杂多关节物体（如多抽屉柜）的关节数量推断仍有挑战；依赖多视角RGB输入且物体需在不同配置下拍摄；优化可能陷入局部最优。
\textbf{方法信息：} Microsoft 参与了这篇论文。
\textbf{设计动机与边界：} 相对原始 3DGS 以独立高斯和图像重建误差为主的设定，ScrewSplat 把问题落在以下可实现的改动上：ScrewSplat将铰接物体识别建模为联合优化问题，同时优化部分感知几何、关节轴/类型和关节角度。核心设计：(1)3DGS几何表示：使用3DGS重建物体外观和几何，为每个高斯基元附加分割标签实现部分分割；(2)螺旋轴模型：为每个运动关节参数化为螺旋轴（旋量），提供关节类型的连续参数化（旋转、平移、螺旋）。引入螺旋轴的置信度分数表示关节类型和数量的不确定性，避免离散优化；(3)联合优化：渲染损失驱动几何重建，运动学约束确保分割部分的运动一致性，角度参数编码不同配置下的运动状态；(4)优化过程中，螺旋轴从随机初始化收敛到正确的运动学结构。无需任何中间步骤（如点对应匹配或部分聚类），整体框架简洁高效。其设计目的不是单独增加一个模块，而是用上述表示、监督或渲染约束消除原流程在目标任务中的具体失配。提出ScrewSplat，首个纯RGB输入的端到端铰接物体识别方法。随机初始化螺旋轴并迭代优化恢复运动学结构，集成3DGS同时重建3D几何并分割刚体可移动部分。无需已知关节类型、关节数量或深度输入。实现SOTA识别精度并支持零样本文本引导操控。适用时还应保留论文给出的边界：复杂多关节物体（如多抽屉柜）的关节数量推断仍有挑战；依赖多视角RGB输入且物体需在不同配置下拍摄；优化可能陷入局部最优。
% METHOD: 2509.07774
\paragraph{HairGS（丝级头发几何高斯）（2025.09）}
HairGS~\cite{2509_07774} 把 3DGS 扩展到 strand-level hair reconstruction。第一阶段用可微高斯光栅化和自适应 densification 得到密集头发几何，并用 Gabor 方向图加入双向方向损失。第二阶段根据距离和角度启发式把高斯段合并成长丝束；第三阶段用光度损失和平滑正则细化关节位置，输出变量长度的 3D strand 集合。项目实现围绕 hair strand reconstruction。相对 MeGA 的头发高斯渲染，HairGS 的目标更几何化：恢复可用的头发丝拓扑。
\textbf{实现细节：} HairGS（2025.09） 回到静态重建路线，但与 Gaussian Haircut 的策略形成鲜明对比。锚点回顾：Gaussian Haircut 的核心思路是"先建非结构化高斯 → 再拟合结构化发丝"，发丝是从高斯几何中间接恢复的。HairGS 同样追求发丝级重建，但走了不同的路径：它直接将 3D 高斯作为发丝几何的基元，通过合并（merging）将短段拼接为长发丝，而非从高斯拟合到独立的发丝参数化。改了什么：HairGS 提出三阶段优化管线，每个阶段都围绕 3DGS 的可微分光栅化器构建，全程无需预训练数据。标准 3DGS 自适应致密化机制在此阶段确保几何充分覆盖。(2) 发丝生成阶段：第一阶段的高斯被转化为由两个关节点构成的短发丝段，每段用圆柱体高斯表示——中心取端点中值，旋转用 Rodrigues 公式将 $x$-轴对齐到段方向 $R = I + K + K^{}2/(1 + \mathbf{v} \cdot d)$，缩放为 $(\|p_{j+1} - p_j\|_2, \tau_j, \tau_j)$，其中 $\tau_j$ 是可学习的厚度参数。
\textbf{方法信息：} 代码：\url{https://github.com/yimin-pan/hair-gs}。
\textbf{设计动机与边界：} 相对原始 3DGS 以独立高斯和图像重建误差为主的设定，HairGS 把问题落在以下可实现的改动上：HairGS多阶段管线：(1)阶段I-几何重建：使用可微光栅化器和自适应致密化获取头发几何的密集详细表示。引入掩膜值限制合并过程在有效头发区域，使用0阶SH避免视角依赖颜色。优化使用L1+DSSIM光度损失和双向方向损失(基于Gabor滤波器的2D方向图)。(2)阶段II-丝束生成：基于距离和角度启发式的合并方案，将分离的头发段组合为更长的连贯丝束。不需要预设每条丝的点数，依靠优化动态确定。(3)阶段III-细化：使用光度损失和平滑正则化细化关节位置，结合方向约束和角度正则化防止形成尖锐角度。最终输出为变量长度链接3D点的丝束集合。预处理使用COLMAP获取相机姿态和FLAME模型顶点作为初始化。其设计目的不是单独增加一个模块，而是用上述表示、监督或渲染约束消除原流程在目标任务中的具体失配。首次将3DGS扩展到丝级别头发几何重建，提出多阶段管线：可微高斯光栅化重建几何，基于距离和角度启发式的合并方案将高斯段合并为连贯丝束，以及光度监督下的丝束细化和生长。还提出评估丝束拓扑准确性的新指标。适用时还应保留论文给出的边界：对初始化质量敏感；极端卷曲发型的合并可能失败；遮挡区域的重建依赖多视角覆盖。
% METHOD: 2510.10492
\paragraph{Human Avatar Compression（先验引导的超低码率压缩）（2025.10）}
Towards Efficient 3D Gaussian Human Avatar Compression~\cite{2510_10492} 处理人体高斯视频传输。它把外观和时间演化解耦：规范高斯 avatar 只需压缩一次，逐帧只传很少的时间参数，例如 SMPL 参数，再通过 LBS 从规范化身变形到目标姿态。训练规范化身时使用人体先验模板和 articulated splatting。相对 PointSplat 的前馈紧凑预测，这里更像编码框架：在超低码率下复用同一个高斯外观，把运动压成人体参数。
\textbf{实现细节：} 框架首先以无网络方式训练规范高斯化身，使用铰接溅射进行外观建模。同时利用人体先验模板通过紧凑参数化表示捕捉时间体运动。外观与时间演化的解耦最小化了冗余：规范化身在序列中共享，仅需一次压缩；时间参数由每帧仅94个参数组成（如SMPL参数），传输码率极低。每帧的目标人体化身通过线性混合蒙皮（LBS）变换从规范化身形变生成，实现时间连贯的视频重建和新视角合成。该方法将人体先验知识有效融入3DGS压缩框架。 提出了一个高效3D化身编码框架，利用紧凑人体先验和规范-目标变换实现超低码率下高质量3D人体化身视频压缩。通过将外观和时间演化解耦，规范化身在序列中仅需压缩一次，时间参数每帧仅94个参数，显著降低了码率需求。 仅适用于人体化身的压缩，对一般动态场景不适用；依赖人体先验模板（如SMPL）的准确性；对宽松衣物的建模可能受限。
\textbf{实现补充：} 实现上还要核对输入数据、输出指标、基线配置、评测协议、系统接口和适用边界；若调用 3DGS 渲染或重建结果，也应说明相机参数、训练设置、优化流程和指标计算如何对齐。
\textbf{方法信息：} Alibaba 参与了这篇论文。
\textbf{设计动机与边界：} 相对原始 3DGS 以独立高斯和图像重建误差为主的设定，Towards Efficient 3D Gaussian Human Avatar Compression 把问题落在以下可实现的改动上：框架首先以无网络方式训练规范高斯化身，使用铰接溅射进行外观建模。同时利用人体先验模板通过紧凑参数化表示捕捉时间体运动。外观与时间演化的解耦最小化了冗余：规范化身在序列中共享，仅需一次压缩；时间参数由每帧仅94个参数组成（如SMPL参数），传输码率极低。每帧的目标人体化身通过线性混合蒙皮（LBS）变换从规范化身形变生成，实现时间连贯的视频重建和新视角合成。该方法将人体先验知识有效融入3DGS压缩框架。其设计目的不是单独增加一个模块，而是用上述表示、监督或渲染约束消除原流程在目标任务中的具体失配。提出了一个高效3D化身编码框架，利用紧凑人体先验和规范-目标变换实现超低码率下高质量3D人体化身视频压缩。通过将外观和时间演化解耦，规范化身在序列中仅需压缩一次，时间参数每帧仅94个参数，显著降低了码率需求。适用时还应保留论文给出的边界：仅适用于人体化身的压缩，对一般动态场景不适用；依赖人体先验模板（如SMPL）的准确性；对宽松衣物的建模可能受限。
% METHOD: 2511.17059
\paragraph{REArtGS++（螺旋运动与平面高斯约束）（2025.11）}
REArtGS++~\cite{2511_17059} 去掉关节类型先验，用 decoupled screw motion 同时表示旋转和平移。每个高斯的运动由部件运动混合得到，混合权重来自零件分割概率。它还把 3D Gaussian 压成 planar Gaussian，以获得可靠法线和无偏深度图，并用 Taylor 一阶展开在平面法线和深度之间建立时间一致几何正则。相对 REArtGS，它的增量是关节类型未知时的统一运动参数化，以及更强的连续时间几何约束。
\textbf{实现细节：} REArtGS++包含三个核心设计：(1)解耦螺旋运动建模：将每个关节参数分解为旋转(θ,a)和平移(t,o)，无需预先确定关节类型。Gaussian运动通过零件运动混合推导，以零件分割概率为混合权重；(2)平面Gaussian表示：引入缩放损失将3D Gaussian压扁为2D平面，获得准确法线和无偏深度图，便于后续时间几何约束。每个平面Gaussian附加分割掩码表示属于各零件的概率，通过Mahalanobis距离和MLP计算；(3)时间几何约束：通过Taylor一阶展开在平面法线和深度间建立时间一致正则化，将两状态监督的几何约束扩展到连续时间域，增强关节参数估计和零件级重建。 提出REArtGS++，通过解耦螺旋运动建模消除关节类型先验需求，通过零件运动混合联合优化零件感知Gaussian和关节参数。引入平面Gaussian约束和时间一致正则化，通过Taylor一阶展开在平面法线和深度间建立时间连续几何约束。在合成和真实铰接物体上优于现有方法。 需要两状态的多视角RGB输入；对复杂多自由度铰接物体的建模仍有限制；零件数量需要预先指定或估计。
\textbf{方法信息：} Tencent 参与了这篇论文。
\textbf{设计动机与边界：} 相对原始 3DGS 以独立高斯和图像重建误差为主的设定，REArtGS++ 把问题落在以下可实现的改动上：REArtGS++包含三个核心设计：(1)解耦螺旋运动建模：将每个关节参数分解为旋转(θ,a)和平移(t,o)，无需预先确定关节类型。Gaussian运动通过零件运动混合推导，以零件分割概率为混合权重；(2)平面Gaussian表示：引入缩放损失将3D Gaussian压扁为2D平面，获得准确法线和无偏深度图，便于后续时间几何约束。每个平面Gaussian附加分割掩码表示属于各零件的概率，通过Mahalanobis距离和MLP计算；(3)时间几何约束：通过Taylor一阶展开在平面法线和深度间建立时间一致正则化，将两状态监督的几何约束扩展到连续时间域，增强关节参数估计和零件级重建。其设计目的不是单独增加一个模块，而是用上述表示、监督或渲染约束消除原流程在目标任务中的具体失配。提出REArtGS++，通过解耦螺旋运动建模消除关节类型先验需求，通过零件运动混合联合优化零件感知Gaussian和关节参数。引入平面Gaussian约束和时间一致正则化，通过Taylor一阶展开在平面法线和深度间建立时间连续几何约束。在合成和真实铰接物体上优于现有方法。适用时还应保留论文给出的边界：需要两状态的多视角RGB输入；对复杂多自由度铰接物体的建模仍有限制；零件数量需要预先指定或估计。
% METHOD: 2511.17092
\paragraph{SPAGS（单状态稀疏视角铰接物体）（2025.11）}

SPAGS~\cite{2511_17092} 面向单状态稀疏视角下的铰接物体重建。它先用 3D 生成模型获得初始 Gaussian，并建立 Gaussian information field 为每个高斯估计可靠性，通过路径积分选择最有信息量的稀疏视角。表示上，它把 3D Gaussian 压成 planar Gaussian，以获得更准确法线和深度。优化从深度平滑正则的粗阶段过渡到少样本扩散先验的细阶段，最后用 VLM 做开放词汇部件分割和关节参数估计。相对 SplArt 的双状态输入，SPAGS 更强调单状态主动视角和开放词汇零件理解。
\textbf{设计动机与边界：} 相对原始 3DGS 以独立高斯和图像重建误差为主的设定，SPAGS 把问题落在以下可实现的改动上：SPAGS框架包含四个核心模块：(1)最优视角感知：建立Gaussian信息场(GIF)，为每个Gaussian引入可靠性参数P，定义信息场强度(IFI)通过沿相机射线的路径积分计算，迭代选择最大化信息潜力的视角。利用3D生成模型SPAR3D获取初始Gaussian并配准；(2)平面Gaussian表示：通过缩放损失将3D Gaussian压扁为平面，获得准确法线和深度估计；(3)粗到细优化：粗阶段使用深度平滑正则化，细阶段引入少样本扩散先验提升细节；(4)开放词汇零件分割：利用视觉语言模型(VLM)通过视觉提示实现开放词汇零件分割，然后估计关节参数进行零件级网格提取。其设计目的不是单独增加一个模块，而是用上述表示、监督或渲染约束消除原流程在目标任务中的具体失配。提出SPAGS，首个无需类别先验的单状态稀疏视角铰接物体重建框架。引入Gaussian信息场感知最优稀疏视角，将3D Gaussian约束为平面原语实现精确法线和深度估计，通过粗到细策略配合深度平滑和扩散先验优化，利用VLM实现开放词汇零件分割和关节参数估计。适用时还应保留论文给出的边界：依赖VLM的分割质量；对非常规形状铰接物体效果受限；最优视角感知增加预处理时间。
\textbf{方法信息：} Tencent 参与了这篇论文；代码：\url{https://github.com/nerficg-project/SPaGS}。



% METHOD: 2512.05113
\paragraph{Splannequin（单目假人挑战视频冻结渲染）（2025.12）}
Splannequin~\cite{2512_05113} 是动态高斯训练正则化方法。单目 Mannequin Challenge 视频中，相机向前运动会导致部分高斯在某些时间隐藏或监督不足，冻结渲染时产生鬼影和模糊。方法用双检测机制把这些高斯分成 hidden 和 defective：hidden 锚到过去良好观测状态，defective 锚到未来强监督状态；一致性损失按时间距离衰减。它不改动态高斯架构，只增加训练损失。相对 4DGS 直接建模动态，Splannequin 解决的是单目冻结场景中监督时序不均的问题。
\textbf{实现细节：} Splannequin是针对单目Mannequin Challenge视频冻结渲染的正则化框架：（1）问题分析：在单目动态高斯中，当相机主要向前运动时，某些高斯在特定时间戳未被观测（隐藏）或监督不足（缺陷），导致冻结渲染时出现鬼影和模糊。（2）双检测机制：将每个监督不足的高斯分类为隐藏（已被相机移过，当前不可见）或缺陷（相机尚未到达，未来才有强监督）。（3）时间锚定正则化：隐藏高斯锚定到过去良好观测状态（$t_{ref}<t$），缺陷高斯锚定到未来强监督状态（$t_{ref}>t$）。一致性损失可写为 $L_{\mathrm{consistency}}(k,t)=\lambda(t,t_{\mathrm{ref}})D(\theta_k(t),\theta_k(t_{\mathrm{ref}}))$，其中 $\lambda=\exp(-|t-t_{\mathrm{ref}}|)$ 按时间距离指数衰减加权。随机采样参考时间戳避免昂贵的锚点搜索。（4）零推理开销：仅作为训练损失项添加，不修改架构。 提出双检测正则化方法，识别动态高斯中的隐藏和缺陷状态并分别锚定到良好监督的参考状态，96
\textbf{实现补充：} 为避免只把它读成名称或数据点，还应把实现接口明确拆开：输入是什么，输出是高斯参数、图像、轨迹、标注还是评测指标，哪些参数、网络模块、训练损失或渲染接口真正接到原始 3DGS / Scaffold-GS / 4D-GS 管线上，以及剪枝、分裂、采样、解码、优化或评测协议在哪一步发生。若该条目是survey、benchmark 或 dataset，这里的“实现”对应论文集合、任务划分、相机/轨迹/标注组织、评价指标和基线配置，而不是虚构新的 rasterizer；若是算法条目，则对应高斯位置、尺度、旋转、不透明度、特征或 MLP 解码如何进入训练和渲染流程。
\textbf{方法信息：} NVIDIA 参与了这篇论文。



% METHOD: 2512.08478
\paragraph{Visionary（WebGPU 高斯运行时平台）（2025.12）}

Visionary~\cite{2512_08478} 不是具体人体重建算法，而是 Web-native 3DGS runtime。它用 WebGPU 做计算和图形后端，并定义 Gaussian Generator 合约：ONNX I/O 模式和元数据规定每帧如何生成或更新高斯位置、尺度、方向和颜色。这样 MLP-based 3DGS、4DGS、神经化身、风格迁移和增强网络都可作为插件在浏览器内运行。它的代码公开且 GitHub 关注度较高，说明这类浏览器端运行时已有一定实践基础。相对 StreamME 的单一流媒体头像，Visionary 提供的是平台层抽象：让动态高斯 avatar 进入 Web 应用。
\textbf{实现细节：} Visionary平台的核心设计：(1) WebGPU渲染器：利用WebGPU作为统一计算和图形后端，通过GPU-based原始排序实现高效的3DGS渲染，在相同3DGS资产下渲染效率优于现有Web查看器。(2) Gaussian Generator合约：定义严格的ONNX I/O模式和元数据，规范算法如何生成或更新高斯属性（位置、缩放、方向、颜色）。任何3DGS系列方法可作为插件加载，无需修改底层运行时。(3) 每帧ONNX推理：在浏览器中直接运行ONNX模型推理，支持动态神经处理和前馈生成后处理。移除原生依赖开销，实现"点击即运行"的浏览器体验。(4) 已支持多种变体：MLP-based 3DGS、4DGS、神经化身、风格变换和增强网络。(5) 提供concise TypeScript API库，便于集成到现有Web应用。统一了重建和生成两种范式的World Model Carrier。
\textbf{设计动机与边界：} 相对原始 3DGS 以独立高斯和图像重建误差为主的设定，Visionary 把问题落在以下可实现的改动上：Visionary平台的核心设计：(1) WebGPU渲染器：利用WebGPU作为统一计算和图形后端，通过GPU-based原始排序实现高效的3DGS渲染，在相同3DGS资产下渲染效率优于现有Web查看器。(2) Gaussian Generator合约：定义严格的ONNX I/O模式和元数据，规范算法如何生成或更新高斯属性（位置、缩放、方向、颜色）。任何3DGS系列方法可作为插件加载，无需修改底层运行时。(3) 每帧ONNX推理：在浏览器中直接运行ONNX模型推理，支持动态神经处理和前馈生成后处理。移除原生依赖开销，实现"点击即运行"的浏览器体验。(4) 已支持多种变体：MLP-based 3DGS、4DGS、神经化身、风格变换和增强网络。(5) 提供concise TypeScript API库，便于集成到现有Web应用。统一了重建和生成两种范式的World Model Carrier。其设计目的不是单独增加一个模块，而是用上述表示、监督或渲染约束消除原流程在目标任务中的具体失配。提出Visionary，首个基于WebGPU的开源Web原生3DGS渲染平台，通过标准化Gaussian Generator合约和每帧ONNX推理，统一支持多种3DGS变体（MLP-3DGS、4DGS、神经化身、风格迁移等）的浏览器内实时渲染和推理，显著降低部署门槛。适用时还应保留论文给出的边界：WebGPU浏览器兼容性仍有限制；大型场景的浏览器内存管理挑战；ONNX模型转换需要额外工作。
\textbf{方法信息：} 代码：\url{https://github.com/visionary-laboratory/visionary}；GitHub 关注度高，约 511 stars。



% METHOD: 2603.02893
\paragraph{MVD-HuGaS（单图扩散多视角到人体高斯）（2026.03）}
Intrinsic Geometry-Appearance Consistency Optimization for Sparse-View Gaussian Splatting~\cite{2603_02893} 也称 MVD-HuGaS。它从单张人体图像生成多视角图像，再优化高斯。多视角扩散模型在人体数据上微调，并加入 3D 几何和人体结构先验；alignment module 从生成的稀疏多视角中推断相机位姿，支持高斯和相机联合优化；面部畸变缓解模块基于深度精化生成的人脸区域。相对 HuGDiffusion 直接扩散高斯属性，MVD-HuGaS 先生成可用多视图和相机，再走传统 3DGS 优化。
\textbf{实现细节：} 方法流程：(1)多视角生成：在高质量3D人体数据集上微调多视角扩散模型，融入3D几何先验和人体结构先验，从单张参考图像生成多视角图像；(2)相机位姿推断：引入对齐模块从生成的稀疏多视角图像推断准确相机位姿，支持3DGS和相机位姿的联合优化；(3)面部畸变缓解：基于深度的面部畸变缓解模块精化生成面部的面部区域，提升整体重建保真度；(4)3DGS优化：利用精化后的多视角图像及其准确相机位姿优化目标人体的3D Gaussians，实现高保真自由视角渲染。 提出MVD-HuGaS，通过增强的多视角人体扩散模型从单张图像生成多视角图像，结合对齐模块联合优化3D Gaussian和相机位姿，引入基于深度的面部畸变缓解模块，在Thuman2.0和2K2K数据集上达到SOTA单视角3D人体渲染性能。 多视角扩散模型的生成质量限制了最终重建上限。面部区域修复可能在极端姿态下失效。
\textbf{实现补充：} 实现上还要核对输入、输出、参数、特征、初始化、训练损失、正则、渲染或推理接口，以及优化中的剪枝、分裂、采样、编码或解码步骤；这些抓手决定该方法到底改了 3DGS 管线的哪一环。
\textbf{方法信息：} Alibaba 参与了这篇论文。
\textbf{设计动机与边界：} 相对原始 3DGS 以独立高斯和图像重建误差为主的设定，Intrinsic Geometry-Appearance Consistency Optimization for Sparse-View Gaussian Splatting 把问题落在以下可实现的改动上：方法流程：(1)多视角生成：在高质量3D人体数据集上微调多视角扩散模型，融入3D几何先验和人体结构先验，从单张参考图像生成多视角图像；(2)相机位姿推断：引入对齐模块从生成的稀疏多视角图像推断准确相机位姿，支持3DGS和相机位姿的联合优化；(3)面部畸变缓解：基于深度的面部畸变缓解模块精化生成面部的面部区域，提升整体重建保真度；(4)3DGS优化：利用精化后的多视角图像及其准确相机位姿优化目标人体的3D Gaussians，实现高保真自由视角渲染。其设计目的不是单独增加一个模块，而是用上述表示、监督或渲染约束消除原流程在目标任务中的具体失配。提出MVD-HuGaS，通过增强的多视角人体扩散模型从单张图像生成多视角图像，结合对齐模块联合优化3D Gaussian和相机位姿，引入基于深度的面部畸变缓解模块，在Thuman2.0和2K2K数据集上达到SOTA单视角3D人体渲染性能。适用时还应保留论文给出的边界：多视角扩散模型的生成质量限制了最终重建上限。面部区域修复可能在极端姿态下失效。
% METHOD: 2604.25466
\paragraph{GHGS-MVSC（多视角语义一致人体高斯）（2026.04）}
Generalizable Human Gaussian Splatting via Multi-view Semantic Consistency~\cite{2604_25466} 面向稀疏视角泛化人体渲染。它先用 VGGT 编码多视角输入并预测深度，把语义特征和潜在嵌入反投影到共享 3D 空间；再在 3D 空间中找邻近嵌入对，并用 DINO 特征判断是否属于同一身体部位，通过 cross-view attention 重新校准潜在嵌入。最后从校准嵌入回归高斯参数。相对 GHG 的 UV 回归，它在 3D 空间中处理跨视角语义一致性。
\textbf{实现细节：} 方法包含三个关键步骤：(1) 潜在嵌入的3D反投影：使用预训练VGGT编码器(DINO+多视角Transformer)编码多视角输入，通过DPT-style解码器预测深度图，将语义特征和潜在嵌入反投影到共享3D空间。(2) 基于语义一致性的跨视角重新校准：在3D空间中找到空间邻近的嵌入对，使用DINO特征判断是否属于相同身体部位，通过跨视角注意力重新校准潜在嵌入。(3) 高斯回归：从重新校准的嵌入直接回归3DGS参数(位置、协方差、颜色、不透明度)。关键创新是在3D空间而非2D图像域建立跨视角关系。 提出利用3D空间中的语义一致性重新校准多视角潜在嵌入，通过跨视角注意力解决人体高纹理区域和遮挡部位的空间歧义，实现更准确的3D高斯定位和更高的稀疏视角人体渲染质量。 依赖SMPL参数估计的准确性；对极端姿态和严重遮挡的处理仍有不足；推理速度未充分优化。
\textbf{方法信息：} 代码：\url{https://github.com/dcvl-3d/ghgs-mvsc_release}。
\textbf{设计动机与边界：} 相对原始 3DGS 以独立高斯和图像重建误差为主的设定，Generalizable Human Gaussian Splatting via Multi-view Semantic Consistency 把问题落在以下可实现的改动上：方法包含三个关键步骤：(1) 潜在嵌入的3D反投影：使用预训练VGGT编码器(DINO+多视角Transformer)编码多视角输入，通过DPT-style解码器预测深度图，将语义特征和潜在嵌入反投影到共享3D空间。(2) 基于语义一致性的跨视角重新校准：在3D空间中找到空间邻近的嵌入对，使用DINO特征判断是否属于相同身体部位，通过跨视角注意力重新校准潜在嵌入。(3) 高斯回归：从重新校准的嵌入直接回归3DGS参数(位置、协方差、颜色、不透明度)。关键创新是在3D空间而非2D图像域建立跨视角关系。其设计目的不是单独增加一个模块，而是用上述表示、监督或渲染约束消除原流程在目标任务中的具体失配。提出利用3D空间中的语义一致性重新校准多视角潜在嵌入，通过跨视角注意力解决人体高纹理区域和遮挡部位的空间歧义，实现更准确的3D高斯定位和更高的稀疏视角人体渲染质量。适用时还应保留论文给出的边界：依赖SMPL参数估计的准确性；对极端姿态和严重遮挡的处理仍有不足；推理速度未充分优化。
% METHOD: 2605.24304
\paragraph{ArtSplat（前馈铰接高斯与逐像素关节图）（2026.05）}
ArtSplat~\cite{2605_24304} 从稀疏多状态、无标定视图直接前馈重建铰接 3DGS。它扩展 VGGT，用 DINOv2 patch 编码、相机和状态 tokens、交替帧注意力与全局注意力处理两个铰接状态。核心输出是逐像素 11 维 articulation map，包括关节类型、轴、枢轴、旋转角和位移，完全可微。每个像素同时预测深度、高斯参数和关节参数，形成状态条件高斯。相对 SPLATART 的分步关节估计，ArtSplat 更像端到端前馈版本。
\textbf{实现细节：} ArtSplat（2026.05） 是首个前馈式铰接 3DGS 框架，从稀疏多视角多状态图像在单次前向传播中同时重建几何和关节参数，速度比基线快 400 倍以上。改了什么：ArtSplat 基于 VGGT 架构扩展，将关节参数建模为 11 维逐像素密集表示（关节类型 logits、轴方向、枢轴位置、旋转角度、平移位移），完全可微且端到端训练。输入为覆盖两个铰接状态的稀疏多视角图像，经 DINOv2（自监督视觉Transformer）patch 化编码后加入可学习的相机和状态 token，通过交替帧注意力和全局注意力处理。关键设计是跨状态注意力（cross-state attention），带状态 token 的注意力块捕获两个铰接状态间的离散运动差异。每个像素同时预测深度、高斯参数和关节参数，形成状态条件化高斯集。
\textbf{技术实现补充：} ArtSplat（2026.05） 是首个前馈式铰接 3DGS 框架，从稀疏多视角多状态图像在单次前向传播中同时重建几何和关节参数，速度比基线快 400 倍以上。改了什么：ArtSplat 基于 VGGT 架构扩展，将关节参数建模为 11 维逐像素密集表示（关节类型 logits、轴方向、枢轴位置、旋转角度、平移位移），完全可微且端到端训练。输入为覆盖两个铰接状态的稀疏多视角图像，经 DINOv2（自监督视觉Transformer）patch 化编码后加入可学习的相机和状态 token，通过交替帧注意力和全局注意力处理。关键设计是跨状态注意力（cross-state attention），带状态 token 的注意力块捕获两个铰接状态间的离散运动差异。每个像素同时预测深度、高斯参数和关节参数，形成状态条件化高斯集。
\textbf{方法信息：} 代码：\url{https://github.com/deniz-dentsu/artsplat-3D}。
\textbf{设计动机与边界：} 相对原始 3DGS 以独立高斯和图像重建误差为主的设定，ArtSplat 把问题落在以下可实现的改动上：ArtSplat基于VGGT架构扩展：(1)输入：稀疏多视图图像覆盖两个铰接状态。(2)特征编码：DINOv2 patch化编码+可学习相机和状态token，经L层交替帧注意力和全局注意力处理。(3)逐像素关节图：将关节参数建模为11维逐像素密集表示，包含不变组（关节类型logits、轴方向、枢轴位置）和可变组（旋转角度、平移位移），完全可微且端到端训练。(4)跨状态注意力(CSA)：带状态token的跨状态注意力块捕获两个铰接状态间的离散运动差异。(5)输出：每个像素同时预测深度、高斯参数和关节参数，形成状态条件化高斯集。其设计目的不是单独增加一个模块，而是用上述表示、监督或渲染约束消除原流程在目标任务中的具体失配。首个前向传播铰接3DGS重建框架，逐像素关节图表示实现端到端可微的关节参数预测，比基线快400倍以上。适用时还应保留论文给出的边界：目前仅支持两个铰接状态输入；关节类型限于旋转和滑动；复杂多关节物体的精度有提升空间

\section{说话人头像与语音驱动}

语音驱动头像的锚点是“音频到运动”。早期 NeRF 或高斯 talking head 常直接改颜色和密度，容易造成嘴部和眼睑扭曲；更稳定的路线是建立持久高斯结构，用音频、情感、AU 或 FLAME 参数只驱动几何变形。另一条路线把说话人扩展到全身或移动设备，要求模型足够轻。
% METHOD: 2404.15264
\paragraph{TalkingGaussian（持久高斯结构与面口分支）（2024.04）}
TalkingGaussian~\cite{2404_15264} 是语音驱动 3DGS 头像的一个清晰锚点。它针对的是早期 talking head 方法常见的外观编辑问题：如果每帧直接改高斯颜色、密度或 radiance，嘴唇、牙齿、眼睑这类区域的外观变化并不平滑，网络容易把运动解释成颜色跳变，产生嘴部混乱、眼睑透明或局部扭曲。TalkingGaussian 的改变是把头部建立为持久高斯结构，颜色和不透明度作为身份与静态外观尽量保持稳定，说话运动主要通过高斯位置、尺度和旋转的连续变形来表达。

实现上，它不是训练一个统一运动场去覆盖整张脸，而是把面部外侧和口腔内部分成两个分支。代码中的训练流程也体现了这一点：先分别训练 face 分支和 mouth 分支，再在 fuse 阶段把两套高斯和运动网络合成。面部分支处理表情、脸颊、眼周等较连续的运动；口腔分支处理嘴唇开合、牙齿和舌头等快速且不连续的区域。推理时，音频特征经运动网络预测每帧高斯变形，两个分支分别渲染后再按 alpha 与 mask 融合。相对 GaussianTalker 等通过共享隐式特征控制高斯属性的路线，TalkingGaussian 的锚点意义在于把“学外观变化”改成“学持久结构上的几何运动”，并明确把口腔作为独立难点拆出来。

论文在 HDTF、VOCASET 等数据集上报告高质量唇同步渲染和 100+ FPS 级实时性。代码发布在 \url{https://github.com/fictionarry/talkinggaussian}，GitHub 约 383 stars；Semantic Scholar 约 79 次引用。
\textbf{实现细节：} TalkingGaussian（2024.04） 同样采用了基于变形的范式，但提出了与 GaussianTalker 不同的解耦策略。Baseline 与问题：直接在动态面部区域生成逐点外观（appearance generation）时，外观空间的跳变和不连续性使得连续平滑的神经场难以精确拟合，导致面部扭曲伪影。改了什么：TalkingGaussian 纯粹使用变形（deformation）来表示面部运动，通过平滑连续的变形驱动持久的高斯基元，而非直接修改点的外观颜色。实现原理：关键发现是面部运动和口腔运动在空间和时间上具有不同的特性：面部表情变化是渐进且空间连续的，而口腔（嘴唇开合、牙齿暴露）是快速且空间不连续的。为此，该方法将模型分解为面部和口腔两个独立分支，每个分支有独立的变形场。面部分支专注于全局表情变化和唇形同步，口腔分支专注于嘴内部的几何变化。设计动机：面部和口腔的解耦避免了单一变形场需要同时学习两种截然不同的运动模式，显著提高了各自的拟合精度。
\textbf{实现补充：} 实现上还要核对输入、输出、参数、特征、初始化、训练损失、正则、渲染或推理接口，以及优化中的剪枝、分裂、采样、编码或解码步骤；这些抓手决定该方法到底改了 3DGS 管线的哪一环。
\textbf{锚点细化：} TalkingGaussian：针对直接在动态面部区域生成逐点外观产生扭曲伪影的局限，用\textbf{面部/口腔解耦的双分支变形场}替代单一变形场，核心思想是\textbf{面部运动和口腔运动在空间和时间上具有不同特性}——面部表情变化是渐进且空间连续的，而口腔（嘴唇开合、牙齿暴露）是快速且空间不连续的，两个独立分支分别学习两种截然不同的运动模式、显著提高拟合精度。
\textbf{方法信息：} 代码：\url{https://github.com/fictionarry/talkinggaussian}；GitHub 约 383 stars；Semantic Scholar 约 79 次引用。



% METHOD: 2404.16012
\paragraph{GaussianTalker（三平面空间特征与音频注意力）（2024.04）}

GaussianTalker~\cite{2404_16012} 解决 3DGS 中相邻点缺乏空间信息共享的问题。它用多分辨率 triplane 编码规范 3D 位置，解码每个高斯的尺度、旋转、SH 和不透明度；音频驱动阶段用 spatial-audio attention 融合高斯空间特征与音频特征，预测每帧属性偏移。训练先学规范形状，再学变形模块。项目仓库包含两个相关仓库。相对 TalkingGaussian 的持久结构加分支变形，GaussianTalker 的增量是共享隐式特征空间，使大量高斯能被稳定地音频控制。
\textbf{设计动机与边界：} 相对原始 3DGS 以独立高斯和图像重建误差为主的设定，GaussianTalker 把问题落在以下可实现的改动上：GaussianTalker构建规范的3DGS头部表示，并通过音频驱动变形。核心创新在于将3D高斯属性编码到共享隐式特征空间：对每个规范3D位置，从多分辨率三平面表示中提取特征嵌入，用于计算每个点的缩放、旋转、球谐函数和不透明度，构成规范3D高斯。随后设计空间-音频注意力模块，将3D高斯特征与音频特征融合，预测每个高斯属性的逐帧偏移量。相比之前的拼接或乘法方法，该设计更稳定地操控大量高斯及其复杂参数。训练采用分阶段策略，先学习规范形状再学习变形模块。该方法利用了3DGS的快速渲染能力，同时通过隐式特征编码解决了3DGS表示中相邻点缺乏空间信息共享的问题。其设计目的不是单独增加一个模块，而是用上述表示、监督或渲染约束消除原流程在目标任务中的具体失配。提出将3D高斯属性编码到共享隐式特征表示中，与音频特征融合后通过空间-音频注意力模块预测每帧偏移量；采用多分辨率三平面表示学习规范3D高斯，解决了3DGS缺乏空间信息共享的问题；实现高达120 FPS的渲染速度，在面部保真度、唇形同步准确性和渲染速度上超越先前方法。适用时还应保留论文给出的边界：需要每人单独训练；对极端头部姿态泛化能力有限；语音特征提取依赖预训练ASR模型。
\textbf{方法信息：} 代码：\url{https://github.com/cvlab-kaist/gaussiantalker}；GitHub 关注度较高，约 406 stars。



% METHOD: 2412.20148
\paragraph{DEGSTalk（长发保持的说话脸）（2024.12）}
DEGSTalk~\cite{2412_20148} 专门处理长发说话人脸。它基于 TalkingGaussian 设置，引入 deformable pre-embedding Gaussian field：每个高斯有可学习嵌入，经隐式表情系数动态调整位置、尺度和旋转。triplane hash encoder 与 MLP 解码点级变形。另一个关键模块是 dynamic hair-preserving rendering，增强长发运动和边界稳定。相对 TalkingGaussian 的面部与口腔分支，DEGSTalk 的局部增量是把头发作为独立难点保留下来，不让说话变形破坏长发外观。
\textbf{方法信息：} 代码：\url{https://github.com/cvi-szu/degstalk}。
\textbf{设计动机与边界：} 相对原始 3DGS 以独立高斯和图像重建误差为主的设定，DEGSTalk 把问题落在以下可实现的改动上：DEGSTalk的核心设计：(1) 可变形预嵌入高斯场：为每个高斯基元关联一个可学习嵌入，编码面部区域的细微信息，通过隐式表情系数动态调整高斯基元的位置、缩放和旋转，精确捕捉面部动态和微妙表情变化；(2) 使用三平面哈希编码器和MLP解码器预测每个基元的点级变形；(3) 动态头发保持肖像渲染：专门设计的渲染策略增强长发运动的真实感，解决现有方法在长发个体上的失真问题；(4) 基于TalkingGaussian的设置，通过音频特征和表情特征驱动变形。其设计目的不是单独增加一个模块，而是用上述表示、监督或渲染约束消除原流程在目标任务中的具体失配。提出DEGSTalk，首个专门解决长发保持的3DGS说话人脸合成方法，通过可变形预嵌入高斯场精确捕捉面部动态，引入动态头发保持渲染技术，在复杂面部动态和头发保持方面优于现有方法。适用时还应保留论文给出的边界：长发运动建模仍面临复杂头发物理模拟的挑战；训练需要高质量的多视角视频数据；对极端表情的泛化能力有限。
% METHOD: 2503.14198
\paragraph{RoGSplat（稀疏多视角人体高斯的鲁棒泛化）（2025.03）}

RoGSplat~\cite{2503_14198} 虽被排在本节，但任务更接近稀疏多视角人体高斯重建。它把 SMPL 顶点通过 Snowflake Point Deconvolution 致密化为 3D 先验点，再结合像素级特征和体素级特征预测 image-aligned 3D prior points。随后粗到细做 pixel-wise Gaussian regression：先基于先验点回归粗高斯，再渲染深度辅助细粒度高斯回归。相对 GPS-Gaussian 的双目深度加逐像素高斯，RoGSplat 更强调 SMPL 先验点致密化和跨数据泛化。
\textbf{实现细节：} RoGSplat包含两个核心组件：图像对齐的人体点预测和粗到细像素级高斯回归。首先，利用Snowflake Point Deconvolution（SPD）网络对稀疏SMPL顶点进行致密化，并结合像素级特征（通过U-Net提取并投影聚合）和体素级特征（通过深度精炼网络获取）来预测图像对齐的3D先验点。然后在粗到细框架中，先基于先验点回归粗高斯，再渲染深度图辅助回归细粒度像素级高斯，以捕获高频细节。整个方法以前馈方式运行，无需逐主体优化，实现了跨数据集的泛化能力。 提出RoGSplat，通过将SMPL顶点提升为密集3D先验点，并结合像素级和体素级特征预测图像对齐的3D先验点，再采用粗到细的像素级高斯回归策略，实现从稀疏多视角图像的鲁棒可泛化人体高斯泼溅重建，无需逐主体优化。 对SMPL模型估计质量有一定依赖；在极端稀疏（如单视角）场景下性能仍有下降；主要针对人体场景，未验证在一般场景中的适用性。
\textbf{设计动机与边界：} 相对原始 3DGS 以独立高斯和图像重建误差为主的设定，RoGSplat 把问题落在以下可实现的改动上：RoGSplat包含两个核心组件：图像对齐的人体点预测和粗到细像素级高斯回归。首先，利用Snowflake Point Deconvolution（SPD）网络对稀疏SMPL顶点进行致密化，并结合像素级特征（通过U-Net提取并投影聚合）和体素级特征（通过深度精炼网络获取）来预测图像对齐的3D先验点。然后在粗到细框架中，先基于先验点回归粗高斯，再渲染深度图辅助回归细粒度像素级高斯，以捕获高频细节。整个方法以前馈方式运行，无需逐主体优化，实现了跨数据集的泛化能力。其设计目的不是单独增加一个模块，而是用上述表示、监督或渲染约束消除原流程在目标任务中的具体失配。提出RoGSplat，通过将SMPL顶点提升为密集3D先验点，并结合像素级和体素级特征预测图像对齐的3D先验点，再采用粗到细的像素级高斯回归策略，实现从稀疏多视角图像的鲁棒可泛化人体高斯泼溅重建，无需逐主体优化。适用时还应保留论文给出的边界：对SMPL模型估计质量有一定依赖；在极端稀疏（如单视角）场景下性能仍有下降；主要针对人体场景，未验证在一般场景中的适用性。
\textbf{方法信息：} Adobe 参与了这篇论文；代码：\url{https://github.com/isee-laboratory/rogsplat}。



% METHOD: 2503.17032
\paragraph{TaoAvatar（移动 AR 全身说话化身）（2025.03）}
TaoAvatar~\cite{2503_17032} 把 talking avatar 推到全身和移动 AR。它先构建 SMPLX++ 混合穿衣参数化模型，用 NeuS2 重建衣物和头发几何，并把高斯绑定到网格三角形局部坐标系。教师-学生框架中，复杂姿态相关非刚性变形由 StyleUNet 教师建模，再蒸馏到轻量 MLP；面部表情和身体姿态的 blendshape 补偿蒸馏损失的细节。相对只生成头部的 talking Gaussian，TaoAvatar 的增量是全身穿衣参数化和移动设备实时运行。
\textbf{实现细节：} TaoAvatar（2025.03） 实现了在 Apple Vision Pro 上以 90 FPS 运行的全身体说话人化身。与上述纯 3DGS 方法不同，TaoAvatar 采用教师-学生框架：教师模型使用 StyleUNet 生成高保真非刚性变形，然后通过知识蒸馏将 StyleUNet 的变形"烘焙"为轻量级 MLP 网络（非刚性变形烘焙+双层混合形状补偿），结合扩展的 SMPLX++ 参数化模型（包含衣物和头发的完整几何）和绑定在网格三角面片上的 3D 高斯实现高效移动端渲染。在情绪控制方面，除唇形同步外，说话人头像还需具备根据情绪状态调整面部表情的能力。
\textbf{方法信息：} Alibaba 参与了这篇论文。
\textbf{设计动机与边界：} 相对原始 3DGS 以独立高斯和图像重建误差为主的设定，TaoAvatar 把问题落在以下可实现的改动上：TaoAvatar包含三个核心设计：(1) 混合穿衣参数化表示（SMPLX++）：扩展SMPLX为穿衣模型，使用NeuS2重建穿衣几何，将非身体组件（衣服、头发）的蒙皮权重从身体传播过来，得到完整的个性化参数化模型。然后将3D高斯绑定到网格三角形上，定义局部坐标系属性，使高斯随网格同步运动。(2) 教师-学生框架：预训练StyleUNet网络处理复杂姿态相关的非刚性变形以捕获高频外观细节。然后通过蒸馏技术将非刚性变形「烘焙」到轻量MLP网络中。(3) 混合形状补偿：为弥补蒸馏过程中的细节损失，开发面部表情和身体姿态的混合形状，补偿高频细节。整体模型在保持高质量的同时足够轻量，适合移动AR设备。其设计目的不是单独增加一个模块，而是用上述表示、监督或渲染约束消除原流程在目标任务中的具体失配。提出TaoAvatar，通过混合穿衣参数化表示将高斯绑定到网格、教师-学生框架烘焙非刚性变形和轻量混合形状补偿，实现可在Apple Vision Pro上以90FPS运行的高保真全身说话虚拟形象。适用时还应保留论文给出的边界：需要多视角RGB视频作为训练输入；个性化模板构建需要额外处理步骤；对极端姿态的泛化受限。
% METHOD: 2510.08587
\paragraph{EGSTalker（KAN 与高效空间音频注意力）（2025.10）}
EGSTalker~\cite{2510_08587} 追求说话人脸质量、唇形和速度的平衡。静态阶段用多分辨率哈希 triplane 编码 3D 点，并用 KAN 网络替代传统 MLP 捕获高维非线性映射。音频驱动阶段，Efficient Spatial-Audio Attention 使用代理机制融合音频和空间特征，并用周期位置编码稳定时序。项目实现围绕 efficient Gaussian deformation。相对 GaussianTalker，EGSTalker 的增量是把空间编码和音频融合都做轻量化，适合短训练视频和实时生成。
\textbf{实现细节：} EGSTalker（2025.10） 从三个维度系统性地优化了 3DGS 说话人头像的效率与质量权衡。读者已从上一节了解到 GaussianTalker 的变形范式——维护规范 3D 高斯并通过音频驱动变形场预测逐帧偏移量。EGSTalker 继承了这一变形范式，但精准定位了其中三个效率瓶颈并逐一重新设计：用 Kolmogorov-Arnold 网络替代 MLP 以更少的参数实现更强的非线性映射能力，用轻量级代理注意力替代全交叉注意力以降低计算开销，并引入两阶段训练策略使训练数据需求从 10 分钟缩短到 3--5 分钟。Baseline 与问题。GaussianTalker 等 3DGS 说话人方法使用标准 MLP 将空间特征映射到高斯参数、用完整交叉注意力融合音频与空间特征。这两处设计存在效率与质量的双重瓶颈。在静态高斯构建阶段，标准 MLP 的高维非线性映射存在信息损失——面部特征到高斯参数之间的复杂关系难以被常规 MLP 的分段线性激活函数（如 ReLU）精确捕捉。在动态变形阶段，标准交叉注意力的计算复杂度为 $O(N^{}2 d)$（$N$ 为空间点数，$d$ 为特征维度），当面部分辨率提高时计算量急剧膨胀，阻碍实时推理。
\textbf{方法信息：} 代码：\url{https://github.com/zhutianheng/egstalker}。
\textbf{设计动机与边界：} 相对原始 3DGS 以独立高斯和图像重建误差为主的设定，EGSTalker 把问题落在以下可实现的改动上：EGSTalker采用两阶段训练策略：（1）静态高斯初始化：设计空间结构编码器，包含多分辨率哈希三平面编码器（将3D点投影到三个正交2D网格，通过Hadamard积融合多分辨率特征）和KAN网络（使用多项式近似精确捕获面部特征与3D高斯参数之间的复杂映射关系）。传统MLP在高维非线性映射中存在信息损失和计算低效问题，KAN通过多项式近似解决了这一问题。（2）音频驱动高斯变形：设计高效高斯变形解码器，包含ESAA（Efficient Spatial-Audio Attention）模块（使用代理机制融合音频和空间特征，保持低计算复杂度）和周期位置编码（PPE）。仅需3-5分钟训练视频。其设计目的不是单独增加一个模块，而是用上述表示、监督或渲染约束消除原流程在目标任务中的具体失配。提出EGSTalker框架，通过多分辨率哈希三平面编码器和KAN网络构建紧凑3DGS表示，结合高效空间-音频注意力（ESAA）模块实现音频-空间特征快速融合，在渲染质量、唇形同步和推理速度上达到SOTA平衡。适用时还应保留论文给出的边界：依赖训练视频的质量和多样性；对极端头部运动可能效果下降；未考虑情感和语音内容的细粒度关联。
% METHOD: 2512.05991
\paragraph{EmoDiffTalk（AU 提示高斯扩散）（2025.12）}
EmoDiffTalk~\cite{2512_05991} 用 AU code bridge 连接情感和表情。先从多视角图像重建规范 3DGS 骨架，并把头部分成面部和非面部区域；然后 Speech-to-AU encoder 把语音映射到 action units，AU 作为提示引导高斯扩散预测动态高斯的位置、尺度、旋转、不透明度和颜色；文本到 AU 情感控制器允许用文本编辑情感状态。相对 EmoTaG 的 GRMN 回归，EmoDiffTalk 用扩散过程建模更细的情感面部动画。
\textbf{实现细节：} EmoDiffTalk（2025.12） 引入动作单元（AU）提示的高斯扩散过程实现细粒度面部动画器，与 EmoTalkingGaussian 在效价-唤醒度空间操作不同，EmoDiffTalk 先用三平面预测高斯颜色属性替代球谐系数以获得更好的颜色归因，然后将音频编码为 AU 码驱动高斯位置等属性，配合文本到 AU 的情绪控制器建立从文本情绪描述到 AU 激活向量的直接映射，实现可扩展的动态情绪编辑。
\textbf{方法信息：} Tencent 参与了这篇论文。
\textbf{设计动机与边界：} 相对原始 3DGS 以独立高斯和图像重建误差为主的设定，EmoDiffTalk 把问题落在以下可实现的改动上：EmoDiffTalk框架包含三个核心阶段：（1）规范高斯骨架重建：从多视角图像重建规范3DGS骨架，将3D头部分为面部和非面部区域，建立运动绑定。采用三平面（triplane）高斯颜色预测替代球谐函数，在AU提示扩散过程中实现更好的颜色属性获取。（2）AU提示高斯扩散：Speech-to-AU编码器将语音特征映射为面部动作单元（AU）码，AU码作为提示引导高斯扩散过程预测动态高斯原语属性（位置、缩放、旋转、不透明度、颜色），实现精细面部动画。（3）文本到AU情感控制器：建立从文本情感提示到AU激活向量的直接映射，使文本输入能自由编辑情感状态。关键创新在于通过AU码空间建模情感到表情的映射，而非在隐特征空间或直接扩散空间中处理，使情感编辑更加准确和可解释。其设计目的不是单独增加一个模块，而是用上述表示、监督或渲染约束消除原流程在目标任务中的具体失配。提出首个支持连续多模态情感编辑的3DGS说话人生成框架，通过AU码空间桥接情感到表情的映射，实现精细面部动画和文本驱动的扩展情感编辑，在EmoTalk3D和RenderMe-360上超越现有方法。适用时还应保留论文给出的边界：AU码空间的表达能力上限可能限制极端情感表现；依赖高质量多视角训练数据；文本到AU映射可能存在歧义。
% METHOD: 2601.01847
\paragraph{ESGaussianFace（情感与风格双变形预测）（2026.01）}
ESGaussianFace~\cite{2601_01847} 同时处理 3D 一致性、情感表达和风格化。它在 3DMM 标准脸上初始化高斯，用多分辨率 triplane 解码规范高斯参数；音频 DeepSpeech 特征驱动唇部，情感视频提取的 3DMM 表情系数提供情感条件。情感-音频引导空间注意力融合两者，两个高斯 deformation predictors 分别预测情感和风格控制下的变形。训练分阶段学习唇形同步、情感变化和风格。相对 EmoTaG 的 FLAME 参数运动，ESGaussianFace 更偏特征融合和多风格生成。
\textbf{实现细节：} 框架包含三个核心模块：(1) 基于三平面的3D高斯生成器：在3DMM标准面部模型上初始化3D高斯，通过多分辨率三平面编码位置特征，解码器获取规范高斯参数；(2) 音视频特征提取与融合：提取音频Deep Speech特征驱动唇部运动，从情感视频提取3DMM表情系数作为情感特征，设计情感-音频引导空间注意力融合两者；(3) ESGaussian变形预测：两个变形预测器分别预测情感和风格控制下的高斯参数变形，神经渲染器生成最终视频。多阶段训练：第一阶段学唇部同步，第二阶段学情感变化，第三阶段学风格特征。 首个同时集成3D一致性、情感表达和风格特征的音频驱动人脸动画框架。核心创新包括：(1) 情感-音频引导的空间注意力方法，有效融合情感和音频内容特征；(2) 两个3D高斯变形预测器分别处理情感和风格变形；(3) 多阶段训练策略逐步学习唇部运动→情感变化→风格特征，单模型处理多种情感。 情感-音频解耦可能不完美；对极端姿态和表情的泛化有限；多阶段训练增加整体训练时间。
\textbf{方法信息：} Microsoft 参与了这篇论文。
\textbf{设计动机与边界：} 相对原始 3DGS 以独立高斯和图像重建误差为主的设定，ESGaussianFace 把问题落在以下可实现的改动上：框架包含三个核心模块：(1) 基于三平面的3D高斯生成器：在3DMM标准面部模型上初始化3D高斯，通过多分辨率三平面编码位置特征，解码器获取规范高斯参数；(2) 音视频特征提取与融合：提取音频Deep Speech特征驱动唇部运动，从情感视频提取3DMM表情系数作为情感特征，设计情感-音频引导空间注意力融合两者；(3) ESGaussian变形预测：两个变形预测器分别预测情感和风格控制下的高斯参数变形，神经渲染器生成最终视频。多阶段训练：第一阶段学唇部同步，第二阶段学情感变化，第三阶段学风格特征。其设计目的不是单独增加一个模块，而是用上述表示、监督或渲染约束消除原流程在目标任务中的具体失配。首个同时集成3D一致性、情感表达和风格特征的音频驱动人脸动画框架。核心创新包括：(1) 情感-音频引导的空间注意力方法，有效融合情感和音频内容特征；(2) 两个3D高斯变形预测器分别处理情感和风格变形；(3) 多阶段训练策略逐步学习唇部运动→情感变化→风格特征，单模型处理多种情感。适用时还应保留论文给出的边界：情感-音频解耦可能不完美；对极端姿态和表情的泛化有限；多阶段训练增加整体训练时间。
% METHOD: 2603.21332
\paragraph{EmoTaG（少样本情感感知说话头像）（2026.03）}
EmoTaG~\cite{2603_21332} 把音频驱动重新定义到结构化 FLAME 参数空间，而不是直接预测高斯变形。FLAME-Gaussian 模型把高斯绑定到 FLAME 三角面，口内区域通过唇部标志点选择和约束增长得到。门控残差运动网络融合韵律音频、上面部 AU 和身份条件，基础分支学习语音到运动，残差分支补情感和身份特定变化，门控分支平衡二者。相对 GaussianTalker，EmoTaG 的关键是让情感通过显式几何参数和门控残差进入运动。
\textbf{实现细节：} EmoTaG（2026.03） 在 FLAME 参数空间而非直接在 3D 高斯上预测运动，与 EmoDiffTalk 的 AU（Action Unit，面部动作单元）驱动策略不同，EmoTaG 引入门控残差运动网络（GRMN），通过身份条件编码器融合音频和情绪信号的 AdaIN 调制，专家运动解码器的三个互补分支（残差、门控和语义情绪引导）协同预测 FLAME 表情参数和下颌姿态。
\textbf{实现补充：} 实现上还要核对输入、输出、参数、特征、初始化、训练损失、正则、渲染或推理接口，以及优化中的剪枝、分裂、采样、编码或解码步骤；这些抓手决定该方法到底改了 3DGS 管线的哪一环。
\textbf{技术实现补充：} EmoTaG（2026.03） 在 FLAME 参数空间而非直接在 3D 高斯上预测运动，与 EmoDiffTalk 的 AU（Action Unit，面部动作单元）驱动策略不同，EmoTaG 引入门控残差运动网络（GRMN），通过身份条件编码器融合音频和情绪信号的 AdaIN 调制，专家运动解码器的三个互补分支（残差、门控和语义情绪引导）协同预测 FLAME 表情参数和下颌姿态。
\textbf{方法信息：} Qualcomm 参与了这篇论文。
\textbf{设计动机与边界：} 相对原始 3DGS 以独立高斯和图像重建误差为主的设定，EmoTaG 把问题落在以下可实现的改动上：EmoTaG框架基于预训练-适配(Pretrain-and-Adapt)范式。核心设计包括：(1)FLAME-Gaussian模型：以FLAME参数模型为显式几何先验，预测表情和下颌参数驱动绑定式3D Gaussians。Gaussians通过重心插值绑定到FLAME网格三角形上，口内区域通过唇部标志点选择并约束增长获得。(2)门控残差运动网络(GRMN)：包含身份条件编码器和专家运动解码器。编码器通过AdaIN调制融合韵律感知音频特征和上面部AU线索，解码器包含三个协作分支：基础分支捕获语音音频-运动映射，残差分支建模情感相关和身份特定变化，门控分支自适应平衡两者。语义情感引导从预训练情感识别器蒸馏情感分布和强度分数。其设计目的不是单独增加一个模块，而是用上述表示、监督或渲染约束消除原流程在目标任务中的具体失配。提出EmoTaG，首个少样本情感感知3D说话人合成框架。关键创新包括：(1)将运动预测重新定义为结构化FLAME参数空间而非直接变形Gaussians，引入显式几何先验提升运动稳定性；(2)提出门控残差运动网络(GRMN)，捕获音频中的情感韵律并补充头部姿态和上面部线索。在情感表现力、唇形同步、视觉真实感和运动稳定性上实现SOTA。适用时还应保留论文给出的边界：依赖FLAME模型的面部表达能力上限；对极端情感表情的建模可能受限；少样本适配需要至少几秒的目标身份视频。

\section{手部与局部建模}

手部方向的锚点是 MANO 或 MANO-HD。高斯方法要同时处理手部细长结构、强遮挡、手-物体接触、双手交互和野外视频扰动。相比全身人体，手部更依赖局部几何先验、交互损失和遮挡补全。
% METHOD: 2311.08581
\paragraph{Drivable 3D Gaussian Avatars（人体、头像与可动画角色）（2023.11）}
Drivable 3D Gaussian Avatars~\cite{2311_08581} 关注的问题是：用 3DGS 建模人体、头像或可动画角色 Drivable 3D Gaussian Avatars 是宽召回补全中发现的论文，主要根据标题、摘要和 arXiv 类别判定其与 3DGS 的关系。
相对原始 3DGS，Drivable 3D Gaussian Avatars 的变化不是只替换损失函数，而是把表示、初始化、约束或网络预测路径重新组织到任务流程中：该工作被召回是因为论文文本包含 Gaussian/3D Gaussian 相关表述；当前分类先确定其任务域和与 3DGS 的相关性，后续证据流水线再补代码、企业和引用信息。
\textbf{实现细节：} 具体实现可以按输入、核心模块和渲染/优化路径来读。输入端首先取得任务所需的图像、相机、文本、时间序列、语义或几何线索，并把这些线索转换成可以约束高斯中心、尺度、旋转、不透明度、颜色或附加特征的条件。核心结构中，方法使用的主要机制是：从摘要层面看，Drivable 3D Gaussian Avatars 围绕“用 3DGS 建模人体、头像或可动画角色”展开。若其正文或代码后续可读，应进一步核对它具体改动的是表示、初始化、渲染、优化、损失、动态建模、压缩或应用接口中的哪一环。 训练或推理时，这些模块通常先给出初始化或中间特征，再通过可微 Gaussian rasterization、重建损失、感知/语义约束、几何正则或任务损失回传到高斯参数；如果方法包含前馈网络，则网络输出的是高斯属性、变形场、压缩码、语义特征或控制变量，而不是只输出最终 2D 图像。渲染路径仍以显式高斯投影、深度排序和 alpha 合成为基础，只是在采样、属性预测、局部编辑、动态变形、压缩编码或下游任务接口处加入额外约束。这样写入本章时，它与前后方法的关系是：相对原始 3DGS 和本章已有锚点方法，它可以理解为 3DGS extends\_or\_uses 宽召回分类阶段的关系判断，后续需要按论文方法部分复核。
\textbf{证据与边界：} 本轮召回分类尚未逐项核对实验表格；后续证据补全阶段需要读取论文源码/PDF或项目页确认数据集、指标和对比方法。 该条来自宽召回批次；源码下载受 arXiv e-print 网络限制，部分论文暂时只能依据标题、摘要和类别分类，需要在入可信书前做人工/LLM精读复核。 本书只采用本地 paper pack、数据库分类、源码/引用登记和论文摘要中可核验的信息；没有本地证据支撑的速度、排行榜或开源完整性不在正文中强化。\textbf{方法信息：} 公开链接：\url{https://github.com/cvpr-org/author-kit}，完整性按本地证据保守处理；引用信号：160 citations (5.0/month, domain-relative cutoff 3.0/month, semantic\_scholar)。
\textbf{设计动机与边界：} 相对原始 3DGS 以独立高斯和图像重建误差为主的设定，Drivable 3D Gaussian Avatars 把问题落在以下可实现的改动上：从摘要层面看，Drivable 3D Gaussian Avatars 围绕“用 3DGS 建模人体、头像或可动画角色”展开。若其正文或代码后续可读，应进一步核对它具体改动的是表示、初始化、渲染、优化、损失、动态建模、压缩或应用接口中的哪一环。其设计目的不是单独增加一个模块，而是用上述表示、监督或渲染约束消除原流程在目标任务中的具体失配。该工作被召回是因为论文文本包含 Gaussian/3D Gaussian 相关表述；当前分类先确定其任务域和与 3DGS 的相关性，后续证据流水线再补代码、企业和引用信息。适用时还应保留论文给出的边界：该条来自宽召回批次；源码下载受 arXiv e-print 网络限制，部分论文暂时只能依据标题、摘要和类别分类，需要在入可信书前做人工/LLM精读复核。
% METHOD: 2311.17113
\paragraph{Human Gaussian Splatting（人体、头像与可动画角色）（2023.11）}
Human Gaussian Splatting~\cite{2311_17113} 关注的问题是：用 3DGS 建模人体、头像或可动画角色 Human Gaussian Splatting 是宽召回补全中发现的论文，主要根据标题、摘要和 arXiv 类别判定其与 3DGS 的关系。
相对原始 3DGS，Human Gaussian Splatting 的变化不是只替换损失函数，而是把表示、初始化、约束或网络预测路径重新组织到任务流程中：该工作被召回是因为论文文本包含 Gaussian/3D Gaussian 相关表述；当前分类先确定其任务域和与 3DGS 的相关性，后续证据流水线再补代码、企业和引用信息。
\textbf{实现细节：} 具体实现可以按输入、核心模块和渲染/优化路径来读。输入端首先取得任务所需的图像、相机、文本、时间序列、语义或几何线索，并把这些线索转换成可以约束高斯中心、尺度、旋转、不透明度、颜色或附加特征的条件。核心结构中，方法使用的主要机制是：从摘要层面看，Human Gaussian Splatting 围绕“用 3DGS 建模人体、头像或可动画角色”展开。若其正文或代码后续可读，应进一步核对它具体改动的是表示、初始化、渲染、优化、损失、动态建模、压缩或应用接口中的哪一环。 训练或推理时，这些模块通常先给出初始化或中间特征，再通过可微 Gaussian rasterization、重建损失、感知/语义约束、几何正则或任务损失回传到高斯参数；如果方法包含前馈网络，则网络输出的是高斯属性、变形场、压缩码、语义特征或控制变量，而不是只输出最终 2D 图像。渲染路径仍以显式高斯投影、深度排序和 alpha 合成为基础，只是在采样、属性预测、局部编辑、动态变形、压缩编码或下游任务接口处加入额外约束。这样写入本章时，它与前后方法的关系是：相对原始 3DGS 和本章已有锚点方法，它可以理解为 3DGS extends\_or\_uses 宽召回分类阶段的关系判断，后续需要按论文方法部分复核。
\textbf{证据与边界：} 本轮召回分类尚未逐项核对实验表格；后续证据补全阶段需要读取论文源码/PDF或项目页确认数据集、指标和对比方法。 该条来自宽召回批次；源码下载受 arXiv e-print 网络限制，部分论文暂时只能依据标题、摘要和类别分类，需要在入可信书前做人工/LLM精读复核。 本书只采用本地 paper pack、数据库分类、源码/引用登记和论文摘要中可核验的信息；没有本地证据支撑的速度、排行榜或开源完整性不在正文中强化。\textbf{方法信息：} 公开链接：\url{https://github.com/human-vdm/human-vdm}，完整性按本地证据保守处理；引用信号：135 citations (4.3/month, domain-relative cutoff 3.0/month, semantic\_scholar)。
\textbf{设计动机与边界：} 相对原始 3DGS 以独立高斯和图像重建误差为主的设定，Human Gaussian Splatting 把问题落在以下可实现的改动上：从摘要层面看，Human Gaussian Splatting 围绕“用 3DGS 建模人体、头像或可动画角色”展开。若其正文或代码后续可读，应进一步核对它具体改动的是表示、初始化、渲染、优化、损失、动态建模、压缩或应用接口中的哪一环。其设计目的不是单独增加一个模块，而是用上述表示、监督或渲染约束消除原流程在目标任务中的具体失配。该工作被召回是因为论文文本包含 Gaussian/3D Gaussian 相关表述；当前分类先确定其任务域和与 3DGS 的相关性，后续证据流水线再补代码、企业和引用信息。适用时还应保留论文给出的边界：该条来自宽召回批次；源码下载受 arXiv e-print 网络限制，部分论文暂时只能依据标题、摘要和类别分类，需要在入可信书前做人工/LLM精读复核。
% METHOD: 2311.17857
\paragraph{Gaussian Shell Maps for Efficient 3D Human Generation（人体、头像与可动画角色）（2023.11）}
Gaussian Shell Maps for Efficient 3D Human Generation~\cite{2311_17857} 关注的问题是：用 3DGS 建模人体、头像或可动画角色 Gaussian Shell Maps for Efficient 3D Human Generation 是宽召回补全中发现的论文，主要根据标题、摘要和 arXiv 类别判定其与 3DGS 的关系。
相对原始 3DGS，Gaussian Shell Maps for Efficient 3D Human Generation 的变化不是只替换损失函数，而是把表示、初始化、约束或网络预测路径重新组织到任务流程中：该工作被召回是因为论文文本包含 Gaussian/3D Gaussian 相关表述；当前分类先确定其任务域和与 3DGS 的相关性，后续证据流水线再补代码、企业和引用信息。
\textbf{实现细节：} 具体实现可以按输入、核心模块和渲染/优化路径来读。输入端首先取得任务所需的图像、相机、文本、时间序列、语义或几何线索，并把这些线索转换成可以约束高斯中心、尺度、旋转、不透明度、颜色或附加特征的条件。核心结构中，方法使用的主要机制是：从摘要层面看，Gaussian Shell Maps for Efficient 3D Human Generation 围绕“用 3DGS 建模人体、头像或可动画角色”展开。若其正文或代码后续可读，应进一步核对它具体改动的是表示、初始化、渲染、优化、损失、动态建模、压缩或应用接口中的哪一环。 训练或推理时，这些模块通常先给出初始化或中间特征，再通过可微 Gaussian rasterization、重建损失、感知/语义约束、几何正则或任务损失回传到高斯参数；如果方法包含前馈网络，则网络输出的是高斯属性、变形场、压缩码、语义特征或控制变量，而不是只输出最终 2D 图像。渲染路径仍以显式高斯投影、深度排序和 alpha 合成为基础，只是在采样、属性预测、局部编辑、动态变形、压缩编码或下游任务接口处加入额外约束。这样写入本章时，它与前后方法的关系是：相对原始 3DGS 和本章已有锚点方法，它可以理解为 3DGS extends\_or\_uses 宽召回分类阶段的关系判断，后续需要按论文方法部分复核。
\textbf{证据与边界：} 本轮召回分类尚未逐项核对实验表格；后续证据补全阶段需要读取论文源码/PDF或项目页确认数据集、指标和对比方法。 该条来自宽召回批次；源码下载受 arXiv e-print 网络限制，部分论文暂时只能依据标题、摘要和类别分类，需要在入可信书前做人工/LLM精读复核。 本书只采用本地 paper pack、数据库分类、源码/引用登记和论文摘要中可核验的信息；没有本地证据支撑的速度、排行榜或开源完整性不在正文中强化。\textbf{方法信息：} 代码：\url{https://github.com/facebookresearch/detectron2}；引用信号：79 citations (2.5/month, domain-relative cutoff 3.0/month, semantic\_scholar)；代码关注：computational-imaging/gsm has 227 stars (7.2/month, domain-relative cutoff 32.9/month)。
\textbf{设计动机与边界：} 相对原始 3DGS 以独立高斯和图像重建误差为主的设定，Gaussian Shell Maps for Efficient 3D Human Generation 把问题落在以下可实现的改动上：从摘要层面看，Gaussian Shell Maps for Efficient 3D Human Generation 围绕“用 3DGS 建模人体、头像或可动画角色”展开。若其正文或代码后续可读，应进一步核对它具体改动的是表示、初始化、渲染、优化、损失、动态建模、压缩或应用接口中的哪一环。其设计目的不是单独增加一个模块，而是用上述表示、监督或渲染约束消除原流程在目标任务中的具体失配。该工作被召回是因为论文文本包含 Gaussian/3D Gaussian 相关表述；当前分类先确定其任务域和与 3DGS 的相关性，后续证据流水线再补代码、企业和引用信息。适用时还应保留论文给出的边界：该条来自宽召回批次；源码下载受 arXiv e-print 网络限制，部分论文暂时只能依据标题、摘要和类别分类，需要在入可信书前做人工/LLM精读复核。
% METHOD: 2312.02069
\paragraph{GaussianAvatars（人体、头像与可动画角色）（2023.12）}
GaussianAvatars~\cite{2312_02069} 关注的问题是：用 3DGS 建模人体、头像或可动画角色 GaussianAvatars 是宽召回补全中发现的论文，主要根据标题、摘要和 arXiv 类别判定其与 3DGS 的关系。
相对原始 3DGS，GaussianAvatars 的变化不是只替换损失函数，而是把表示、初始化、约束或网络预测路径重新组织到任务流程中：该工作被召回是因为论文文本包含 Gaussian/3D Gaussian 相关表述；当前分类先确定其任务域和与 3DGS 的相关性，后续证据流水线再补代码、企业和引用信息。
\textbf{实现细节：} 具体实现可以按输入、核心模块和渲染/优化路径来读。输入端首先取得任务所需的图像、相机、文本、时间序列、语义或几何线索，并把这些线索转换成可以约束高斯中心、尺度、旋转、不透明度、颜色或附加特征的条件。核心结构中，方法使用的主要机制是：从摘要层面看，GaussianAvatars 围绕“用 3DGS 建模人体、头像或可动画角色”展开。若其正文或代码后续可读，应进一步核对它具体改动的是表示、初始化、渲染、优化、损失、动态建模、压缩或应用接口中的哪一环。 训练或推理时，这些模块通常先给出初始化或中间特征，再通过可微 Gaussian rasterization、重建损失、感知/语义约束、几何正则或任务损失回传到高斯参数；如果方法包含前馈网络，则网络输出的是高斯属性、变形场、压缩码、语义特征或控制变量，而不是只输出最终 2D 图像。渲染路径仍以显式高斯投影、深度排序和 alpha 合成为基础，只是在采样、属性预测、局部编辑、动态变形、压缩编码或下游任务接口处加入额外约束。这样写入本章时，它与前后方法的关系是：相对原始 3DGS 和本章已有锚点方法，它可以理解为 3DGS extends\_or\_uses 宽召回分类阶段的关系判断，后续需要按论文方法部分复核。
\textbf{证据与边界：} 本轮召回分类尚未逐项核对实验表格；后续证据补全阶段需要读取论文源码/PDF或项目页确认数据集、指标和对比方法。 该条来自宽召回批次；源码下载受 arXiv e-print 网络限制，部分论文暂时只能依据标题、摘要和类别分类，需要在入可信书前做人工/LLM精读复核。 本书只采用本地 paper pack、数据库分类、源码/引用登记和论文摘要中可核验的信息；没有本地证据支撑的速度、排行榜或开源完整性不在正文中强化。\textbf{方法信息：} Official GaussianAvatars repository license states Toyota Motor Europe NV/SA retains IP/proprietary rights 参与；代码：\url{https://github.com/ShenhanQian/GaussianAvatars}；引用信号：313 citations (10.0/month, domain-relative cutoff 3.0/month, semantic\_scholar)；代码关注：shenhanqian/gaussianavatars has 1027 stars (32.9/month, domain-relative cutoff 32.9/month)。
\textbf{原始方法补充：} GaussianAvatars 解决的是原始 3DGS 的高斯没有骨架连接、表情或姿态变化时会漂移的问题。输入是带相机参数的多视图人头视频；每帧先用光度人脸跟踪器拟合 FLAME 网格，再把每个高斯局部绑定到一个 FLAME 三角形，随该三角形从局部空间变换到全局空间后进行 tile rasterization。训练用渲染颜色损失优化局部高斯，同时仍执行增密和剪枝；新生高斯通过绑定继承策略保留三角形关系，并对局部位置和尺度加入正则以抑制新姿态下的伪影。输出是由 FLAME 参数驱动、可实时 splatting 的头像，因此可动画性来自绑定与正则，而不是把每个表情单独重建。
% METHOD: 2312.02902
\paragraph{HeadGaS（人体、头像与可动画角色）（2023.12）}
HeadGaS~\cite{2312_02902} 关注的问题是：用 3DGS 建模人体、头像或可动画角色 HeadGaS 是宽召回补全中发现的论文，主要根据标题、摘要和 arXiv 类别判定其与 3DGS 的关系。
相对原始 3DGS，HeadGaS 的变化不是只替换损失函数，而是把表示、初始化、约束或网络预测路径重新组织到任务流程中：该工作被召回是因为论文文本包含 Gaussian/3D Gaussian 相关表述；当前分类先确定其任务域和与 3DGS 的相关性，后续证据流水线再补代码、企业和引用信息。
\textbf{实现细节：} 具体实现可以按输入、核心模块和渲染/优化路径来读。输入端首先取得任务所需的图像、相机、文本、时间序列、语义或几何线索，并把这些线索转换成可以约束高斯中心、尺度、旋转、不透明度、颜色或附加特征的条件。核心结构中，方法使用的主要机制是：从摘要层面看，HeadGaS 围绕“用 3DGS 建模人体、头像或可动画角色”展开。若其正文或代码后续可读，应进一步核对它具体改动的是表示、初始化、渲染、优化、损失、动态建模、压缩或应用接口中的哪一环。 训练或推理时，这些模块通常先给出初始化或中间特征，再通过可微 Gaussian rasterization、重建损失、感知/语义约束、几何正则或任务损失回传到高斯参数；如果方法包含前馈网络，则网络输出的是高斯属性、变形场、压缩码、语义特征或控制变量，而不是只输出最终 2D 图像。渲染路径仍以显式高斯投影、深度排序和 alpha 合成为基础，只是在采样、属性预测、局部编辑、动态变形、压缩编码或下游任务接口处加入额外约束。这样写入本章时，它与前后方法的关系是：相对原始 3DGS 和本章已有锚点方法，它可以理解为 3DGS extends\_or\_uses 宽召回分类阶段的关系判断，后续需要按论文方法部分复核。
\textbf{证据与边界：} 本轮召回分类尚未逐项核对实验表格；后续证据补全阶段需要读取论文源码/PDF或项目页确认数据集、指标和对比方法。 该条来自宽召回批次；源码下载受 arXiv e-print 网络限制，部分论文暂时只能依据标题、摘要和类别分类，需要在入可信书前做人工/LLM精读复核。 本书只采用本地 paper pack、数据库分类、源码/引用登记和论文摘要中可核验的信息；没有本地证据支撑的速度、排行榜或开源完整性不在正文中强化。\textbf{方法信息：} Huawei 参与；公开链接：\url{https://github.com/cvpr-org/author-kit}，完整性按本地证据保守处理；引用信号：71 citations (2.3/month, domain-relative cutoff 3.0/month, semantic\_scholar)。
\textbf{设计动机与边界：} 相对原始 3DGS 以独立高斯和图像重建误差为主的设定，HeadGaS 把问题落在以下可实现的改动上：从摘要层面看，HeadGaS 围绕“用 3DGS 建模人体、头像或可动画角色”展开。若其正文或代码后续可读，应进一步核对它具体改动的是表示、初始化、渲染、优化、损失、动态建模、压缩或应用接口中的哪一环。其设计目的不是单独增加一个模块，而是用上述表示、监督或渲染约束消除原流程在目标任务中的具体失配。该工作被召回是因为论文文本包含 Gaussian/3D Gaussian 相关表述；当前分类先确定其任务域和与 3DGS 的相关性，后续证据流水线再补代码、企业和引用信息。适用时还应保留论文给出的边界：该条来自宽召回批次；源码下载受 arXiv e-print 网络限制，部分论文暂时只能依据标题、摘要和类别分类，需要在入可信书前做人工/LLM精读复核。
% METHOD: 2312.04558
\paragraph{MonoGaussianAvatar（人体、头像与可动画角色）（2023.12）}
MonoGaussianAvatar~\cite{2312_04558} 关注的问题是：用 3DGS 建模人体、头像或可动画角色 MonoGaussianAvatar 是宽召回补全中发现的论文，主要根据标题、摘要和 arXiv 类别判定其与 3DGS 的关系。
相对原始 3DGS，MonoGaussianAvatar 的变化不是只替换损失函数，而是把表示、初始化、约束或网络预测路径重新组织到任务流程中：该工作被召回是因为论文文本包含 Gaussian/3D Gaussian 相关表述；当前分类先确定其任务域和与 3DGS 的相关性，后续证据流水线再补代码、企业和引用信息。
\textbf{实现细节：} 具体实现可以按输入、核心模块和渲染/优化路径来读。输入端首先取得任务所需的图像、相机、文本、时间序列、语义或几何线索，并把这些线索转换成可以约束高斯中心、尺度、旋转、不透明度、颜色或附加特征的条件。核心结构中，方法使用的主要机制是：从摘要层面看，MonoGaussianAvatar 围绕“用 3DGS 建模人体、头像或可动画角色”展开。若其正文或代码后续可读，应进一步核对它具体改动的是表示、初始化、渲染、优化、损失、动态建模、压缩或应用接口中的哪一环。 训练或推理时，这些模块通常先给出初始化或中间特征，再通过可微 Gaussian rasterization、重建损失、感知/语义约束、几何正则或任务损失回传到高斯参数；如果方法包含前馈网络，则网络输出的是高斯属性、变形场、压缩码、语义特征或控制变量，而不是只输出最终 2D 图像。渲染路径仍以显式高斯投影、深度排序和 alpha 合成为基础，只是在采样、属性预测、局部编辑、动态变形、压缩编码或下游任务接口处加入额外约束。这样写入本章时，它与前后方法的关系是：相对原始 3DGS 和本章已有锚点方法，它可以理解为 3DGS extends\_or\_uses 宽召回分类阶段的关系判断，后续需要按论文方法部分复核。
\textbf{证据与边界：} 本轮召回分类尚未逐项核对实验表格；后续证据补全阶段需要读取论文源码/PDF或项目页确认数据集、指标和对比方法。 该条来自宽召回批次；源码下载受 arXiv e-print 网络限制，部分论文暂时只能依据标题、摘要和类别分类，需要在入可信书前做人工/LLM精读复核。 本书只采用本地 paper pack、数据库分类、源码/引用登记和论文摘要中可核验的信息；没有本地证据支撑的速度、排行榜或开源完整性不在正文中强化。\textbf{方法信息：} 代码：\url{https://github.com/zju3dv/easymocap}；引用信号：126 citations (4.1/month, domain-relative cutoff 3.0/month, semantic\_scholar)；代码关注：yufan1012/monogaussianavatar has 145 stars (4.7/month, domain-relative cutoff 32.9/month)。
\textbf{原始方法补充：} MonoGaussianAvatar 面向单目视频，避免原始 3DGS 需要多视图静态初始化且无法以 FLAME 姿态稳定驱动。它从人物的表情和姿态视频开始，在初始化空间维护可插入和删除的高斯均值；参数预测网络以均值为输入，生成不透明度、旋转、尺度和颜色。可学习 offset 将均值映到 canonical space，FLAME 的表情与姿态再决定其变形空间位置；渲染同时使用预测属性和变形场，并以真实帧图像监督。插删策略会随训练收敛把点分配到头发和细节区，设计动机是兼顾单目可获得性、细节和大幅运动时的结构稳定性，而不是依赖隐式场的平滑几何。
% METHOD: 2312.05941
\paragraph{ASH（人体、头像与可动画角色）（2023.12）}
ASH~\cite{2312_05941} 关注的问题是：用 3DGS 建模人体、头像或可动画角色 ASH 是宽召回补全中发现的论文，主要根据标题、摘要和 arXiv 类别判定其与 3DGS 的关系。
相对原始 3DGS，ASH 的变化不是只替换损失函数，而是把表示、初始化、约束或网络预测路径重新组织到任务流程中：该工作被召回是因为论文文本包含 Gaussian/3D Gaussian 相关表述；当前分类先确定其任务域和与 3DGS 的相关性，后续证据流水线再补代码、企业和引用信息。
\textbf{实现细节：} 具体实现可以按输入、核心模块和渲染/优化路径来读。输入端首先取得任务所需的图像、相机、文本、时间序列、语义或几何线索，并把这些线索转换成可以约束高斯中心、尺度、旋转、不透明度、颜色或附加特征的条件。核心结构中，方法使用的主要机制是：从摘要层面看，ASH 围绕“用 3DGS 建模人体、头像或可动画角色”展开。若其正文或代码后续可读，应进一步核对它具体改动的是表示、初始化、渲染、优化、损失、动态建模、压缩或应用接口中的哪一环。 训练或推理时，这些模块通常先给出初始化或中间特征，再通过可微 Gaussian rasterization、重建损失、感知/语义约束、几何正则或任务损失回传到高斯参数；如果方法包含前馈网络，则网络输出的是高斯属性、变形场、压缩码、语义特征或控制变量，而不是只输出最终 2D 图像。渲染路径仍以显式高斯投影、深度排序和 alpha 合成为基础，只是在采样、属性预测、局部编辑、动态变形、压缩编码或下游任务接口处加入额外约束。这样写入本章时，它与前后方法的关系是：相对原始 3DGS 和本章已有锚点方法，它可以理解为 3DGS extends\_or\_uses 宽召回分类阶段的关系判断，后续需要按论文方法部分复核。
\textbf{证据与边界：} 本轮召回分类尚未逐项核对实验表格；后续证据补全阶段需要读取论文源码/PDF或项目页确认数据集、指标和对比方法。 该条来自宽召回批次；源码下载受 arXiv e-print 网络限制，部分论文暂时只能依据标题、摘要和类别分类，需要在入可信书前做人工/LLM精读复核。 本书只采用本地 paper pack、数据库分类、源码/引用登记和论文摘要中可核验的信息；没有本地证据支撑的速度、排行榜或开源完整性不在正文中强化。\textbf{方法信息：} 代码：\url{https://github.com/kv2000/ash}；引用信号：126 citations (4.1/month, domain-relative cutoff 3.0/month, semantic\_scholar)；代码关注：kv2000/ash has 137 stars (4.4/month, domain-relative cutoff 32.9/month)。
\textbf{设计动机与边界：} 相对原始 3DGS 以独立高斯和图像重建误差为主的设定，ASH 把问题落在以下可实现的改动上：从摘要层面看，ASH 围绕“用 3DGS 建模人体、头像或可动画角色”展开。若其正文或代码后续可读，应进一步核对它具体改动的是表示、初始化、渲染、优化、损失、动态建模、压缩或应用接口中的哪一环。其设计目的不是单独增加一个模块，而是用上述表示、监督或渲染约束消除原流程在目标任务中的具体失配。该工作被召回是因为论文文本包含 Gaussian/3D Gaussian 相关表述；当前分类先确定其任务域和与 3DGS 的相关性，后续证据流水线再补代码、企业和引用信息。适用时还应保留论文给出的边界：该条来自宽召回批次；源码下载受 arXiv e-print 网络限制，部分论文暂时只能依据标题、摘要和类别分类，需要在入可信书前做人工/LLM精读复核。
% METHOD: 2312.15059
\paragraph{Deformable 3D Gaussian Splatting for Animatable Human Avatars（人体、头像与可动画角色）（2023.12）}
Deformable 3D Gaussian Splatting for Animatable Human Avatars~\cite{2312_15059} 关注的问题是：用 3DGS 建模人体、头像或可动画角色 Deformable 3D Gaussian Splatting for Animatable Human Avatars 是宽召回补全中发现的论文，主要根据标题、摘要和 arXiv 类别判定其与 3DGS 的关系。
相对原始 3DGS，Deformable 3D Gaussian Splatting for Animatable Human Avatars 的变化不是只替换损失函数，而是把表示、初始化、约束或网络预测路径重新组织到任务流程中：该工作被召回是因为论文文本包含 Gaussian/3D Gaussian 相关表述；当前分类先确定其任务域和与 3DGS 的相关性，后续证据流水线再补代码、企业和引用信息。
\textbf{实现细节：} 具体实现可以按输入、核心模块和渲染/优化路径来读。输入端首先取得任务所需的图像、相机、文本、时间序列、语义或几何线索，并把这些线索转换成可以约束高斯中心、尺度、旋转、不透明度、颜色或附加特征的条件。核心结构中，方法使用的主要机制是：从摘要层面看，Deformable 3D Gaussian Splatting for Animatable Human Avatars 围绕“用 3DGS 建模人体、头像或可动画角色”展开。若其正文或代码后续可读，应进一步核对它具体改动的是表示、初始化、渲染、优化、损失、动态建模、压缩或应用接口中的哪一环。 训练或推理时，这些模块通常先给出初始化或中间特征，再通过可微 Gaussian rasterization、重建损失、感知/语义约束、几何正则或任务损失回传到高斯参数；如果方法包含前馈网络，则网络输出的是高斯属性、变形场、压缩码、语义特征或控制变量，而不是只输出最终 2D 图像。渲染路径仍以显式高斯投影、深度排序和 alpha 合成为基础，只是在采样、属性预测、局部编辑、动态变形、压缩编码或下游任务接口处加入额外约束。这样写入本章时，它与前后方法的关系是：相对原始 3DGS 和本章已有锚点方法，它可以理解为 3DGS extends\_or\_uses 宽召回分类阶段的关系判断，后续需要按论文方法部分复核。
\textbf{证据与边界：} 本轮召回分类尚未逐项核对实验表格；后续证据补全阶段需要读取论文源码/PDF或项目页确认数据集、指标和对比方法。 该条来自宽召回批次；源码下载受 arXiv e-print 网络限制，部分论文暂时只能依据标题、摘要和类别分类，需要在入可信书前做人工/LLM精读复核。 本书只采用本地 paper pack、数据库分类、源码/引用登记和论文摘要中可核验的信息；没有本地证据支撑的速度、排行榜或开源完整性不在正文中强化。\textbf{方法信息：} Huawei 参与；代码：\url{https://github.com/zju3dv/easymocap}；引用信号：58 citations (1.9/month, domain-relative cutoff 3.0/month, semantic\_scholar)。
\textbf{设计动机与边界：} 相对原始 3DGS 以独立高斯和图像重建误差为主的设定，Deformable 3D Gaussian Splatting for Animatable Human Avatars 把问题落在以下可实现的改动上：从摘要层面看，Deformable 3D Gaussian Splatting for Animatable Human Avatars 围绕“用 3DGS 建模人体、头像或可动画角色”展开。若其正文或代码后续可读，应进一步核对它具体改动的是表示、初始化、渲染、优化、损失、动态建模、压缩或应用接口中的哪一环。其设计目的不是单独增加一个模块，而是用上述表示、监督或渲染约束消除原流程在目标任务中的具体失配。该工作被召回是因为论文文本包含 Gaussian/3D Gaussian 相关表述；当前分类先确定其任务域和与 3DGS 的相关性，后续证据流水线再补代码、企业和引用信息。适用时还应保留论文给出的边界：该条来自宽召回批次；源码下载受 arXiv e-print 网络限制，部分论文暂时只能依据标题、摘要和类别分类，需要在入可信书前做人工/LLM精读复核。
% METHOD: 2312.02155
\paragraph{GPS-Gaussian（逐像素泛化人体高斯锚点）（2023）}
GPS-Gaussian~\cite{2312_02155} 是人体稀疏视角新视角合成中的前馈锚点。它不再为每个新人物从头优化一套 3DGS，而是把源视图上的每个前景像素都看成一个潜在高斯，直接预测“高斯参数图”。给定多台稀疏相机，系统为目标视角选择两个相邻源视图，用共享图像编码器提取密集特征；随后用 RAFT-Stereo 式迭代双目模块估计视差/深度。深度和相机参数决定像素反投影到 3D 的位置，源图 RGB 给出颜色，其余属性如 opacity、scale、rotation 由网络在 2D 图像平面逐像素回归。

训练分两步更容易理解。代码仓库中 stage1 先训练深度/视差预测分支，使相邻源视图能稳定建立几何对应；stage2 再打开 Gaussian rendering，用预测出的高斯参数经 \texttt{pts2render} 聚合到目标视角，并用 L1 与 SSIM 等图像损失端到端优化。这样得到的是一组从两张源图反投影出的 3D 高斯，而不是某个主体专属的长期 avatar。推理时，输入新人物的稀疏视图即可前馈产生高斯并渲染目标视角，不需要 per-subject fine-tuning。相对 HUGS/3DGS-Avatar 的“单人视频优化”，GPS-Gaussian 的锚点意义是把人体 3DGS 改成了像素对齐的泛化表示。

论文报告支持 2K 分辨率实时人体渲染，并在 ZJU-MoCap、H36M、DNA-Rendering 等数据集上优于多种隐式人体 NVS 方法。代码发布在 \url{https://github.com/aipixel/gps-gaussian}，GitHub 约 622 stars；Semantic Scholar 约 212 次引用。
\textbf{实现细节：} GPS-Gaussian（2023.12） 是该方向的先驱工作。Baseline 与问题：传统基于优化的 3DGS 方法需要对每个新身份从头优化高斯参数，训练耗时数分钟到数十分钟，完全无法满足实时应用需求。改了什么：GPS-Gaussian 提出了无需优化的即时新视角合成框架：给定稀疏源视图，直接通过前馈网络预测 3D 高斯属性，无需任何逐场景优化。实现原理：核心流程分为三步。(1) 视角选择与深度估计：从 $N$ 个输入图像中选择与目标视角最邻近的两个源视图（左视图 $v_l$ 和右视图 $v_r$），通过共享图像编码器提取多尺度特征图，然后使用受 RAFT-Stereo 启发的迭代式双目深度估计器预测视差图（等价于深度图）。选择迭代式估计的原因是它避免了构建高计算成本的代价体积（cost volume），更适合高分辨率输入。(2) 高斯参数预测：对于源视图中的每个前景像素，其深度值和 RGB 值分别确定对应 3D 高斯的位置和颜色，其余参数（缩放、旋转、不透明度）由像素级回归网络预测。这些 2D 参数图通过相机参数反投影到 3D 空间，形成 3D 高斯表示。(3) 跨视图聚合与可微分渲染：来自两个源视图的 3D 高斯在 3D 空间中聚合，通过可微分光栅化渲染到目标视角，整个流程端到端可训练。设计动机：将 3D 高斯的属性预测分解为"从深度+RGB确定位置和颜色"和"从回归网络预测其余参数"两个步骤，充分利用了双目立体视觉提供的几何约束，同时让回归网络仅需学习高斯的形状和透明度属性，大幅降低了学习难度。公式与参数：通过在大规模人体扫描数据上训练的回归模块，支持 2K 分辨率的实时渲染。
\textbf{方法信息：} 代码：\url{https://github.com/aipixel/gps-gaussian}；GitHub 约 622 stars；Semantic Scholar 约 212 次引用。



% METHOD: 2402.06149
\paragraph{HeadStudio（人体、头像与可动画角色）（2024.02）}
HeadStudio~\cite{2402_06149} 关注的问题是：用 3DGS 建模人体、头像或可动画角色 HeadStudio 是宽召回补全中发现的论文，主要根据标题、摘要和 arXiv 类别判定其与 3DGS 的关系。
相对原始 3DGS，HeadStudio 的变化不是只替换损失函数，而是把表示、初始化、约束或网络预测路径重新组织到任务流程中：该工作被召回是因为论文文本包含 Gaussian/3D Gaussian 相关表述；当前分类先确定其任务域和与 3DGS 的相关性，后续证据流水线再补代码、企业和引用信息。
\textbf{实现细节：} 具体实现可以按输入、核心模块和渲染/优化路径来读。输入端首先取得任务所需的图像、相机、文本、时间序列、语义或几何线索，并把这些线索转换成可以约束高斯中心、尺度、旋转、不透明度、颜色或附加特征的条件。核心结构中，方法使用的主要机制是：从摘要层面看，HeadStudio 围绕“用 3DGS 建模人体、头像或可动画角色”展开。若其正文或代码后续可读，应进一步核对它具体改动的是表示、初始化、渲染、优化、损失、动态建模、压缩或应用接口中的哪一环。 训练或推理时，这些模块通常先给出初始化或中间特征，再通过可微 Gaussian rasterization、重建损失、感知/语义约束、几何正则或任务损失回传到高斯参数；如果方法包含前馈网络，则网络输出的是高斯属性、变形场、压缩码、语义特征或控制变量，而不是只输出最终 2D 图像。渲染路径仍以显式高斯投影、深度排序和 alpha 合成为基础，只是在采样、属性预测、局部编辑、动态变形、压缩编码或下游任务接口处加入额外约束。这样写入本章时，它与前后方法的关系是：相对原始 3DGS 和本章已有锚点方法，它可以理解为 3DGS extends\_or\_uses 宽召回分类阶段的关系判断，后续需要按论文方法部分复核。
\textbf{证据与边界：} 本轮召回分类尚未逐项核对实验表格；后续证据补全阶段需要读取论文源码/PDF或项目页确认数据集、指标和对比方法。 该条来自宽召回批次；源码下载受 arXiv e-print 网络限制，部分论文暂时只能依据标题、摘要和类别分类，需要在入可信书前做人工/LLM精读复核。 本书只采用本地 paper pack、数据库分类、源码/引用登记和论文摘要中可核验的信息；没有本地证据支撑的速度、排行榜或开源完整性不在正文中强化。\textbf{方法信息：} 代码：\url{https://github.com/zhenglinzhou/headstudio}；代码关注：zhenglinzhou/headstudio has 215 stars (7.4/month, domain-relative cutoff 32.9/month)。
\textbf{设计动机与边界：} 相对原始 3DGS 以独立高斯和图像重建误差为主的设定，HeadStudio 把问题落在以下可实现的改动上：从摘要层面看，HeadStudio 围绕“用 3DGS 建模人体、头像或可动画角色”展开。若其正文或代码后续可读，应进一步核对它具体改动的是表示、初始化、渲染、优化、损失、动态建模、压缩或应用接口中的哪一环。其设计目的不是单独增加一个模块，而是用上述表示、监督或渲染约束消除原流程在目标任务中的具体失配。该工作被召回是因为论文文本包含 Gaussian/3D Gaussian 相关表述；当前分类先确定其任务域和与 3DGS 的相关性，后续证据流水线再补代码、企业和引用信息。适用时还应保留论文给出的边界：该条来自宽召回批次；源码下载受 arXiv e-print 网络限制，部分论文暂时只能依据标题、摘要和类别分类，需要在入可信书前做人工/LLM精读复核。
% METHOD: 2403.05087
\paragraph{SplattingAvatar（人体、头像与可动画角色）（2024.03）}
SplattingAvatar~\cite{2403_05087} 关注的问题是：用 3DGS 建模人体、头像或可动画角色 SplattingAvatar 是宽召回补全中发现的论文，主要根据标题、摘要和 arXiv 类别判定其与 3DGS 的关系。
相对原始 3DGS，SplattingAvatar 的变化不是只替换损失函数，而是把表示、初始化、约束或网络预测路径重新组织到任务流程中：该工作被召回是因为论文文本包含 Gaussian/3D Gaussian 相关表述；当前分类先确定其任务域和与 3DGS 的相关性，后续证据流水线再补代码、企业和引用信息。
\textbf{实现细节：} 具体实现可以按输入、核心模块和渲染/优化路径来读。输入端首先取得任务所需的图像、相机、文本、时间序列、语义或几何线索，并把这些线索转换成可以约束高斯中心、尺度、旋转、不透明度、颜色或附加特征的条件。核心结构中，方法使用的主要机制是：从摘要层面看，SplattingAvatar 围绕“用 3DGS 建模人体、头像或可动画角色”展开。若其正文或代码后续可读，应进一步核对它具体改动的是表示、初始化、渲染、优化、损失、动态建模、压缩或应用接口中的哪一环。 训练或推理时，这些模块通常先给出初始化或中间特征，再通过可微 Gaussian rasterization、重建损失、感知/语义约束、几何正则或任务损失回传到高斯参数；如果方法包含前馈网络，则网络输出的是高斯属性、变形场、压缩码、语义特征或控制变量，而不是只输出最终 2D 图像。渲染路径仍以显式高斯投影、深度排序和 alpha 合成为基础，只是在采样、属性预测、局部编辑、动态变形、压缩编码或下游任务接口处加入额外约束。这样写入本章时，它与前后方法的关系是：相对原始 3DGS 和本章已有锚点方法，它可以理解为 3DGS extends\_or\_uses 宽召回分类阶段的关系判断，后续需要按论文方法部分复核。
\textbf{证据与边界：} 本轮召回分类尚未逐项核对实验表格；后续证据补全阶段需要读取论文源码/PDF或项目页确认数据集、指标和对比方法。 该条来自宽召回批次；源码下载受 arXiv e-print 网络限制，部分论文暂时只能依据标题、摘要和类别分类，需要在入可信书前做人工/LLM精读复核。 本书只采用本地 paper pack、数据库分类、源码/引用登记和论文摘要中可核验的信息；没有本地证据支撑的速度、排行榜或开源完整性不在正文中强化。\textbf{方法信息：} 代码：\url{https://github.com/pybind/pybind11}；引用信号：230 citations (8.2/month, domain-relative cutoff 3.0/month, semantic\_scholar)；代码关注：initialneil/splattingavatar has 561 stars (20.0/month, domain-relative cutoff 32.9/month)。
\textbf{设计动机与边界：} 相对原始 3DGS 以独立高斯和图像重建误差为主的设定，SplattingAvatar 把问题落在以下可实现的改动上：从摘要层面看，SplattingAvatar 围绕“用 3DGS 建模人体、头像或可动画角色”展开。若其正文或代码后续可读，应进一步核对它具体改动的是表示、初始化、渲染、优化、损失、动态建模、压缩或应用接口中的哪一环。其设计目的不是单独增加一个模块，而是用上述表示、监督或渲染约束消除原流程在目标任务中的具体失配。该工作被召回是因为论文文本包含 Gaussian/3D Gaussian 相关表述；当前分类先确定其任务域和与 3DGS 的相关性，后续证据流水线再补代码、企业和引用信息。适用时还应保留论文给出的边界：该条来自宽召回批次；源码下载受 arXiv e-print 网络限制，部分论文暂时只能依据标题、摘要和类别分类，需要在入可信书前做人工/LLM精读复核。
% METHOD: 2405.12663
\paragraph{LAGA（分层高斯穿衣头像生成）（2024）}
LAGA~\cite{2405_12663} 从文本生成可分解穿衣 avatar。它把人体和服装表示为多层高斯点集，按从内到外顺序生成：先人体，再逐层服装。Dual-SDS loss 同时约束服装自身质量和服装与下层的遮挡连贯性。换装时，形状正则让服装适配目标体型，深度正则处理层间遮挡，外观正则保持纹理连续。相对 PGC 的真实捕获服装，LAGA 更像生成和编辑框架，利用高斯显式点云直接移动服装层。
\textbf{实现细节：} LAGA（2024.05） 实现了从文本描述创建可分解的高保真化身，每层对应特定的服装或人体本身。生成阶段采用由粗到精策略：先用稀疏高斯（5000 点）配合密度引导优化服装粗略形状，再通过带扰动的复制上采样和 SDS 损失细化高频细节；同时提出 dual-SDS 损失，同时优化单层局部图像和多层合成全局图像，确保各层服装之间的视觉连贯性。在服装转移阶段，引入三个正则化损失引导高斯运动：人体拟合损失使服装贴合目标体型轮廓，相似性损失保持转移前后服装的整体形状，可见性损失确保服装不被内层遮挡。该方法支持不同体型之间的服装自由转移，突破了此前仅限相似体型间转移的限制。
\textbf{实现补充：} 实现上还要核对输入、输出、参数、特征、初始化、训练损失、正则、渲染或推理接口，以及优化中的剪枝、分裂、采样、编码或解码步骤；这些抓手决定该方法到底改了 3DGS 管线的哪一环。
\textbf{方法信息：} 代码：\url{https://github.com/sjtu-deepvisionlab/laga}。
\textbf{设计动机与边界：} 相对原始 3DGS 以独立高斯和图像重建误差为主的设定，LAGA 把问题落在以下可实现的改动上：LAGA的核心设计包括三个部分：(1) 分层头像生成：将头像建模为分层结构的高斯点集，N-1层表示服装，1层表示人体。按照从内到外的顺序逐层生成：先生成人体，再依次生成各件服装。每层的高斯点独立优化，通过从粗到细策略确保每件服装的多样性和质量；(2) 双SDS损失(Dual-SDS Loss)：为维持生成的服装与人体（以及其他服装）之间的一致性和遮挡关系，提出双SDS损失，分别约束服装自身的质量和服装与下层之间的视觉连贯性。这确保了叠穿效果的自然性；(3) 服装转移：提出三个正则化损失引导高斯点的移动，使服装可以从一个头像转移到另一个头像：(a) 形状正则化确保服装高斯点与目标体型匹配，(b) 深度正则化处理层间遮挡关系，(c) 外观正则化保持服装纹理的连续性。利用GS的显式点云表示特性，可以直接控制高斯点的位置和属性实现换装。其设计目的不是单独增加一个模块，而是用上述表示、监督或渲染约束消除原流程在目标任务中的具体失配。首次实现基于3DGS的可分解穿衣头像生成框架，将人体和服装分层表示支持自由搭配和换装，提出双SDS损失和三个正则化损失确保生成质量和服装转移的合理性，在3D穿衣人生成质量上超越现有方法适用时还应保留论文给出的边界：分层生成增加了总生成时间；复杂服装（如多层叠穿）的生成质量有待提升；服装转移在体型差异过大时效果下降
% METHOD: 2407.12777
\paragraph{GHG（UV 空间前馈人体高斯）（2024）}
Generalizable Human Gaussians~\cite{2407_12777} 把高斯参数预测转成 SMPL-X UV 空间上的 2D 回归。图像编码器从二到三张稀疏视角提取特征，2D 卷积网络在人体模板 UV 图上预测每个点的高斯位置、颜色、不透明度、尺度和旋转。multi-scaffold 机制表示相对模板表面的多层偏移细节。相对 GPS-Gaussian 的像素反投影，GHG 的增量是利用人体模板 UV 的规则性，减少稀疏视角下 3D 回归歧义。
\textbf{实现细节：} Generalizable Human Gaussians (GHG)（2024.07） 与 GPS-Gaussian 的核心区别在于高斯参数的回归空间：GPS-Gaussian 在像素级 2D 图像空间预测高斯属性，而 GHG 将回归过程转移到人体模板的 2D UV 空间。Baseline 与问题：GPS-Gaussian 等前馈方法在像素空间直接预测高斯参数，缺乏对人体几何结构（关节、自遮挡、复杂表面如头发）的显式建模，导致稀疏视角（如仅 3 视角）下几何和纹理重建严重失真。另一方面，NeRF-based 泛化方法虽然利用了人体先验，但逐像素体渲染的计算瓶颈使其无法实时推理。改了什么：GHG 将 3D 高斯参数的学习重新表述为人体模板 2D UV 空间上的回归过程，利用 SMPL 模板提供的强几何先验和 2D CNN 的邻域信息聚合能力，并引入多脚手架（multi-scaffold）表示弥补模板与真实人体之间的几何差距。实现原理：核心流程分为三步。(1) UV 空间高斯参数化：将 SMPL 模板网格展开为 2D UV 图，每个 UV 像素对应模板上的一个 3D 点，该点的位置即为高斯中心位置的锚定。通过将稀疏输入视图的特征反投影到 UV 空间，获得与人体几何对齐的特征图。
\textbf{方法信息：} Meta 参与了这篇论文。
\textbf{设计动机与边界：} 相对原始 3DGS 以独立高斯和图像重建误差为主的设定，Generalizable Human Gaussians for Sparse View Synthesis 把问题落在以下可实现的改动上：GHG的核心创新在于将3D高斯参数（位置、颜色、不透明度、缩放、旋转）的学习重定义为人体模板（如SMPL-X）的2D UV空间上的回归过程。具体而言，首先通过图像编码器提取稀疏输入视角的特征，然后利用2D卷积网络在UV空间上预测每个UV点对应的3D高斯参数。这种方法的优势是：(1) 利用人体模板提供的强几何先验，解决了稀疏视角下3D回归的歧义性；(2) 2D卷积比3D网络更高效且更容易训练。此外，提出multi-scaffold机制来表示高斯相对于模板表面的偏移细节，通过多尺度脚手架捕捉不同层级的几何细节。整个方法以feed-forward方式运行，给定新的稀疏视角即可直接预测人体3D高斯表示并渲染新视角，无需逐场景优化。其设计目的不是单独增加一个模块，而是用上述表示、监督或渲染约束消除原流程在目标任务中的具体失配。提出GHG方法，首次将3D高斯参数学习转化为人体模板UV空间的2D回归问题，结合人体几何先验实现前馈式稀疏视角人体渲染。多脚手架设计有效表示偏移细节，在数据集内泛化和跨数据集泛化设置上均超越现有方法。适用时还应保留论文给出的边界：依赖人体模板的准确性，对于极端姿势或宽松衣物可能表现不佳；仅针对人体渲染，不适用于一般场景。
% METHOD: 2410.08840
\paragraph{Interaction-aware Hand Avatars（单图交互手部头像）（2024.10）}
Learning Interaction-aware 3D Gaussian Splatting for One-shot Hand Avatars~\cite{2410_08840} 从单张图像创建可动画手部 avatar。它把手部表示解耦为可优化 identity map 和学习到的 latent geometry features 或 neural texture map：前者适应分布外手部，后者提供训练先验。interaction-aware attention module 检测手内和双手交互区域，自适应高斯精炼模块在这些区域增加细节。相对 WildGHand 的视频鲁棒优化，它更强调 one-shot 输入下如何借先验和交互注意力恢复手部。
\textbf{实现细节：} 方法包含三个核心设计：（1）解耦3D手部表示：将手部3D表示解耦为优化身份图（可高效one-shot拟合分布外手部）和学习潜几何特征/神经纹理图（通过网络训练捕获可靠先验）。利用手部网格的显式几何嵌入结合神经纹理图的隐式表示。（2）交互感知注意力模块：检测并处理手部内和手部间的交互区域，增强这些区域的渲染质量。（3）自适应高斯精炼模块：对交互区域的高斯进行针对性精炼。整个框架基于3DGS实现高效渲染。 提出首个基于3DGS的交互手部Avatar框架，通过解耦表示（优化身份图+学习潜特征）和交互感知模块，显著提升单图输入下的手部Avatar图像质量。 依赖MANO等参数化手部模型；对于极端遮挡和极端姿态可能受限；训练需要大规模手部数据集
\textbf{方法信息：} 代码：\url{https://github.com/xuanhuang0/guassianhand}。
\textbf{设计动机与边界：} 相对原始 3DGS 以独立高斯和图像重建误差为主的设定，Learning Interaction-aware 3D Gaussian Splatting for One-shot Hand Avatars 把问题落在以下可实现的改动上：方法包含三个核心设计：（1）解耦3D手部表示：将手部3D表示解耦为优化身份图（可高效one-shot拟合分布外手部）和学习潜几何特征/神经纹理图（通过网络训练捕获可靠先验）。利用手部网格的显式几何嵌入结合神经纹理图的隐式表示。（2）交互感知注意力模块：检测并处理手部内和手部间的交互区域，增强这些区域的渲染质量。（3）自适应高斯精炼模块：对交互区域的高斯进行针对性精炼。整个框架基于3DGS实现高效渲染。其设计目的不是单独增加一个模块，而是用上述表示、监督或渲染约束消除原流程在目标任务中的具体失配。提出首个基于3DGS的交互手部Avatar框架，通过解耦表示（优化身份图+学习潜特征）和交互感知模块，显著提升单图输入下的手部Avatar图像质量。适用时还应保留论文给出的边界：依赖MANO等参数化手部模型；对于极端遮挡和极端姿态可能受限；训练需要大规模手部数据集
% METHOD: 2411.16758
\paragraph{Motion-Aware Animatable Gaussian Avatars Deblurring（人体、头像与可动画角色）（2024.11）}
Motion-Aware Animatable Gaussian Avatars Deblurring~\cite{2411_16758} 关注的问题是：用 3DGS 建模人体、头像或可动画角色 Motion-Aware Animatable Gaussian Avatars Deblurring 是宽召回补全中发现的论文，主要根据标题、摘要和 arXiv 类别判定其与 3DGS 的关系。
作为 supplementary 类条目，Motion-Aware Animatable Gaussian Avatars Deblurring 不应被写成单纯的新渲染算子。它更像是把 3DGS 放进 人体化身与动画 的协议、系统或应用边界里：该工作被召回是因为论文文本包含 Gaussian/3D Gaussian 相关表述；当前分类先确定其任务域和与 3DGS 的相关性，后续证据流水线再补代码、企业和引用信息。
\textbf{实现细节：} 具体实现可以按输入、核心模块和渲染/优化路径来读。输入端首先取得任务所需的图像、相机、文本、时间序列、语义或几何线索，并把这些线索转换成可以约束高斯中心、尺度、旋转、不透明度、颜色或附加特征的条件。核心结构中，方法使用的主要机制是：从摘要层面看，Motion-Aware Animatable Gaussian Avatars Deblurring 围绕“用 3DGS 建模人体、头像或可动画角色”展开。若其正文或代码后续可读，应进一步核对它具体改动的是表示、初始化、渲染、优化、损失、动态建模、压缩或应用接口中的哪一环。 训练或推理时，这些模块通常先给出初始化或中间特征，再通过可微 Gaussian rasterization、重建损失、感知/语义约束、几何正则或任务损失回传到高斯参数；如果方法包含前馈网络，则网络输出的是高斯属性、变形场、压缩码、语义特征或控制变量，而不是只输出最终 2D 图像。渲染路径仍以显式高斯投影、深度排序和 alpha 合成为基础，只是在采样、属性预测、局部编辑、动态变形、压缩编码或下游任务接口处加入额外约束。这样写入本章时，它与前后方法的关系是：相对原始 3DGS 和本章已有锚点方法，它可以理解为 3DGS extends\_or\_uses 宽召回分类阶段的关系判断，后续需要按论文方法部分复核。
\textbf{证据与边界：} 本轮召回分类尚未逐项核对实验表格；后续证据补全阶段需要读取论文源码/PDF或项目页确认数据集、指标和对比方法。 该条来自宽召回批次；源码下载受 arXiv e-print 网络限制，部分论文暂时只能依据标题、摘要和类别分类，需要在入可信书前做人工/LLM精读复核。 本书只采用本地 paper pack、数据库分类、源码/引用登记和论文摘要中可核验的信息；没有本地证据支撑的速度、排行榜或开源完整性不在正文中强化。\textbf{方法信息：} 代码：\url{https://github.com/myniuuu/mad-avatar}。
\textbf{设计动机与边界：} 相对原始 3DGS 以独立高斯和图像重建误差为主的设定，Motion-Aware Animatable Gaussian Avatars Deblurring 把问题落在以下可实现的改动上：从摘要层面看，Motion-Aware Animatable Gaussian Avatars Deblurring 围绕“用 3DGS 建模人体、头像或可动画角色”展开。若其正文或代码后续可读，应进一步核对它具体改动的是表示、初始化、渲染、优化、损失、动态建模、压缩或应用接口中的哪一环。其设计目的不是单独增加一个模块，而是用上述表示、监督或渲染约束消除原流程在目标任务中的具体失配。该工作被召回是因为论文文本包含 Gaussian/3D Gaussian 相关表述；当前分类先确定其任务域和与 3DGS 的相关性，后续证据流水线再补代码、企业和引用信息。适用时还应保留论文给出的边界：该条来自宽召回批次；源码下载受 arXiv e-print 网络限制，部分论文暂时只能依据标题、摘要和类别分类，需要在入可信书前做人工/LLM精读复核。
% METHOD: 2412.06974
\paragraph{MV-DUSt3R+（单阶段多视图重建与高斯头）（2024）}
MV-DUSt3R+~\cite{2412_06974} 从 DUSt3R 和 MASt3R 的双视图模式扩展到任意多视图单阶段处理。多视角解码器同时学习参考视角与所有源视角关系，跨参考视角块缓解参考选择敏感性；扩展版本附加高斯 splatting head，预测像素对齐的 3D Gaussian 参数用于新视角合成。相对人体专用模型，它是通用 sparse-view reconstruction 骨干；放在本章中可作为未来人体多视图前馈重建的底层组件。
\textbf{实现细节：} 核心设计：1）MV-DUSt3R：设计多视角解码器块，同时学习一个参考视角与所有其他源视角的成对关系，以及源视角之间的关系。预测的各视角点图在同一参考坐标系下对齐，完全消除级联全局优化；2）MV-DUSt3R+：为解决参考视角选择的敏感性，引入跨参考视角块融合不同参考视角选择的信息，提升鲁棒性；3）高斯泼溅扩展：为支持新视角合成，在两个网络中附加并联合训练高斯泼溅头，预测像素对齐的3D高斯参数。整个推理过程仅需2秒。相比 DUSt3R/MASt3R 的 $O(n^{}2)$ 两两推理+全局优化，MV-DUSt3R+ 是 $O(n)$单阶段处理。 提出MV-DUSt3R+单阶段网络，可在一次前向传播中处理任意数量视角，无需级联全局优化。核心创新：多视角解码器块同时学习参考视角与所有源视角的关系，跨参考视角块融合多参考视角信息。扩展版本附加高斯泼溅头支持新视角合成。在多视角立体重建、姿态估计和新视角合成上显著超越前作。 计算量随视角数增加而增长；对极大规模视角集的可扩展性有待验证。
\textbf{实现补充：} 实现上还要核对输入、输出、参数、特征、初始化、训练损失、正则、渲染或推理接口，以及优化中的剪枝、分裂、采样、编码或解码步骤；这些抓手决定该方法到底改了 3DGS 管线的哪一环。
\textbf{方法信息：} Meta 参与了这篇论文。
\textbf{设计动机与边界：} 相对原始 3DGS 以独立高斯和图像重建误差为主的设定，MV-DUSt3R+ 把问题落在以下可实现的改动上：核心设计：1）MV-DUSt3R：设计多视角解码器块，同时学习一个参考视角与所有其他源视角的成对关系，以及源视角之间的关系。预测的各视角点图在同一参考坐标系下对齐，完全消除级联全局优化；2）MV-DUSt3R+：为解决参考视角选择的敏感性，引入跨参考视角块融合不同参考视角选择的信息，提升鲁棒性；3）高斯泼溅扩展：为支持新视角合成，在两个网络中附加并联合训练高斯泼溅头，预测像素对齐的3D高斯参数。整个推理过程仅需2秒。相比DUSt3R/MASt3R的O(n\^{}2)两两推理+全局优化，MV-DUSt3R+是O(n)单阶段处理。其设计目的不是单独增加一个模块，而是用上述表示、监督或渲染约束消除原流程在目标任务中的具体失配。提出MV-DUSt3R+单阶段网络，可在一次前向传播中处理任意数量视角，无需级联全局优化。核心创新：多视角解码器块同时学习参考视角与所有源视角的关系，跨参考视角块融合多参考视角信息。扩展版本附加高斯泼溅头支持新视角合成。在多视角立体重建、姿态估计和新视角合成上显著超越前作。适用时还应保留论文给出的边界：计算量随视角数增加而增长；对极大规模视角集的可扩展性有待验证。
% METHOD: 2412.11599
\paragraph{3D2-Actor（人体、头像与可动画角色）（2024.12）}
3D2-Actor~\cite{2412_11599} 关注的问题是：用 3DGS 建模人体、头像或可动画角色 3D\$\textasciicircum{}2\$-Actor 是宽召回补全中发现的论文，主要根据标题、摘要和 arXiv 类别判定其与 3DGS 的关系。
作为 supplementary 类条目，3D\$\textasciicircum{}2\$-Actor 不应被写成单纯的新渲染算子。它更像是把 3DGS 放进 人体化身与动画 的协议、系统或应用边界里：该工作被召回是因为论文文本包含 Gaussian/3D Gaussian 相关表述；当前分类先确定其任务域和与 3DGS 的相关性，后续证据流水线再补代码、企业和引用信息。
\textbf{实现细节：} 具体实现可以按输入、核心模块和渲染/优化路径来读。输入端首先取得任务所需的图像、相机、文本、时间序列、语义或几何线索，并把这些线索转换成可以约束高斯中心、尺度、旋转、不透明度、颜色或附加特征的条件。核心结构中，方法使用的主要机制是：从摘要层面看，3D$\textasciicircum{}2$-Actor 围绕“用 3DGS 建模人体、头像或可动画角色”展开。若其正文或代码后续可读，应进一步核对它具体改动的是表示、初始化、渲染、优化、损失、动态建模、压缩或应用接口中的哪一环。 训练或推理时，这些模块通常先给出初始化或中间特征，再通过可微 Gaussian rasterization、重建损失、感知/语义约束、几何正则或任务损失回传到高斯参数；如果方法包含前馈网络，则网络输出的是高斯属性、变形场、压缩码、语义特征或控制变量，而不是只输出最终 2D 图像。渲染路径仍以显式高斯投影、深度排序和 alpha 合成为基础，只是在采样、属性预测、局部编辑、动态变形、压缩编码或下游任务接口处加入额外约束。这样写入本章时，它与前后方法的关系是：相对原始 3DGS 和本章已有锚点方法，它可以理解为 3DGS extends\_or\_uses 宽召回分类阶段的关系判断，后续需要按论文方法部分复核。
\textbf{证据与边界：} 本轮召回分类尚未逐项核对实验表格；后续证据补全阶段需要读取论文源码/PDF或项目页确认数据集、指标和对比方法。 该条来自宽召回批次；源码下载受 arXiv e-print 网络限制，部分论文暂时只能依据标题、摘要和类别分类，需要在入可信书前做人工/LLM精读复核。 本书只采用本地 paper pack、数据库分类、源码/引用登记和论文摘要中可核验的信息；没有本地证据支撑的速度、排行榜或开源完整性不在正文中强化。\textbf{方法信息：} 代码：\url{https://github.com/silence-tang/gaussianactor}。
\textbf{设计动机与边界：} 相对原始 3DGS 以独立高斯和图像重建误差为主的设定，3D$^2$-Actor 把问题落在以下可实现的改动上：从摘要层面看，3D$^2$-Actor 围绕“用 3DGS 建模人体、头像或可动画角色”展开。若其正文或代码后续可读，应进一步核对它具体改动的是表示、初始化、渲染、优化、损失、动态建模、压缩或应用接口中的哪一环。其设计目的不是单独增加一个模块，而是用上述表示、监督或渲染约束消除原流程在目标任务中的具体失配。该工作被召回是因为论文文本包含 Gaussian/3D Gaussian 相关表述；当前分类先确定其任务域和与 3DGS 的相关性，后续证据流水线再补代码、企业和引用信息。适用时还应保留论文给出的边界：该条来自宽召回批次；源码下载受 arXiv e-print 网络限制，部分论文暂时只能依据标题、摘要和类别分类，需要在入可信书前做人工/LLM精读复核。
% METHOD: 2412.14963
\paragraph{IDOL（人体、头像与可动画角色）（2024.12）}
IDOL~\cite{2412_14963} 关注的问题是：用 3DGS 建模人体、头像或可动画角色 IDOL 是宽召回补全中发现的论文，主要根据标题、摘要和 arXiv 类别判定其与 3DGS 的关系。
作为 supplementary 类条目，IDOL 不应被写成单纯的新渲染算子。它更像是把 3DGS 放进 人体化身与动画 的协议、系统或应用边界里：该工作被召回是因为论文文本包含 Gaussian/3D Gaussian 相关表述；当前分类先确定其任务域和与 3DGS 的相关性，后续证据流水线再补代码、企业和引用信息。
\textbf{实现细节：} 具体实现可以按输入、核心模块和渲染/优化路径来读。输入端首先取得任务所需的图像、相机、文本、时间序列、语义或几何线索，并把这些线索转换成可以约束高斯中心、尺度、旋转、不透明度、颜色或附加特征的条件。核心结构中，方法使用的主要机制是：从摘要层面看，IDOL 围绕“用 3DGS 建模人体、头像或可动画角色”展开。若其正文或代码后续可读，应进一步核对它具体改动的是表示、初始化、渲染、优化、损失、动态建模、压缩或应用接口中的哪一环。 训练或推理时，这些模块通常先给出初始化或中间特征，再通过可微 Gaussian rasterization、重建损失、感知/语义约束、几何正则或任务损失回传到高斯参数；如果方法包含前馈网络，则网络输出的是高斯属性、变形场、压缩码、语义特征或控制变量，而不是只输出最终 2D 图像。渲染路径仍以显式高斯投影、深度排序和 alpha 合成为基础，只是在采样、属性预测、局部编辑、动态变形、压缩编码或下游任务接口处加入额外约束。这样写入本章时，它与前后方法的关系是：相对原始 3DGS 和本章已有锚点方法，它可以理解为 3DGS extends\_or\_uses 宽召回分类阶段的关系判断，后续需要按论文方法部分复核。
\textbf{证据与边界：} 本轮召回分类尚未逐项核对实验表格；后续证据补全阶段需要读取论文源码/PDF或项目页确认数据集、指标和对比方法。 该条来自宽召回批次；源码下载受 arXiv e-print 网络限制，部分论文暂时只能依据标题、摘要和类别分类，需要在入可信书前做人工/LLM精读复核。 本书只采用本地 paper pack、数据库分类、源码/引用登记和论文摘要中可核验的信息；没有本地证据支撑的速度、排行榜或开源完整性不在正文中强化。\textbf{方法信息：} 代码：\url{https://github.com/yiyuzhuang/IDOL}；引用信号：62 citations (3.3/month, domain-relative cutoff 3.0/month, semantic\_scholar)；代码关注：yiyuzhuang/idol has 319 stars (17.1/month, domain-relative cutoff 32.9/month)。
\textbf{设计动机与边界：} 相对原始 3DGS 以独立高斯和图像重建误差为主的设定，IDOL 把问题落在以下可实现的改动上：从摘要层面看，IDOL 围绕“用 3DGS 建模人体、头像或可动画角色”展开。若其正文或代码后续可读，应进一步核对它具体改动的是表示、初始化、渲染、优化、损失、动态建模、压缩或应用接口中的哪一环。其设计目的不是单独增加一个模块，而是用上述表示、监督或渲染约束消除原流程在目标任务中的具体失配。该工作被召回是因为论文文本包含 Gaussian/3D Gaussian 相关表述；当前分类先确定其任务域和与 3DGS 的相关性，后续证据流水线再补代码、企业和引用信息。适用时还应保留论文给出的边界：该条来自宽召回批次；源码下载受 arXiv e-print 网络限制，部分论文暂时只能依据标题、摘要和类别分类，需要在入可信书前做人工/LLM精读复核。
% METHOD: 2501.15008
\paragraph{HuGDiffusion（单图人体高斯属性扩散）（2025）}

HuGDiffusion~\cite{2501_15008} 从单张图像做人体新视角合成。它先提取 SMPL、法线、人解析等人体先验，再用多阶段扩散分别生成位置、尺度和旋转、不透明度、颜色等 3DGS 属性。关键不是只用像素损失，而是构建代理 ground-truth 3D Gaussian 属性作为属性级监督。相对 GPS-Gaussian 的双视图深度路线，HuGDiffusion 用扩散模型补足单视角缺失信息，输出的是完整高斯属性而非只从已见像素反投影。
\textbf{实现细节：} HuGDiffusion（2025.01） 是首个将条件扩散模型引入单视角人体 3DGS 新视角合成的工作。与前面 GPS-Gaussian 等前馈回归方法不同，HuGDiffusion 不再将高斯属性预测视为回归问题，而是视为条件生成问题，利用扩散模型在不可见区域生成合理外观的能力来弥补回归模型的固有缺陷。Baseline 与问题。现有可泛化 3DGS 框架（如 GPS-Gaussian）遵循统一的训练范式：参数化学习模型 $\mathcal{M}(\cdot;\theta)$ 以输入图像为条件，直接回归 3D 高斯属性集 $\mathcal{A}$，再通过可微分渲染施加像素级监督 $\min_\theta \sum_{k=1}^{K} \mathcal{L}_\mathrm{pmet}(\mathcal{R}(\mathcal{M}(\cdot;\theta); \mathbf{V}^{(k)}), \mathbf{I}^{(k)})$。这种回归式监督存在根本性局限：单视角输入中大量人体区域（背面、遮挡部位）不可见，回归模型缺乏对这些区域生成合理外观的能力——即使数值指标尚可，渲染出的不可见区域也经常出现大片黑色伪影和不自然纹理。改了什么。HuGDiffusion 对上述训练范式进行了两层根本性改变。（1）监督信号从像素级变为属性级：不再仅通过渲染图像与真实图像的像素差异来优化，而是直接在 3D 高斯属性空间上构建代理真值（proxy ground-truth），实施属性级监督 $\min_\theta \mathcal{L}_\mathrm{setdiff}(\mathcal{A}, \hat{\mathcal{A}})$。
\textbf{设计动机与边界：} 相对原始 3DGS 以独立高斯和图像重建误差为主的设定，HuGDiffusion 把问题落在以下可实现的改动上：HuGDiffusion的管线设计：(1) 人体特征提取：从单张输入图像提取多模态人体先验作为条件信号，包括SMPL参数估计、法线图、人体解析图等。(2) 多阶段生成策略：观察到联合学习全部3DGS属性困难，将生成分为多个阶段分别获取位置、尺度/旋转、不透明度和颜色/SH等属性。每阶段有专门的扩散模型和条件设计。(3) 属性级监督：构建代理GT 3D高斯属性作为高质量监督信号，而非仅依赖像素级损失。使用预训练3DGS模型拟合获取这些代理属性，提供更直接的属性空间监督。(4) 条件扩散模型：基于提取的人体先验作为条件，生成对应的3DGS属性。其设计目的不是单独增加一个模块，而是用上述表示、监督或渲染约束消除原流程在目标任务中的具体失配。提出HuGDiffusion，首个基于扩散模型的泛化3DGS人体渲染管线。从单视角图像出发，通过精心设计的人体特征提取和多阶段生成策略获取不同类型的3DGS属性。构建代理GT 3D高斯属性作为属性级监督信号，显著优于SOTA方法。适用时还应保留论文给出的边界：生成质量受限于单视角输入的信息量。扩散模型的采样速度限制了实时应用。对SMPL估计精度的依赖性。
\textbf{方法信息：} Tencent 参与了这篇论文；代码：\url{https://github.com/haiantyz/hugdiffusion}。



% METHOD: 2501.17085
\paragraph{Evaluating CrowdSplat（人群高斯感知质量评估）（2025）}
Evaluating CrowdSplat~\cite{2501_17085} 是评估工作。它用双选强制选择实验研究运动复杂度、LOD 级别和像素高度如何影响高斯人群 avatar 的感知质量。结果显示复杂动作比简单循环动作更容易暴露伪影，LOD 降低更容易被察觉，LOD 与像素高度有交互效应。相对 CrowdSplat 的系统实现，这篇工作给出读者选择高斯数量和视距策略的感知依据：人群渲染不能只看帧率，也要知道远距离和复杂运动下用户何时会发现简化。
\textbf{实现细节：} 采用双选强制选择（2AFC）实验范式，参与者观看成对的动画3D高斯头像，其中一个是最高质量的标准参考，另一个是较低LOD版本。实验设计包含三个独立变量：运动类型（2种：简单循环慢跑运动 vs 复杂武术表演）、LOD级别（3种：不同高斯数量）和像素高度（5种：对应不同视距）。参与者任务为选择更详细的头像。使用三因素方差分析（ANOVA）分析数据，并辅以Bonferroni事后检验。实验发现运动主效应显著，复杂运动的伪影更容易被察觉；LOD主效应显著，高斯数量越少伪影越明显；像素高度无主效应但与LOD存在交互效应。 首次系统性地通过用户实验（2AFC范式）评估3D高斯头像的感知质量，研究了运动复杂度、LOD级别（高斯数量）和像素高度（视距）三个因素及其交互效应。发现复杂运动比简单循环运动的伪影更明显，LOD降低导致伪影更易察觉，且LOD与像素高度存在交互效应。 仅评估感知质量而非客观渲染指标；实验场景和头像类型有限；未考虑更复杂的人群交互场景；参与者群体可能存在采样偏差
\textbf{实现补充：} 实现上还要核对输入、输出、参数、特征、初始化、训练损失、正则、渲染或推理接口，以及优化中的剪枝、分裂、采样、编码或解码步骤；这些抓手决定该方法到底改了 3DGS 管线的哪一环。
\textbf{方法信息：} Microsoft 参与了这篇论文。
\textbf{设计动机与边界：} 相对原始 3DGS 以独立高斯和图像重建误差为主的设定，Evaluating CrowdSplat 把问题落在以下可实现的改动上：采用双选强制选择（2AFC）实验范式，参与者观看成对的动画3D高斯头像，其中一个是最高质量的标准参考，另一个是较低LOD版本。实验设计包含三个独立变量：运动类型（2种：简单循环慢跑运动 vs 复杂武术表演）、LOD级别（3种：不同高斯数量）和像素高度（5种：对应不同视距）。参与者任务为选择更详细的头像。使用三因素方差分析（ANOVA）分析数据，并辅以Bonferroni事后检验。实验发现运动主效应显著，复杂运动的伪影更容易被察觉；LOD主效应显著，高斯数量越少伪影越明显；像素高度无主效应但与LOD存在交互效应。其设计目的不是单独增加一个模块，而是用上述表示、监督或渲染约束消除原流程在目标任务中的具体失配。首次系统性地通过用户实验（2AFC范式）评估3D高斯头像的感知质量，研究了运动复杂度、LOD级别（高斯数量）和像素高度（视距）三个因素及其交互效应。发现复杂运动比简单循环运动的伪影更明显，LOD降低导致伪影更易察觉，且LOD与像素高度存在交互效应。适用时还应保留论文给出的边界：仅评估感知质量而非客观渲染指标；实验场景和头像类型有限；未考虑更复杂的人群交互场景；参与者群体可能存在采样偏差
% METHOD: 2502.20220
\paragraph{Avat3r（人体、头像与可动画角色）（2025.02）}
Avat3r~\cite{2502_20220} 关注的问题是：用 3DGS 建模人体、头像或可动画角色 Avat3r 是宽召回补全中发现的论文，主要根据标题、摘要和 arXiv 类别判定其与 3DGS 的关系。
作为 supplementary 类条目，Avat3r 不应被写成单纯的新渲染算子。它更像是把 3DGS 放进 人体化身与动画 的协议、系统或应用边界里：该工作被召回是因为论文文本包含 Gaussian/3D Gaussian 相关表述；当前分类先确定其任务域和与 3DGS 的相关性，后续证据流水线再补代码、企业和引用信息。
\textbf{实现细节：} 具体实现可以按输入、核心模块和渲染/优化路径来读。输入端首先取得任务所需的图像、相机、文本、时间序列、语义或几何线索，并把这些线索转换成可以约束高斯中心、尺度、旋转、不透明度、颜色或附加特征的条件。核心结构中，方法使用的主要机制是：从摘要层面看，Avat3r 围绕“用 3DGS 建模人体、头像或可动画角色”展开。若其正文或代码后续可读，应进一步核对它具体改动的是表示、初始化、渲染、优化、损失、动态建模、压缩或应用接口中的哪一环。 训练或推理时，这些模块通常先给出初始化或中间特征，再通过可微 Gaussian rasterization、重建损失、感知/语义约束、几何正则或任务损失回传到高斯参数；如果方法包含前馈网络，则网络输出的是高斯属性、变形场、压缩码、语义特征或控制变量，而不是只输出最终 2D 图像。渲染路径仍以显式高斯投影、深度排序和 alpha 合成为基础，只是在采样、属性预测、局部编辑、动态变形、压缩编码或下游任务接口处加入额外约束。这样写入本章时，它与前后方法的关系是：相对原始 3DGS 和本章已有锚点方法，它可以理解为 3DGS extends\_or\_uses 宽召回分类阶段的关系判断，后续需要按论文方法部分复核。
\textbf{证据与边界：} 本轮召回分类尚未逐项核对实验表格；后续证据补全阶段需要读取论文源码/PDF或项目页确认数据集、指标和对比方法。 该条来自宽召回批次；源码下载受 arXiv e-print 网络限制，部分论文暂时只能依据标题、摘要和类别分类，需要在入可信书前做人工/LLM精读复核。 本书只采用本地 paper pack、数据库分类、源码/引用登记和论文摘要中可核验的信息；没有本地证据支撑的速度、排行榜或开源完整性不在正文中强化。\textbf{方法信息：} Avat3r project page lists Meta Reality Labs affiliation and code coming soon 参与；引用信号：49 citations (3.0/month, domain-relative cutoff 3.0/month, semantic\_scholar)。
\textbf{设计动机与边界：} 相对原始 3DGS 以独立高斯和图像重建误差为主的设定，Avat3r 把问题落在以下可实现的改动上：从摘要层面看，Avat3r 围绕“用 3DGS 建模人体、头像或可动画角色”展开。若其正文或代码后续可读，应进一步核对它具体改动的是表示、初始化、渲染、优化、损失、动态建模、压缩或应用接口中的哪一环。其设计目的不是单独增加一个模块，而是用上述表示、监督或渲染约束消除原流程在目标任务中的具体失配。该工作被召回是因为论文文本包含 Gaussian/3D Gaussian 相关表述；当前分类先确定其任务域和与 3DGS 的相关性，后续证据流水线再补代码、企业和引用信息。适用时还应保留论文给出的边界：该条来自宽召回批次；源码下载受 arXiv e-print 网络限制，部分论文暂时只能依据标题、摘要和类别分类，需要在入可信书前做人工/LLM精读复核。
% METHOD: 2503.20779
\paragraph{PGC（单姿态物理服装高斯）（2025）}
PGC~\cite{2503_20779} 从单个静态姿态的多视角捕获重建可仿真的服装。它用 mesh-embedded 3DGS：输入多视角服装图像和初始网格后，每个高斯都相对网格三角形局部坐标系参数化，位置、旋转和尺度随网格变形而更新。网格提供可仿真几何，3DGS 捕捉近场皮带、拉链、针织和毛皮等高频细节。外观训练中，方法用服装特定 PBR 着色模型估计反照率、粗糙度、金属度和间接光照；渲染时，高斯-PBR 混合模块让远场依赖物理网格，近场由高斯补细节。相对 LAGA 的生成式换装，PGC 更偏捕获和物理仿真。
\textbf{实现细节：} PGC（2025.03） 仅需要单个静态帧的多视角捕获即可重建可模拟的服装，采用混合网格嵌入 3D 高斯泼溅表示。高斯的位置定义为相对三角面片的偏移量，使其随网格变形自动移动，捕获飞散纤维等近场高频细节。同时提出基于 Disney BRDF 的布料专用 PBR 着色模型（含织物特有的光泽散射分量），从输入图像中分解出空间变化的反照率和粗糙度等材质参数。渲染时采用高斯-PBR 混合策略：低频成分由 PBR 着色提供姿态依赖的光影效果，高频成分由 3DGS 补充织物纹理细节，使单姿态重建即可泛化到新姿态和新光照。
\textbf{方法信息：} Meta 参与了这篇论文。
\textbf{设计动机与边界：} 相对原始 3DGS 以独立高斯和图像重建误差为主的设定，PGC 把问题落在以下可实现的改动上：PGC包含三个关键组件。（1）网格嵌入3DGS重建：每个高斯通过三角形局部坐标帧映射到网格上的三角形。高斯的位置、旋转和尺度相对于该帧参数化，使其与网格耦合。3DGS捕捉精细的体积面料细节（如皮带、口袋、拉链、针织和毛皮），网格提供用于仿真的几何基础。（2）物理渲染（PBR）外观重建：提出服装特定的物理着色模型从真实捕获中重建PBR外观，准确捕捉新姿态和新光照条件下的阴影和间接光照效果。（3）高斯-PBR混合渲染：结合物理渲染处理远场的姿态相关着色，使用3DGS重建填补近场的姿态无关细节。通过偏移量将高斯位置定义在网格上，使网格能辅助仿真新动力学和预测着色。其设计目的不是单独增加一个模块，而是用上述表示、监督或渲染约束消除原流程在目标任务中的具体失配。提出了PGC方法，使用混合网格嵌入3DGS表示服装，仅需单个静态帧的多视角捕获即可重建仿真就绪的服装，高斯捕捉近场着色和高频细节，网格编码远场反照率和反射参数，支持新姿态泛化和重光照。适用时还应保留论文给出的边界：需要多视角捕获设置；仅支持单层服装；对复杂纹理的建模可能受限。
% METHOD: 2503.12886
\paragraph{RGBAvatar（人体、头像与可动画角色）（2025.03）}
RGBAvatar~\cite{2503_12886} 关注的问题是：用 3DGS 建模人体、头像或可动画角色 RGBAvatar 是宽召回补全中发现的论文，主要根据标题、摘要和 arXiv 类别判定其与 3DGS 的关系。
作为 supplementary 类条目，RGBAvatar 不应被写成单纯的新渲染算子。它更像是把 3DGS 放进 人体化身与动画 的协议、系统或应用边界里：该工作被召回是因为论文文本包含 Gaussian/3D Gaussian 相关表述；当前分类先确定其任务域和与 3DGS 的相关性，后续证据流水线再补代码、企业和引用信息。
\textbf{实现细节：} 具体实现可以按输入、核心模块和渲染/优化路径来读。输入端首先取得任务所需的图像、相机、文本、时间序列、语义或几何线索，并把这些线索转换成可以约束高斯中心、尺度、旋转、不透明度、颜色或附加特征的条件。核心结构中，方法使用的主要机制是：从摘要层面看，RGBAvatar 围绕“用 3DGS 建模人体、头像或可动画角色”展开。若其正文或代码后续可读，应进一步核对它具体改动的是表示、初始化、渲染、优化、损失、动态建模、压缩或应用接口中的哪一环。 训练或推理时，这些模块通常先给出初始化或中间特征，再通过可微 Gaussian rasterization、重建损失、感知/语义约束、几何正则或任务损失回传到高斯参数；如果方法包含前馈网络，则网络输出的是高斯属性、变形场、压缩码、语义特征或控制变量，而不是只输出最终 2D 图像。渲染路径仍以显式高斯投影、深度排序和 alpha 合成为基础，只是在采样、属性预测、局部编辑、动态变形、压缩编码或下游任务接口处加入额外约束。这样写入本章时，它与前后方法的关系是：相对原始 3DGS 和本章已有锚点方法，它可以理解为 3DGS extends\_or\_uses 宽召回分类阶段的关系判断，后续需要按论文方法部分复核。
\textbf{证据与边界：} 本轮召回分类尚未逐项核对实验表格；后续证据补全阶段需要读取论文源码/PDF或项目页确认数据集、指标和对比方法。 该条来自宽召回批次；源码下载受 arXiv e-print 网络限制，部分论文暂时只能依据标题、摘要和类别分类，需要在入可信书前做人工/LLM精读复核。 本书只采用本地 paper pack、数据库分类、源码/引用登记和论文摘要中可核验的信息；没有本地证据支撑的速度、排行榜或开源完整性不在正文中强化。\textbf{方法信息：} 代码：\url{https://github.com/gapszju/rgbavatar}；代码关注：gapszju/rgbavatar has 134 stars (8.5/month, domain-relative cutoff 32.9/month)。
\textbf{设计动机与边界：} 相对原始 3DGS 以独立高斯和图像重建误差为主的设定，RGBAvatar 把问题落在以下可实现的改动上：从摘要层面看，RGBAvatar 围绕“用 3DGS 建模人体、头像或可动画角色”展开。若其正文或代码后续可读，应进一步核对它具体改动的是表示、初始化、渲染、优化、损失、动态建模、压缩或应用接口中的哪一环。其设计目的不是单独增加一个模块，而是用上述表示、监督或渲染约束消除原流程在目标任务中的具体失配。该工作被召回是因为论文文本包含 Gaussian/3D Gaussian 相关表述；当前分类先确定其任务域和与 3DGS 的相关性，后续证据流水线再补代码、企业和引用信息。适用时还应保留论文给出的边界：该条来自宽召回批次；源码下载受 arXiv e-print 网络限制，部分论文暂时只能依据标题、摘要和类别分类，需要在入可信书前做人工/LLM精读复核。
% METHOD: 2504.07949
\paragraph{InteractAvatar（手-脸非刚性交互）（2025.04）}
InteractAvatar~\cite{2504_07949} 面向手接触脸部的头像。它把面部用 FLAME 网格绑定高斯，手部用 MANO 模板和 Dynamic Gaussian Hand 表示。手部 MLP 建模姿态依赖的关节褶皱和自阴影；手-面部交互模块再用 MLP 预测接触时的非刚性面部变形和外观变化，例如皮肤凹陷。相对普通头像方法只处理表情，InteractAvatar 把手作为主动变形源；相对 GHOST，它目标不是重建未知物体，而是保留手-脸接触的照片级细节。
\textbf{实现细节：} InteractAvatar（2025.04） 首次建模手-脸交互的光真实外观。其采用混合 mesh-Gaussian 表示，将高斯绑定到 FLAME（面部）和 MANO（手部）网格面上，随网格变形而刚性变换，兼顾 mesh 的变形可控性和 3DGS 的细节表现力。针对手部的高自由度运动，动态精修模块通过两组 MLP 分别建模几何动态（位置、缩放、旋转偏移）和外观动态（颜色、透明度偏移），捕获关节弯曲时的细皱纹和自投射阴影等姿态依赖变化。手-脸交互模块则利用 DECAF 的基于位置动力学（PBD）的碰撞求解器获取顶点偏移，并通过 PointNet++ 提取手-脸全局几何特征，结合相对位姿和面部表情参数，预测交互引起的非刚性变形和外观变化；距离加权的偏移调制确保远离接触点的区域不受交互模块影响。
\textbf{方法信息：} Meta 参与了这篇论文。
\textbf{设计动机与边界：} 相对原始 3DGS 以独立高斯和图像重建误差为主的设定，InteractAvatar 把问题落在以下可实现的改动上：方法基于GaussianAvatar，使用混合网格-高斯表示：(1) 面部头像：FLAME网格参数化面部，每个网格面片绑定高斯，通过局部坐标系随网格运动；(2) 手部头像：MANO模板模型参数化手部，动态高斯手模型使用MLP建模姿态依赖的动态效果（自投射阴影、关节褶皱）；(3) 手-面部交互模块：使用MLP建模手接触面部时的非刚性变形，捕获皮肤凹陷和外观变化。从多视角视频训练，通过多视角模板拟合获取FLAME和MANO参数。支持新视角合成、自重演和跨身份重演。其设计目的不是单独增加一个模块，而是用上述表示、监督或渲染约束消除原流程在目标任务中的具体失配。提出InteractAvatar，首个利用3DGS建模非刚性手-面部变形和交互的方法。核心创新包括：(1) Dynamic Gaussian Hand模型结合模板模型和3DGS加动态精炼模块，捕获姿态依赖变化（如皱纹和阴影）；(2) 手-面部交互模块建模微妙几何和外观动态；(3) 混合网格-高斯表示绑定高斯到网格面片实现动画。支持从单目或多视角视频重建手和手-面部交互的高保真细节，可由新姿态驱动。适用时还应保留论文给出的边界：需要多视角模板拟合获取FLAME/MANO参数；极端手部姿态可能超出训练分布；计算成本较高。
% METHOD: 2504.09097
\paragraph{BIGS（双手未知物体高斯交互）（2025.04）}
BIGS~\cite{2504_09097} 处理双手与未知物体。它只优化右手高斯，再通过左右翻转共享给左手，在统一规范空间累积双手信息，缓解遮挡。物体部分使用预训练 text-to-image diffusion 的 SDS loss 补全未见区域，即使输入视角有限也能恢复完整形状。最后的 interaction subject optimization 加入穿透、接触和渲染损失，联合优化手部高斯、物体高斯和交互参数。项目包含两个相关仓库。相对 GHOST 的检索几何先验，BIGS 更依赖扩散先验补遮挡物体。
\textbf{实现细节：} BIGS（2025.04） 解决了从单目 RGB 视频进行双手类别无关交互重建的问题。其核心设计是双手共享高斯：在规范空间中仅优化单手（右手）的 3D 高斯，通过对 $x$ 轴翻转将其映射到左手，使两只手的 3D 信息在同一规范空间中有效积累，缓解双手严重遮挡下的信息不足。手部和物体各自使用 TriplaneNet 提取特征，分别通过外观、几何和变形 MLP 预测高斯参数，再经 LBS 或刚体变换映射到图像空间。对于被遮挡的物体不可见区域，BIGS 引入基于预训练扩散模型的分数蒸馏采样（SDS）损失，结合文本反转学习物体描述和 PiDi 边界条件 ControlNet 引导，在不依赖物体形状先验的情况下重建完整物体几何。
\textbf{方法信息：} 代码：\url{https://github.com/on-jungwoan/bigs}。
\textbf{设计动机与边界：} 相对原始 3DGS 以独立高斯和图像重建误差为主的设定，BIGS 把问题落在以下可实现的改动上：管线处理单目视频输入：(1) 预处理：使用手部网格回归器获取MANO姿态和形状参数，使用物体网格估计器获取物体网格，作为高斯初始位置；(2) 手部高斯：仅优化右手高斯，通过左右翻转共享给左手，在统一规范空间中有效累积两只手的3D信息，缓解严重遮挡；(3) 物体高斯：使用预训练T2I扩散模型的SDS损失重建被遮挡的物体部分，即使有限视角也能恢复完整3D形状；(4) 交互主体优化：引入多种损失函数确保手-物3D对齐，包括穿透损失、接触损失和渲染损失。联合优化手部高斯、物体高斯和交互参数。其设计目的不是单独增加一个模块，而是用上述表示、监督或渲染约束消除原流程在目标任务中的具体失配。提出BIGS，首个从单目视频重建双手和未知物体3D高斯溅射的方法。核心创新：(1) 利用预训练扩散模型的SDS损失重建严重遮挡下的未见物体部分；(2) 利用MANO手部模型先验，两只手共享单一高斯模型通过翻转有效累积3D信息；(3) 引入交互主体优化步骤在高斯优化中考虑手和物体的3D对齐。在3D手部姿态估计、3D物体重建和渲染质量上达到SOTA。适用时还应保留论文给出的边界：依赖MANO手部模型先验；对极端遮挡场景的物体重建仍有挑战；需要准确的2D手部和物体检测作为预处理。
% METHOD: 2510.07729
\paragraph{ComGS（可重光照物体-场景组合）（2025）}
ComGS~\cite{2510_07729} 使用 2DGS 重建可重光照物体，并用 Surface Octahedral Probes 存储间接光和遮挡信息，实现快速编辑和实时组合渲染。重建阶段先优化辐射场和几何，再分解材质和光照；编辑阶段从场景中提取局部辐射场，并用扩散补全环境光照。它与人体 avatar 的关系在于组合和重光照：当人体、服装或道具需要插入场景时，类似 SOPs 的局部光照缓存可提供更真实阴影。
\textbf{实现细节：} ComGS管线分为三个阶段：（1）重建阶段：使用2DGS进行可微多目标渲染生成G-buffer（RGB、深度、法线、albedo、粗糙度、金属度），分两步优化——先重建辐射场和几何，再分解材质和光照。材质采用简化Disney BRDF模型，光照分解为直接光照（可学习环境贴图）和间接光照。引入SOPs存储间接光照和遮挡信息，通过插值而非光线追踪获取着色信息，实现至少2倍重建加速和实时阴影计算。（2）编辑阶段：在物体放置位置对高斯场景进行360度追踪提取部分辐射场，使用微调的扩散模型补全完整环境光照；同时缓存SOPs的遮挡信息。（3）渲染阶段：结合可重光照物体和场景实现实时渲染。 提出Surface Octahedral Probes（SOPs）实现高效可重光照高斯物体重建（2倍加速），通过局部环境光照估计和扩散模型补全实现逼真的阴影生成，构建了完整的ComGS物体-场景组合框架，支持约28fps实时渲染和36秒快速编辑。 局部光照估计可能无法处理复杂全局光照传输；扩散模型补全光照可能引入不一致性；SOPs分辨率和精度之间存在权衡。
\textbf{实现补充：} 实现上还要核对输入、输出、参数、特征、初始化、训练损失、正则、渲染或推理接口，以及优化中的剪枝、分裂、采样、编码或解码步骤；这些抓手决定该方法到底改了 3DGS 管线的哪一环。
\textbf{方法信息：} Huawei 参与了这篇论文。
\textbf{设计动机与边界：} 相对原始 3DGS 以独立高斯和图像重建误差为主的设定，ComGS 把问题落在以下可实现的改动上：ComGS管线分为三个阶段：（1）重建阶段：使用2DGS进行可微多目标渲染生成G-buffer（RGB、深度、法线、albedo、粗糙度、金属度），分两步优化——先重建辐射场和几何，再分解材质和光照。材质采用简化Disney BRDF模型，光照分解为直接光照（可学习环境贴图）和间接光照。引入SOPs存储间接光照和遮挡信息，通过插值而非光线追踪获取着色信息，实现至少2倍重建加速和实时阴影计算。（2）编辑阶段：在物体放置位置对高斯场景进行360度追踪提取部分辐射场，使用微调的扩散模型补全完整环境光照；同时缓存SOPs的遮挡信息。（3）渲染阶段：结合可重光照物体和场景实现实时渲染。其设计目的不是单独增加一个模块，而是用上述表示、监督或渲染约束消除原流程在目标任务中的具体失配。提出Surface Octahedral Probes（SOPs）实现高效可重光照高斯物体重建（2倍加速），通过局部环境光照估计和扩散模型补全实现逼真的阴影生成，构建了完整的ComGS物体-场景组合框架，支持约28fps实时渲染和36秒快速编辑。适用时还应保留论文给出的边界：局部光照估计可能无法处理复杂全局光照传输；扩散模型补全光照可能引入不一致性；SOPs分辨率和精度之间存在权衡。
% METHOD: 2601.05853
\paragraph{LayerGS（2DGS 多层人体分解与修复）（2026）}
LayerGS~\cite{2601_05853} 从任意姿态人体分解身体和服装。它使用 2DGS 附着在 SMPL 或 SMPL-X 规范网格表面，以获得更准确表面几何。训练分三阶段：先单层重建并用 Stable Diffusion SDS 修复遮挡区域，再冻结外层优化内层身体，最后冻结内层精化外层，处理层间重叠伪影。每层可提取网格并用 LBS 新姿态驱动。项目实现围绕 layered avatars。相对 LAGA 的文本生成，LayerGS 重点是从图像分解已有穿衣人体并支持虚拟试穿。
\textbf{实现细节：} LayerGS（2026.01） 将任意姿态的人体分解为可动画的多层化身，每层编码为一组 2D 高斯实现准确的几何和光真实渲染。该方法采用三阶段训练策略：先训练单层规范化身重建外层服装的粗略形状，再冻结外层优化内层身体，最后冻结内层精修外层以解决层间交叠伪影。通过预训练 2D 扩散模型（Stable Diffusion V1.5 + ControlNet）的分数蒸馏采样修复被遮挡区域。训练完成后从每层 2D 高斯提取 3D 网格并将高斯绑定到对应三角面片，利用线性混合蒙皮驱动网格变形并附加 Laplacian 变形解决层间穿透，实现可动画的 3D 虚拟试穿。
\textbf{设计动机与边界：} 相对原始 3DGS 以独立高斯和图像重建误差为主的设定，LayerGS 把问题落在以下可实现的改动上：三阶段训练流程：Stage 1（单层规范化化身重建与修复）——从多视图RGB图像出发，在SMPL(-X)规范网格表面初始化30k个2DGS，预计算LBS权重。使用SDS损失（Stable Diffusion V1.5）修复遮挡区域生成粗外层服装。应用深度失真损失和法线一致性损失确保几何准确性。Stage 2（多层训练）——冻结Stage 1外层，优化规范内层身体。Stage 3（精细外层训练）——冻结规范内层，训练精化外层解决层间重叠导致的渲染伪影。最终从每层提取3D网格并将对应Gaussian附着其上，通过LBS蒙皮实现新姿态下的3D虚拟试穿重组。关键使用2DGS而非3DGS获得更准确的表面几何。为了消除原始 3DGS 在该目标任务中的具体失配，方法把上述表示、监督或渲染约束作为可学习和可执行的改动，而不是单独堆叠模块。提出LayerGS框架，通过2DGS将任意姿态人体分解为身体和服装的多层可动画化身，利用预训练2D扩散模型通过SDS修复遮挡区域，三阶段训练策略实现内外层的精确重建，支持新视角、新姿态下的虚拟试穿。适用时还应保留论文给出的边界：需要多视图输入；修复质量依赖扩散模型先验；三阶段训练流程较长。
\textbf{方法信息：} 代码：\url{https://github.com/rockyxu66/layergs}。



% METHOD: 2602.20556
\paragraph{WildGHand（野外扰动解耦手部高斯）（2026.02）}
WildGHand~\cite{2602_20556} 面向野外单目手部视频，存在手-物交互、极端姿态、运动模糊和光照变化。它基于 MANO-HD 建模手部，并把扰动显式表示为高斯颜色和位置的时变 bias，从而分离干净手部外观和干扰信息。扰动感知优化策略生成逐帧各向异性权重 mask，低权重处理高扰动区域；SIFT 跟踪提供稀疏几何约束。项目实现包含 benchmark。相对单图手部 avatar，WildGHand 的关键是把“干扰”作为可学习变量，而不是简单拟合输入。
\textbf{实现细节：} 手部、服装与扩展应用共享"部件级解耦 + 先验注入"思路。手部方向上 JGHand 用关节感知高斯增强姿态，BIGS 处理生物力学约束下的双手交互，WildGHand 在野外环境重建可驱动手部。服装方向上 GS-VTON 实现 3D 虚拟试穿，GGAvatar 从单视频分离衣物，GarmentGS 专门建模服装几何，Cloth-HUGS 与 ReWeaver 分别融合 HUGS 与可仿真编织拓扑。扩展应用则把上述能力推向更广场景：OccFusion 用占用融合处理遮挡人体，LAGA 通过语言引导生成可动画化身，HOGSA 用分层高斯优化手-物场景，PGC 用物理引导从单姿态生成可仿真衣物。
\textbf{实现补充：} 实现上还要核对输入、输出、参数、特征、初始化、训练损失、正则、渲染或推理接口，以及优化中的剪枝、分裂、采样、编码或解码步骤；这些抓手决定该方法到底改了 3DGS 管线的哪一环。
\textbf{方法信息：} 代码已公开，仓库为 xuanhuang0/wildghand (2 stars)。入选依据：企业:否；代码:是；星标:否；引用:否。
\textbf{设计动机与边界：} 相对原始 3DGS 以独立高斯和图像重建误差为主的设定，WildGHand 把问题落在以下可实现的改动上：WildGHand基于MANO-HD手部模型参数化，在3DGS框架上构建两个核心组件：(1) 动态扰动解耦模块：将扰动显式建模为3D高斯属性上的时变偏置（bias），每个高斯的颜色和位置都有可学习的时变偏置项，在优化过程中分离干净的手部外观和干扰信息。(2) 扰动感知优化策略：生成逐帧各向异性加权掩码引导优化，根据扰动置信度对不同区域赋予不同权重，抑制高扰动区域的影响。结合MANO-HD模型提供的手部几何先验和SIFT特征点跟踪提供的稀疏几何约束。整体框架采用自适应优化范式，而非简单拟合输入视频。还构建了一个包含多种扰动条件的野外手部视频数据集作为benchmark。其设计目的不是单独增加一个模块，而是用上述表示、监督或渲染约束消除原流程在目标任务中的具体失配。提出WildGHand框架，通过动态扰动解耦和扰动感知优化策略，首次实现从野外单目视频（存在手-物交互、极端姿态、运动模糊、光照变化等干扰）的高保真3D高斯手部化身重建。PSNR相对提升15.8%，LPIPS相对降低23.1%。适用时还应保留论文给出的边界：依赖MANO模型的参数化表达，对极端手部形变（如紧握）可能不够灵活；多手交互场景需要分别处理每只手；长时间视频的优化效率需提升。
% METHOD: 2603.18912
\paragraph{GHOST（类别无关手-物交互重建）（2026.03）}
GHOST~\cite{2603_18912} 从单目 RGB 视频快速重建手-物体交互。预处理用 SAM2 分割物体、SfM 估计相机和稀疏点云、MANO 初始化手部，并通过 InternVL 和 OpenShape 从 Objaverse 检索物体几何先验。抓取感知对齐优化手部平移和物体尺度，随后用 2D Gaussian disks 联合重建手和物体。几何先验一致性损失补全遮挡，手感知背景损失避免手遮挡物体区域被误惩罚。相对 BIGS，GHOST 主要处理单手和类别无关物体。
\textbf{实现细节：} GHOST（2026.03） 修改了 3DGS 在手-物联合重建中的表示与优化环节——用 2D 高斯盘替代 3D 高斯椭球表示物体与手部，引入几何先验一致性损失和抓取感知对齐替代纯光度优化——多视角渲染管线保持不变。Baseline 与问题：从单目 RGB 视频重建手-物交互面临三重挑战。首先，标准 3DGS 使用 3D 高斯椭球表示物体，在稀疏视角和频繁遮挡下，被手指遮挡的物体区域无法获得有效监督，导致重建出现空洞或手指形状的"洞"。其次，SfM 恢复的物体点云与 HaMeR 估计的手部平移存在尺度不一致和空间漂移，导致手与物体在 3D 空间中并不真正接触。最后，手部是高度变形的非刚体，标准 3DGS 的全局高斯集合无法跟随 MANO 网格的姿态变化。改了什么：GHOST 从表示、几何和空间对齐三个层面进行改进。表示层面，物体和手部均采用 2D 高斯盘（2DGS）而非 3D 高斯椭球，提供更紧凑的表面建模。几何层面，引入基于大规模 3D 物体数据库（Objaverse）的几何先验检索与对齐机制，以及几何一致性损失 $L_{geo}$ 替代仅依赖光度监督的范式。空间对齐层面，提出抓取感知（grasp-aware）联合优化，同时精修手部平移 $T_t^h$ 和物体全局尺度 $S_o$。
\textbf{方法信息：} 代码：\url{https://github.com/ataboukhadra/ghost}。
\textbf{设计动机与边界：} 相对原始 3DGS 以独立高斯和图像重建误差为主的设定，GHOST 把问题落在以下可实现的改动上：GHOST分三阶段：(1)预处理：使用SAM2分割物体，SfM估计相机参数和稀疏点云，MANO初始化手部网格，通过InternVL+OpenShape从Objaverse检索几何先验网格；(2)手-物体对齐：通过抓取感知推理优化手部平移和物体尺度，使用射线投射算法和ICP对齐先验网格与SfM点云；(3)GS优化：使用2D高斯圆盘联合重建手和物体，引入几何先验一致性损失补全遮挡区域，手感知背景损失避免对手遮挡物体区域的惩罚。整个过程比HOLD等NeRF方法快一个数量级。其设计目的不是单独增加一个模块，而是用上述表示、监督或渲染约束消除原流程在目标任务中的具体失配。提出GHOST（Gaussian Hand-Object Splatting），一个快速、类别无关的手-物体交互重建框架。使用2D高斯泼溅表示手和物体，引入三个创新：(1)几何先验检索和一致性损失补全遮挡区域；(2)抓取感知对齐优化手部平移和物体尺度；(3)手感知背景损失。比先前方法快一个数量级。适用时还应保留论文给出的边界：依赖Objaverse数据库的覆盖度；对严重形变物体效果可能下降；几何先验检索的准确性受VLM描述影响。
% METHOD: 2606.01493
\paragraph{Splatshot（每步 3D 反馈的单图人脸生成）（2026）}
SplatShot~\cite{2606_01493} 处理单张非受控照片到 3D 人脸高斯头像的问题。它的基线困境很明确：前馈 3DGS 人脸模型能给出显式几何一致性，但遇到野外照片时外观质量和身份保持容易退化；扩散模型能补出高质量多视角图像，却容易让不同视角的鼻梁、眼眶、发际线和侧脸轮廓互相矛盾。SplatShot 的改变是把这两者放进同一个去噪循环，而不是先生成多视角图再事后拟合 3DGS。

具体流程是：给定一张参考照片和一个基础 3DGS 人脸模板，扩散模型在多个目标视角上联合去噪；在每个 denoising step，系统先从当前噪声状态预测干净图像，再用这些多视角预测重拟合一套 3DGS，随后把这套 3DGS 重渲染回各视角，与扩散预测图像计算光度差异，并把差异反传回噪声估计。这个 3D feedback 让扩散采样不断向“同一组高斯能解释的多视角头像”靠拢。代码仓库的入口是 \texttt{inference.py}，输出包括最终 \texttt{avatar.ply}、基座模型、扩散中间图像和 COLMAP 相机；项目结构中也单独列出 3DGS-guided denoising loop、Gaussian trainer 和 diffusion wrapper。

Samsung Research America 参与了这篇论文。代码发布在 \url{https://github.com/hliang2/splatshot}，GitHub 约 43 stars；由于论文很新，目前引用还少。相对“扩散生成后再优化 3DGS”的流程，SplatShot 的重点是每一步都用显式高斯重投影纠正多视角不一致。
\textbf{锚点细化：} Splatshot：在 3DGS 的“人脸头像生成”基础上推进，主要改变是：SplatShot：训练自由框架，在扩散去噪过程中反复重拟合3DGS并用3D反馈约束多视角一致性。它要解决的问题是：提出每个去噪步的3D反馈循环：从多视角预测图像重拟合3DGS，再把光度差异反传到扩散过程，以结合2D先验和显式3D一致性。核心思想是：SplatShot把3DGS几何一致性直接嵌入扩散去噪过程，从单照生成3D人脸头像。
	extbf{实现细节：} Splatshot 在扩散去噪的每一步插入显式三维反馈，而不是只要求多视图图像彼此相似。给定基础人脸 3DGS 和参考照片，去噪器联合预测多个目标视角的干净图像；系统用这些预测重新拟合临时 3DGS，把该模型重渲染回各视角，并将重渲染与当前图像的光度差作为跨视图三维一致性损失反馈给下一步去噪。因此推理是“多视图去噪—拟合高斯—渲染比较—反馈更新”的循环；它依赖人脸先验，且迭代重拟合带来额外计算。
	extbf{实现细节：} Splatshot 在扩散去噪的每一步插入显式三维反馈，而不是只要求多视图图像彼此相似。给定基础人脸 3DGS 和参考照片，去噪器联合预测多个目标视角的干净图像；系统用这些预测重新拟合临时 3DGS，把该模型重渲染回各视角，并将重渲染与当前图像的光度差作为跨视图三维一致性损失反馈给下一步去噪。因此推理是“多视图去噪—拟合高斯—渲染比较—反馈更新”的循环；它依赖人脸先验，且迭代重拟合带来额外计算。
\textbf{方法信息：} Samsung Research America 参与；代码：\url{https://github.com/hliang2/splatshot}；GitHub 约 43 stars；目前引用还少。



% METHOD: 2606.13655
\paragraph{Flex4DHuman（相对位姿编码的动态多视角补全）（2026）}
Flex4DHuman~\cite{2606_13655} 从单目或稀疏多视角动态人物视频生成密集同步多视角视频，用于后续 4D 人体重建。它不用骨架、深度、法线或目标视角几何先验，而是用相对相机位姿编码驱动 Wan 视频扩散骨干，并通过五轴位置编码扩展时空 RoPE，注入视角索引和连续 $SE(3)$ 位姿。相对 AniGS，它不直接输出高斯，而是为 4DGS 重建生成更完整多视角视频条件。
\textbf{实现细节：} Flex4DHuman（2026.06） 延续 4D Gaussian splats 在"动态人体多视角生成"上的基本设定，主要改变是：Flex4DHuman：多视角视频扩散模型，仅用相对相机位姿条件生成可送入4D Gaussian splats重建的视图。高斯基元仍承担可渲染动态场景表示。其设计动机是：用相对相机位姿编码替代骨架、深度、法线或目标视角几何先验，使动态人物多视角补全更灵活。具体而言，模型基于Wan视频扩散骨干，通过五轴位置编码扩展时空RoPE以注入视角索引和连续SE(3)位姿，输出同步多视角视频。适用边界是：本身不是3DGS重建方法；下游4D Gaussian质量依赖生成视频一致性。
\textbf{实现补充：} 实现上还要核对输入、输出、参数、特征、初始化、训练损失、正则、渲染或推理接口，以及优化中的剪枝、分裂、采样、编码或解码步骤；这些抓手决定该方法到底改了 3DGS 管线的哪一环。
\textbf{方法信息：} 代码仓库：\url{https://github.com/andy-cheng/flex4dhuman}；GitHub 关注度较高，约 89 stars。
\textbf{设计动机与边界：} 相对原始 3DGS 以独立高斯和图像重建误差为主的设定，Flex4DHuman 把问题落在以下可实现的改动上：模型基于Wan视频扩散骨干，通过五轴位置编码扩展时空RoPE以注入视角索引和连续SE(3)位姿，输出同步多视角视频。其设计目的不是单独增加一个模块，而是用上述表示、监督或渲染约束消除原流程在目标任务中的具体失配。用相对相机位姿编码替代骨架、深度、法线或目标视角几何先验，使动态人物多视角补全更灵活。适用时还应保留论文给出的边界：本身不是3DGS重建方法；下游4D Gaussian质量依赖生成视频一致性。
% METHOD: 2606.19156
\paragraph{Hand-4DGS（第一视角前馈 4D 手部）（2026）}
Hand-4DGS~\cite{2606_19156} 把前馈 3DGS 推到第一视角动态手部重建。它面对的输入不是稳定外部相机拍摄的人体，而是 egocentric video：相机随头部移动，手部快速运动，还经常出现双手和物体遮挡。HUGS、EVA 这类逐序列优化方法可以改成手部模板，但仍依赖每段视频的优化和手部姿态估计；一般前馈高斯方法又多针对静态人体或稀疏视角图像，不能直接建模 4D 手部动态。

Hand-4DGS 的做法是把手部 mesh 先验前馈预测出来，再把高斯锚到 mesh 结构上。给定视频帧，网络先提取视觉特征，并用 temporal convolution 聚合相邻帧上下文；随后浅层 MLP 预测 MANO 形状、姿态和全局平移，得到每帧手部网格。方法在网格顶点上学习位置嵌入，并在网格面内插值得到高斯嵌入；高斯位置和旋转由 mesh 几何显式计算，颜色与尺度由 MLP 解码，不透明度则被正则到接近 1，因为手本身应是不透明物体。训练时先用现成估计器产生的伪 3D 顶点监督稳定结构，再逐步引入 2D 渲染监督，使预测 mesh 和高斯外观一起被图像误差修正。论文报告约 60 FPS 推理，并在 H2O、ARCTIC 等第一视角手物交互数据上评估。

Microsoft 参与了这篇论文；截至本次核查，未看到完整公开代码。相对逐视频优化的动态手部高斯，Hand-4DGS 的增量是把 MANO 结构先验、时序卷积和可微高斯渲染合成一个可泛化的前馈 4D 手部表示。
\textbf{实现细节：} Hand-4DGS。Hand-4DGS（2026.06） 延续 4D Gaussian Splatting 在"手部与第一视角重建"上的基本设定，主要改变是：Hand-4DGS：mesh-guided表示和temporal convolutions结合的前馈4D手部Gaussian Splatting框架。高斯基元仍承担可渲染动态场景表示。其设计动机是：将mesh结构先验和时间卷积引入前馈4DGS手部重建，实现约60 FPS推理和跨场景泛化。具体而言，方法使用手部mesh指导高斯布局和结构约束，并通过时序卷积建模快速手部运动，直接从第一视角视频输出动态高斯手部。适用边界是：依赖手部mesh先验；严重遮挡和手物交互仍困难。
\textbf{设计动机与边界：} 相对原始 3DGS 以独立高斯和图像重建误差为主的设定，Hand-4DGS 把问题落在以下可实现的改动上：方法使用手部mesh指导高斯布局和结构约束，并通过时序卷积建模快速手部运动，直接从第一视角视频输出动态高斯手部。为了消除原始 3DGS 在该目标任务中的具体失配，方法把上述表示、监督或渲染约束作为可学习和可执行的改动，而不是单独堆叠模块。将mesh结构先验和时间卷积引入前馈4DGS手部重建，实现约60 FPS推理和跨场景泛化。适用时还应保留论文给出的边界：依赖手部mesh先验；严重遮挡和手物交互仍困难。
\textbf{方法信息：} Microsoft 参与；截至本次核查，未看到完整公开代码。



% METHOD: 2606.26092
\paragraph{TryOnCrafter（相机可控视频虚拟试穿）（2026）}
TryOnCrafter~\cite{2606_26092} 解决 camera-controllable video virtual try-on。传统视频试穿常被源视频相机轨迹锁住，难以自由改变视角；该方法强调结构同步、纹理幻化和背景一致性，并使用可渲染 4D try-on proxy 支持任意相机轨迹合成。当前素材未直接展开代理是否为标准 4D Gaussian 表示，因此正文只把它作为可渲染人体代理和试穿方向的边界案例。相对 LayerGS 的分层 3D avatar，它更偏视频生成和相机轨迹控制。
\textbf{实现细节：} TryOnCrafter（2026.06） 延续 4DGS 在"Camera-controllable virtual try-on"上的基本设定，主要改变是：构建可渲染 4D 试衣代理以支持任意视角动态渲染。不改变 3DGS 主线的表示、渲染或优化管线。其设计动机是：突破传统 VVT 只能跟随源相机轨迹的问题，让服装试穿视频可在任意相机轨迹下合成。具体而言，摘要未直接确认代理是否为 4D Gaussian 表示，但任务与动态可渲染人体代理高度接近 4DGS 应用；方法强调结构同步、纹理幻化和背景一致性。适用边界是：3DGS 证据来自关键词召回和可渲染 4D proxy 描述，若全文未使用 Gaussian 表示则相关性需下调。
\textbf{实现补充：} 实现上还要核对输入、输出、参数、特征、初始化、训练损失、正则、渲染或推理接口，以及优化中的剪枝、分裂、采样、编码或解码步骤；这些抓手决定该方法到底改了 3DGS 管线的哪一环。
\textbf{方法信息：} Alibaba 参与了这篇论文。
\textbf{设计动机与边界：} 相对原始 3DGS 以独立高斯和图像重建误差为主的设定，TryOnCrafter 把问题落在以下可实现的改动上：摘要未直接确认代理是否为 4D Gaussian 表示，但任务与动态可渲染人体代理高度接近 4DGS 应用；方法强调结构同步、纹理幻化和背景一致性。其设计目的不是单独增加一个模块，而是用上述表示、监督或渲染约束消除原流程在目标任务中的具体失配。突破传统 VVT 只能跟随源相机轨迹的问题，让服装试穿视频可在任意相机轨迹下合成。适用时还应保留论文给出的边界：3DGS 证据来自关键词召回和可渲染 4D proxy 描述，若全文未使用 Gaussian 表示则相关性需下调。
% METHOD: 2606.32036
\paragraph{PointSplat（面向直播的紧凑人体高斯）（2026）}

PointSplat~\cite{2606_32036} 面向沉浸式直播中的紧凑人体表示。它把预测从 view-centric 图像空间转为 3D human-centric representation，直接学习和生成紧凑 Gaussian 或 point-based 人体表示，以减少跨视角冗余并降低带宽。相对 Human Avatar Compression 的编码框架，PointSplat 更像前馈预测模型：从输入中直接得到适合传输和渲染的人体高斯。
\textbf{实现细节：} PointSplat（2026.06） 延续 3DGS avatar 在"Compact human Gaussian representations"上的基本设定，主要改变是：PointSplat：在 3D 空间进行 human-centric prediction 以减少多视图冗余。高斯基元仍承担可渲染动态场景表示。其设计动机是：将预测从 view-centric 转为 3D 空间中的 human-centric representation，以同时保持高保真和低带宽。具体而言，网络直接学习并生成紧凑 Gaussian/point-based 人体表示，减少跨视角冗余，适配有限计算和传输带宽的 live streaming。适用边界是：依赖人体先验和输入视角质量；宽松衣物、遮挡和快速动作可能增加难度。
\textbf{实现补充：} 为避免只把它读成名称或数据点，还应把实现接口明确拆开：输入是什么，输出是高斯参数、图像、轨迹、标注还是评测指标，哪些参数、网络模块、训练损失或渲染接口真正接到原始 3DGS / Scaffold-GS / 4D-GS 管线上，以及剪枝、分裂、采样、解码、优化或评测协议在哪一步发生。若该条目是survey、benchmark 或 dataset，这里的“实现”对应论文集合、任务划分、相机/轨迹/标注组织、评价指标和基线配置，而不是虚构新的 rasterizer；若是算法条目，则对应高斯位置、尺度、旋转、不透明度、特征或 MLP 解码如何进入训练和渲染流程。
\textbf{设计动机与边界：} 相对原始 3DGS 以独立高斯和图像重建误差为主的设定，PointSplat 把问题落在以下可实现的改动上：网络直接学习并生成紧凑 Gaussian/point-based 人体表示，减少跨视角冗余，适配有限计算和传输带宽的 live streaming。其设计目的不是单独增加一个模块，而是用上述表示、监督或渲染约束消除原流程在目标任务中的具体失配。将预测从 view-centric 转为 3D 空间中的 human-centric representation，以同时保持高保真和低带宽。适用时还应保留论文给出的边界：依赖人体先验和输入视角质量；宽松衣物、遮挡和快速动作可能增加难度。
\textbf{方法信息：} ByteDance 参与了这篇论文；代码：\url{https://github.com/anhthuan1999/pointsplat}。
